\providecommand{\Bun}{\operatorname{Bun}}
\providecommand{\Loc}{\operatorname{Loc}}
\providecommand{\Hecke}{\operatorname{Hecke}}

\chapter{Derived structures and the geometric Langlands correspondence}
\label{ch:derived-langlands}

\index{geometric Langlands!derived structures|textbf}
\index{oper|textbf}
\index{bar complex!critical level|textbf}

The geometric Langlands programme produces a categorical equivalence
\[
\mathcal{D}\text{-mod}(\Bun_G)
\;\simeq\;
\mathrm{IndCoh}\bigl(\Loc_{G^\vee}\bigr)
\]
between D-modules on the moduli of $G$-bundles and ind-coherent
sheaves on the stack of $G^\vee$-local systems. Categories, however,
are opaque: the equivalence exists, but the mechanism that forces it
remains invisible at the categorical level. What \emph{produces} the
oper space? What structural property of the vertex algebra at
critical level \emph{compels} the center to jump? The categorical
framework registers the outcome; it does not exhibit the cause.

Bar-cobar duality exhibits the cause. The bar complex
$\barB(\widehat{\fg}_k)$ depends on the level~$k$ through a single
scalar: the modular characteristic
$\kappa(\widehat{\fg}_k) = \dim(\fg)(k + h^\vee)/(2h^\vee)$.
%: formula from CLAUDE.md C3; k=0 -> dim(g)/2; k=-h^v -> 0 verified
At generic level, $\kappa(\widehat{\fg}_k) \neq 0$ and the bar complex is curved:
$d^2 \neq 0$, no cohomology, no geometry. At the critical level
$k = -h^\vee$, the numerator $k + h^\vee$ vanishes, $\kappa(\widehat{\fg}_{-h^\vee}) = 0$,
and the bar complex becomes an honest chain complex. The cohomology
that appears is the algebra of differential forms on opers:
\[
H^n\bigl(\barB(\widehat{\fg}_{-h^\vee})\bigr)
\;\cong\;
\Omega^n\bigl(\mathrm{Op}_{\fg^\vee}(D)\bigr)
\qquad\text{for all } n \ge 0.
\]
This is Theorem~\ref{thm:oper-bar-dl}. It does not prove the full
categorical localization; it identifies the chain-level mechanism from
which the oper side of the Frenkel--Gaitsgory programme emerges.

To see why no other construction yields this identification, observe
the dichotomy. At generic level, the curvature $\kappa(\widehat{\fg}_k) \neq 0$
obstructs the formation of cohomology: the bar complex is a curved
dg-coalgebra with $d^2 = \kappa(\widehat{\fg}_k) \cdot \omega_1 \neq 0$, and no
functor from curved coalgebras to graded algebras recovers opers.
At critical level, $\kappa(\widehat{\fg}_{-h^\vee}) = 0$ is the \emph{unique} value at which
the bar differential squares to zero. The cobar construction
$\Omega(\barB(\widehat{\fg}_{-h^\vee}))$ then inverts to recover the
vertex algebra (Theorem~B), and the bar cohomology computes the oper
package. Any approach that avoids the bar complex must independently
reconstruct this vanishing; the Feigin--Frenkel proof of the center
theorem does exactly that, using the Segal--Sugawara construction.
Bar-cobar duality is the unique homological framework in which the
center theorem, the oper identification, and the Koszul structure
appear as facets of a single $d^2 = 0$ condition.

The $\mathfrak{sl}_2$ case makes the mechanism concrete. Here
$h^\vee = 2$, $\dim(\mathfrak{sl}_2) = 3$, and the critical level
is $k = -2$. The modular characteristic evaluates to
$\kappa(\widehat{\mathfrak{sl}_2}_{-2}) = 3 \cdot (-2 + 2)/(2 \cdot
2) = 0$: uncurved bar complex, honest cohomology. The classical
r-matrix $r(z) = k\Omega_{\mathrm{tr}}/z$ at $k = -2$ gives
$r(z) = -2\Omega_{\mathrm{tr}}/z \neq 0$
(equivalently, $r(z) = \Omega/((k{+}h^\vee)z)$ diverges at $k = -h^\vee$: this
is the Sugawara singularity in the KZ normalisation);
% Two conventions: k*Omega_tr/z (finite at critical level, trace-form)
% vs Omega/((k+h^v)*z) (diverges at critical level, KZ). Same physics.
the r-matrix does not vanish, the OPE structure persists, but the
\emph{curvature} vanishes. The center jumps precisely because
$\kappa(\widehat{\fg}_{-h^\vee}) = 0$ uncurves the bar complex while the r-matrix continues to
encode the full vertex algebra structure. This separation of
curvature from the r-matrix is the structural content of the critical
level.

At admissible levels $k = -h^\vee + p/q$, the curvature is nonzero
but rational, and the bar complex conjecturally acquires a periodic
CDG structure with period~$2p$
(Conjecture~\ref{conj:periodic-cdg}). Given this periodicity,
Positselski's Koszul duality for CDG-algebras would recover the
Kazhdan--Lusztig equivalence
$\cO_k^{\mathrm{int}}(\widehat{\fg}) \simeq
\mathcal{C}(U_q(\fg))$
on the semisimplified tilting quotient. Any lift beyond the
semisimplified target to the full quantum-group category, or to a
completed/coderived enlargement, is a separate problem
requiring additional inputs beyond the MC3 evaluation-generated core proved in
Corollary~\ref{cor:mc3-all-types}.

\begin{theorem}[Critical bar-to-oper bridge; \ClaimStatusProvedHere]
\label{thm:langlands-bar-bridge}
\index{geometric Langlands!bridge theorem|textbf}
\index{oper!bridge from bar complex|textbf}
Let $\fg$ be a simple finite-dimensional Lie algebra with dual
Coxeter number~$h^\vee$, and let $\fg^\vee$ denote the Langlands
dual Lie algebra. At the critical level $k = -h^\vee$:
\begin{enumerate}[label=\textup{(\roman*)}]
\item \textup{(Curvature vanishing)}\quad
 $\kappa(\widehat{\fg}_{-h^\vee}) = 0$
 \textup{(}Corollary~\textup{\ref{cor:critical-level-universality}}\textup{)},
 so $\barB(\widehat{\fg}_{-h^\vee})$ is an honest chain complex
 \textup{(}uncurved: $m_0 = 0$, hence $d^2 = 0$\textup{)}.
\item \textup{(Oper differential-form identification)}\quad
 The bar cohomology computes the graded algebra of differential forms
 on the oper space:
 \[
 H^n\bigl(\barB(\widehat{\fg}_{-h^\vee})\bigr)
 \;\cong\;
 \Omega^n\bigl(\mathrm{Op}_{\fg^\vee}(D)\bigr)
 \quad\text{for all } n \ge 0.
 \]
 At $n = 0$ this is the Feigin--Frenkel center
 $Z(\widehat{\fg}_{-h^\vee}) \cong
 \mathrm{Fun}(\mathrm{Op}_{\fg^\vee}(D))$;
 at $n = 1$ this gives K\"ahler differentials of the oper space;
 at $n \ge 2$ this gives the higher differential-form package.
\item \textup{(Derived center)}\quad
 The derived center of the completed enveloping algebra is
 \[
 Z_{\mathrm{der}}(\widehat{\fg}_{-h^\vee})
 \;:=\;
 RHH^0_{\mathrm{ch}}(\widehat{\fg}_{-h^\vee})
 \]
 and is computed in this chapter using the critical bar complex
 $\barB(\widehat{\fg}_{-h^\vee})$ as a resolution input. Separately,
 the cohomology of that bar complex recovers the oper differential
 forms above.
 This supplies the manuscript's internal bridge to the oper side of
 the critical geometric Langlands picture; any stronger
 automorphic/spectral categorical identification remains external to
 the theorem proved here.
\end{enumerate}
\end{theorem}

\begin{proof}
Part~(i) is Corollary~\ref{cor:critical-level-universality}: the
modular characteristic $\kappa(\widehat{\fg}_k)$ vanishes at the critical level
because the Killing form $(\cdot,\cdot)_k$ degenerates at
$k = -h^\vee$. Part~(ii) is proved degree by degree:
$H^0$ is Theorem~\ref{thm:oper-bar-h0-dl} (via the PBW spectral
sequence and the Feigin--Frenkel theorem), $H^1$ is
Proposition~\ref{prop:oper-bar-h1-dl}, $H^2$ is
Proposition~\ref{prop:oper-bar-h2-dl}, and the general case
follows from the Whitehead spectral decomposition
(Proposition~\ref{prop:whitehead-spectral-decomposition})
together with formal smoothness of the oper space
(Frenkel--Gaitsgory~\cite{FG06}). Part~(iii) is the internal
bookkeeping statement that the derived center
$RHH^0_{\mathrm{ch}}(\widehat{\fg}_{-h^\vee})$ is computed here via
the critical bar complex as a bar-Ext resolution; the bar complex
itself remains distinct from the derived center, and its cohomology
recovers the oper differential-form package. This uses the bar-Ext identification
together with the Frenkel--Teleman theorem~\cite{FT06}; it does
not identify the bar complex with the derived global sections of the
oper structure sheaf.
\end{proof}

\begin{remark}[Periodicity open problem]
The theorematic core of this chapter is asymmetric. The critical-level
bar/oper identification is proved, and the manuscript establishes
rank-$1$ chiral Hochschild periodicity at critical level in
Proposition~\ref{prop:sl2-periodicity-dl}. The stronger admissible-level
statement that the full bar complex carries a periodic CDG structure,
especially in the higher-rank form needed for Kazhdan--Lusztig
transport, remains conjectural and belongs to the MC3/periodicity
extension problem \textup{(}MC3 categorical Clebsch--Gordan closure and
the generated-core DK comparison are proved for all simple types, with
post-CG completion open beyond type~$A$; the periodicity aspect is
open\textup{)}. Even on the KL side, the ordinary semisimplified target
$\mathcal{C}(U_q(\fg))$ and any larger non-semisimple or
completed/coderived lift must be kept separate.
\end{remark}

\begin{construction}[Kazhdan--Lusztig as categorical MC3]
\label{constr:kl-categorical-mc3}
\ClaimStatusProvedHere
\index{Kazhdan--Lusztig equivalence!categorical MC3|textbf}
\index{MC3!Kazhdan--Lusztig realization|textbf}
\index{thick generation!prefundamental|textbf}
The Kazhdan--Lusztig equivalence $\mathrm{KL}_k \colon \mathcal{O}^{\mathrm{int}}_k(\hat{\mathfrak{g}}) \xrightarrow{\sim} \mathrm{Rep}^{\mathrm{fd}}(U_q(\mathfrak{g}))$ at $q = e^{\pi i/(k+h^\vee)}$ is best read as a categorical analogue of the MC3 strategy rather than as a theorematic export of the repaired Yangian MC3 chain. On the Yangian side, the proved comparison is the DK transport on the evaluation-generated core \textup{(}Corollary~\ref{cor:dk23-all-types}\textup{)}, with the post-core prefundamental/completion packets kept separate \textup{(}Corollary~\ref{cor:mc3-all-types}, Theorem~\ref{thm:mc3-type-a-resolution}\textup{)}.

The shared mechanism is generator-and-transport on a controlled subcategory. In the KL setting, the category $\mathcal{O}^{\mathrm{int}}_k(\hat{\mathfrak{g}})$ is generated by Weyl modules $V_\lambda$, and the quantum group $U_q(\mathfrak{g})$ acts on that generating set via the $R$-matrix. In the MC3 setting, the proved Yangian comparison is on the evaluation-generated core, where the dg-shifted Yangian acts via the spectral kernel. What is not claimed here is that the repaired Yangian MC3 theorem itself proves KL, or that full prefundamental thick generation of $\mathcal{O}^{\mathrm{sh}}$ is already available on the general Yangian surface.
\end{construction}

%================================================================
% SECTION 1: DERIVED ALGEBRAIC GEOMETRY PREREQUISITES
%================================================================

\section{Derived algebraic geometry: prerequisites}
\label{sec:dag-prerequisites}
\index{derived algebraic geometry}

The exposition follows
Lurie~\cite{HA} and To\"en--Vezzosi.

\begin{definition}[Derived stack]
\label{def:derived-stack}
\index{derived stack|textbf}
A \emph{derived stack} (over a field~$k$ of characteristic zero) is
a functor
$X\colon \mathrm{cdga}_k^{\mathrm{op}} \to \mathcal{S}$
from commutative dg-algebras to the $\infty$-category of spaces,
satisfying \'etale descent. A derived stack~$X$ has:
\begin{enumerate}[label=\textup{(\roman*)}]
\item A \emph{classical truncation} $t_0(X)$, which is an ordinary
 algebraic stack obtained by restricting to ordinary (non-dg)
 commutative algebras.
\item A \emph{cotangent complex} $\mathbb{L}_X \in \mathrm{QCoh}(X)$,
 a quasi-coherent sheaf on~$X$ (in the derived sense) controlling
 deformations and obstructions. For a smooth classical
 scheme~$Y$ viewed as a derived stack, $\mathbb{L}_Y = \Omega^1_Y[0]$
 is the ordinary sheaf of differentials.
\item A \emph{derived structure sheaf} $\cO_X^{\mathrm{der}}$,
 a sheaf of cdga's whose cohomology sheaves
 $\mathcal{H}^i(\cO_X^{\mathrm{der}})$ encode the derived
 information beyond the classical truncation.
\end{enumerate}
\end{definition}

\begin{definition}[Formal smoothness]
\label{def:formal-smoothness}
\index{formal smoothness}
A derived stack~$X$ is \emph{formally smooth} if its cotangent
complex $\mathbb{L}_X$ is concentrated in cohomological degree~$0$
(i.e., $\mathbb{L}_X \simeq \Omega^1_{t_0(X)}[0]$). Equivalently,
$X$ has no higher obstructions: every infinitesimal deformation
extends.
\end{definition}

\begin{remark}[Formal smoothness vs.\ derived richness]
\label{rem:formal-smooth-derived}
The oper space $\mathrm{Op}_{\fg^\vee}(D)$ on a formal disk is
formally smooth (Frenkel--Gaitsgory~\cite{FG06}), so its cotangent
complex is classical. Nevertheless, the bar complex
$\barB(\widehat{\fg}_{-h^\vee})$ produces a chain complex
whose cohomology recovers
$\Omega^\bullet(\mathrm{Op}_{\fg^\vee}(D))$: the extra chain-level
structure is in the bar complex, not in a nontrivial derived
structure sheaf on the oper space. The bar complex is better viewed
as a chain-level enhancement of the oper differential-form algebra.
\end{remark}


%================================================================
% SECTION 2: THE OPER SPACE
%================================================================

\section{The oper space}
\label{sec:oper-space}
\index{oper!definition|textbf}

\begin{definition}[Opers; Beilinson--Drinfeld]
\label{def:opers}
\index{oper!Beilinson--Drinfeld}
Let $G^\vee$ be the Langlands dual group of~$G$, with
Lie algebra $\fg^\vee$, Borel subalgebra
$\mathfrak{b}^\vee \subset \fg^\vee$, and nilpotent radical
$\mathfrak{n}^\vee \subset \mathfrak{b}^\vee$. Let $X$ be a smooth algebraic
curve. A \emph{$G^\vee$-oper} on~$X$ is a triple
$(F_{G^\vee}, F_{B^\vee}, \nabla)$ where:
\begin{enumerate}[label=\textup{(\roman*)}]
\item $F_{G^\vee}$ is a principal $G^\vee$-bundle on~$X$;
\item $F_{B^\vee} \subset F_{G^\vee}$ is a reduction to the Borel
 subgroup~$B^\vee$;
\item $\nabla$ is a connection on $F_{G^\vee}$ satisfying the
 \emph{oper condition}: the induced map
 $TX \to F_{G^\vee}/F_{B^\vee} = F_{G^\vee} \times_{G^\vee} (G^\vee/B^\vee)$
 given by the connection lies in the open $B^\vee$-orbit.
\end{enumerate}
The moduli space of $G^\vee$-opers on~$X$ is denoted
$\mathrm{Op}_{G^\vee}(X)$.
\end{definition}

\begin{example}[SL2 opers as projective connections]
\label{ex:sl2-opers}
\index{oper!$\mathrm{SL}_2$}
For $G = \mathrm{SL}_2$, opers on a curve~$X$ are
\emph{projective connections}: second-order differential operators
$\partial_z^2 + t(z)$ on~$X$. On a formal disk $D = \mathrm{Spec}\, k[[z]]$, the space of
$\mathrm{PGL}_2$-opers is $\mathrm{Op}_{\mathrm{PGL}_2}(D) \cong k[[z]]$,
parametrized by the ``projective connection'' $t(z) = \sum_{n \geq 0} t_n z^n$.
In general, $\mathrm{Op}_{G^\vee}(D) \cong k[[z]]^{\mathrm{rank}(\fg)}$
is a formal power series space of dimension $\mathrm{rank}(\fg)$.
\end{example}

\begin{theorem}[Feigin--Frenkel center; \ClaimStatusProvedElsewhere]
\label{thm:ff-center-dl}
\index{Feigin--Frenkel center|textbf}
\textup{(}Feigin--Frenkel~\cite{Feigin-Frenkel},
\cite{FG06}\textup{)}\quad
The center of the completed enveloping algebra of
$\widehat{\fg}$ at the critical level $k = -h^\vee$ is
canonically isomorphic to the algebra of functions on
$\fg^\vee$-opers:
\begin{equation}\label{eq:ff-center}
Z(\widehat{\fg}_{-h^\vee})
\;\cong\;
\mathrm{Fun}\bigl(\mathrm{Op}_{\fg^\vee}(D)\bigr).
\end{equation}
The center $Z(\widehat{\fg}_{-h^\vee})$ is a polynomial
algebra on generators in conformal weights
$d_1, d_2, \ldots, d_r$ \textup{(}the exponents of $\fg$
plus~$1$\textup{)}.
\end{theorem}


%================================================================
% SECTION 3: BAR COMPLEX AT CRITICAL LEVEL
%================================================================

\section{Bar complex at critical level and the oper differential-form package}
\label{sec:bar-critical}
\index{bar complex!critical level}

\begin{theorem}[Zeroth bar cohomology = oper functions;
\ClaimStatusProvedHere]
\label{thm:oper-bar-h0-dl}
\index{bar complex!$H^0 = \mathrm{Fun}(\mathrm{Op})$}
For $\fg$ a simple finite-dimensional Lie algebra, the zeroth
cohomology of the reduced bar complex at critical level is:
\begin{equation}\label{eq:h0-oper}
H^0\bigl(\barB(\widehat{\fg}_{-h^\vee})\bigr)
\;\cong\;
\mathrm{Fun}\bigl(\mathrm{Op}_{\fg^\vee}(D)\bigr).
\end{equation}
\end{theorem}

\begin{proof}
At the critical level, the curvature vanishes
($\kappa(\widehat{\fg}_{-h^\vee}) = 0$,
Corollary~\ref{cor:critical-level-universality}), so the bar
complex is an honest chain complex.
The PBW spectral sequence
(Theorem~\ref{thm:pbw-koszulness-criterion}) has
\[
E_1^{p,q}
= H^q\bigl(\fg;\, S^p(\fg[t^{-1}]t^{-1})\bigr)
\;\Longrightarrow\;
H^{p+q}\bigl(\barB(\widehat{\fg}_{-h^\vee})\bigr).
\]
At $p + q = 0$: the Whitehead lemma gives
$H^0(\fg; S^*(\fg[t^{-1}]t^{-1})) = S^*(\fg[t^{-1}]t^{-1})^\fg$,
the invariant subalgebra. By the Feigin--Frenkel theorem
(Theorem~\ref{thm:ff-center-dl}), this invariant subalgebra is exactly
$\mathrm{Fun}(\mathrm{Op}_{\fg^\vee}(D))$. The higher
differentials $d_r$ for $r \geq 1$ vanish on total degree~$0$
by degree considerations, so $E_\infty^{0,0} = E_1^{0,0}$.
\end{proof}

\begin{proposition}[First bar cohomology as oper one-forms;
\ClaimStatusProvedHere]
\label{prop:oper-bar-h1-dl}
\index{bar complex!$H^1 = \Omega^1(\mathrm{Op})$}
\begin{equation}\label{eq:h1-oper}
H^1\bigl(\barB(\widehat{\fg}_{-h^\vee})\bigr)
\;\cong\;
\Omega^1\bigl(\mathrm{Op}_{\fg^\vee}(D)\bigr).
\end{equation}
\end{proposition}

\begin{proof}
On the $E_1$ page of the PBW spectral sequence, the total degree-$1$
component is $E_1^{1,0} \oplus E_1^{0,1}$. The term
$E_1^{0,1} = H^1(\fg; S^*(\fg[t^{-1}]t^{-1}))$ vanishes by the
Whitehead lemma ($H^1(\fg; M) = 0$ for any finite-dimensional
$\fg$-module~$M$ when $\fg$ is semisimple). The surviving term
$E_1^{1,0} = \Omega^1(S^*(\fg[t^{-1}]t^{-1})^\fg)$ is the module of
K\"ahler differentials of the invariant subalgebra. By formal
smoothness of the oper space (Frenkel--Gaitsgory~\cite{FG06}),
this is $\Omega^1(\mathrm{Op}_{\fg^\vee}(D))$. Again, higher
differentials vanish by degree.
\end{proof}

\begin{proposition}[Second bar cohomology as oper two-forms;
\ClaimStatusProvedHere]
\label{prop:oper-bar-h2-dl}
\index{bar complex!$H^2 = \Omega^2(\mathrm{Op})$}
\begin{equation}\label{eq:h2-oper}
H^2\bigl(\barB(\widehat{\fg}_{-h^\vee})\bigr)
\;\cong\;
\Omega^2\bigl(\mathrm{Op}_{\fg^\vee}(D)\bigr).
\end{equation}
\end{proposition}

\begin{proof}
On the $E_1$ page of the PBW spectral sequence, the total
degree-$2$ component is
$E_1^{2,0} \oplus E_1^{1,1} \oplus E_1^{0,2}$.

The coefficient module $S^p(\fg[t^{-1}]t^{-1})$ decomposes as a
direct sum of finite-dimensional $\fg$-modules, graded by conformal
weight: each component
$S^p(\fg[t^{-1}]t^{-1})_w$ is a finite-dimensional representation
of the semisimple Lie algebra~$\fg$. By Whitehead's first and
second lemmas, $H^q(\fg; V) = 0$ for $q = 1, 2$ and any
finite-dimensional $\fg$-module~$V$ when $\fg$ is semisimple.
Since Lie algebra cohomology commutes with direct sums:
\begin{align*}
E_1^{1,1} &= H^1\bigl(\fg;\, S^1(\fg[t^{-1}]t^{-1})\bigr) = 0, \\
E_1^{0,2} &= H^2\bigl(\fg;\, S^0(\fg[t^{-1}]t^{-1})\bigr)
 = H^2(\fg; k) = 0.
\end{align*}
The surviving term is $E_1^{2,0} = \Omega^2(S^*(\fg[t^{-1}]t^{-1})^\fg)
= \Omega^2(\mathrm{Op}_{\fg^\vee}(D))$ by formal smoothness.
Higher differentials $d_r$ ($r \geq 2$) on $E_r^{2,0}$ land in
$E_r^{2+r, 1-r}$ where $1 - r < 0$, hence vanish; differentials
$d_r\colon E_r^{2-r, r-1} \to E_r^{2,0}$ originate from terms with
$q = r - 1 \geq 1$, which are zero by the Whitehead vanishing above.
\end{proof}

\begin{proposition}[Whitehead spectral decomposition;
\ClaimStatusProvedHere]
\label{prop:whitehead-spectral-decomposition}
\index{bar complex!Whitehead decomposition|textbf}
\index{Whitehead lemma!spectral decomposition}
On the PBW spectral sequence for~$\barB(\widehat{\fg}_{-h^\vee})$,
the $E_1$ page decomposes as a tensor product:
\begin{multline}\label{eq:whitehead-spectral-decomposition}
E_1^{p,q}
= H^q\bigl(\fg;\, S^p(\fg[t^{-1}]t^{-1})\bigr)
\cong
S^p(\fg[t^{-1}]t^{-1})^{\fg}
\otimes
H^q(\fg;\, k)\\
\cong
\mathrm{Fun}^p\bigl(\mathrm{Op}_{\fg^\vee}(D)\bigr)
\otimes H^q(\fg;\, k).
\end{multline}
In particular, $E_1^{p,q} = 0$ unless~$q$ belongs to the
set of degrees of the exterior algebra
$\Lambda(\omega_{2d_1-1},\allowbreak \ldots,\allowbreak \omega_{2d_r-1})$,
where $d_1, \ldots, d_r$ are the exponents of~$\fg$
incremented by~$1$. For
$\fg = \mathfrak{sl}_2$ ($r = 1$, $d_1 = 2$),
the only nonzero rows are $q = 0$ and $q = 3$.
\end{proposition}

\begin{proof}
The coefficient module $S^p(\fg[t^{-1}]t^{-1})$ decomposes
as a $\fg$-module into a direct sum of finite-dimensional
representations, graded by conformal weight. By the
\emph{generalized Whitehead lemma}: for any simple~$\fg$ and any
nontrivial finite-dimensional irreducible $\fg$-module~$V$,
$H^q(\fg;\, V) = 0$ for \emph{all}~$q$. \textup{(}Proof: the
Casimir element~$C$ acts on~$V$ by the nonzero scalar
$\langle \lambda,\, \lambda + 2\rho\rangle$ but trivially on
cohomology, forcing $H^q(\fg;\, V) = 0$.\textup{)}

Decompose $S^p(\fg[t^{-1}]t^{-1})
= M^{\mathrm{triv}} \oplus M^{\mathrm{nontriv}}$
into trivial and nontrivial isotypic components. Since Lie
algebra cohomology commutes with direct sums:
\[
H^q\bigl(\fg;\, S^p(\fg[t^{-1}]t^{-1})\bigr)
= H^q(\fg;\, M^{\mathrm{triv}})
= S^p(\fg[t^{-1}]t^{-1})^{\fg} \otimes H^q(\fg;\, k).
\]
The identification
$S^*(\fg[t^{-1}]t^{-1})^{\fg}
\cong \mathrm{Fun}(\mathrm{Op}_{\fg^\vee}(D))$
is the Feigin--Frenkel theorem~\textup{(}Theorem~\ref{thm:ff-center-dl}\textup{)}.
By Chevalley's theorem,
$H^*(\fg;\, k) \cong
\Lambda(\omega_{2d_1-1}, \ldots, \omega_{2d_r-1})$, which is
nonzero precisely in degrees that are sums of distinct elements
from $\{2d_1 - 1, \ldots, 2d_r - 1\}$.

\emph{Remark on earlier proofs.}
The Whitehead vanishing results used in the proofs of
Propositions~\ref{prop:oper-bar-h1-dl}
and~\ref{prop:oper-bar-h2-dl} ($H^1 = H^2 = 0$ with
coefficients) are special cases of this decomposition: the
classical Whitehead lemmas ($q = 1, 2$) follow from the
generalized version, and the vanishing of $E_1^{p,q}$ for
$q = 1, 2$ is immediate from $H^1(\fg; k) = H^2(\fg; k) = 0$.
\end{proof}

\begin{remark}[Spectral consequences by rank]
\label{rem:spectral-by-rank}
\index{bar complex!rank dependence}
The Whitehead decomposition gives $2^r$ nonzero rows on the $E_1$
page for rank~$r$: $\mathfrak{sl}_2$ has $2$ rows ($d_4$ only),
$\mathfrak{sl}_3$ has $4$ rows ($d_4$, $d_6$, composites).
Degeneration occurs after $\leq 2d_r$ steps ($d_r$ the largest
exponent), with the abutment controlled by transgression
differentials
$d_{2d_i}\colon \omega_{2d_i-1} \mapsto \tau_i \in
\mathrm{Fun}^{2d_i}(\mathrm{Op})$.
\end{remark}


\begin{proposition}[Differential analysis at degree 3;
\ClaimStatusProvedHere]
\label{prop:h3-differential-analysis}
\index{bar complex!$H^3$ obstruction}
On the PBW spectral sequence for~$\barB(\widehat{\fg}_{-h^\vee})$,
the total degree-$3$ terms on the $E_1$ page are:
\begin{align*}
E_1^{3,0} &= S^3(\fg[t^{-1}]t^{-1})^{\fg}
 \;\cong\; \Omega^3(\mathrm{Op}_{\fg^\vee}(D))
 \quad\text{(formal smoothness)}, \\
E_1^{2,1} &= H^1(\fg;\, S^2(\fg[t^{-1}]t^{-1})) = 0
 \quad\text{(Whitehead)}, \\
E_1^{1,2} &= H^2(\fg;\, S^1(\fg[t^{-1}]t^{-1})) = 0
 \quad\text{(Whitehead)}, \\
E_1^{0,3} &= H^3(\fg; k) \cong k
 \quad\text{(Cartan)}.
\end{align*}
The differentials $d_r\colon E_r^{0,3} \to E_r^{r, 4-r}$
vanish for $r = 1, 2, 3$:
\begin{enumerate}[label=\textup{(\roman*)}]
\item $d_1\colon E_1^{0,3} \to E_1^{1,3}
 = H^3(\fg;\, \fg[t^{-1}]t^{-1})
 = \bigoplus_{n \geq 1} H^3(\fg;\, \fg)$.
 By Poincar\'e duality for Lie algebra cohomology,
 $H^3(\fg;\, \fg) \cong H^{\dim\fg - 3}(\fg;\, \fg)^*$
 \textup{(}using $\fg \cong \fg^*$ via the Killing form\textup{)}.
 For $\fg = \mathfrak{sl}_2$, this gives
 $H^3(\mathfrak{sl}_2;\, \mathfrak{sl}_2) \cong H^0(\mathfrak{sl}_2;\, \mathfrak{sl}_2)^*
 = (\mathfrak{sl}_2^{\mathfrak{sl}_2})^* = 0$.
 More generally, $H^q(\fg;\, \fg) = 0$ for \emph{all}~$q$
 and all simple~$\fg$, since $\fg$ is a nontrivial irreducible
 $\fg$-module and the Casimir element acts by a nonzero scalar
 on~$\fg$ but trivially on cohomology
 \textup{(}generalized Whitehead lemma\textup{)}.
 Hence $d_1 = 0$.
\item $d_2\colon E_2^{0,3} \to E_2^{2,2}$: the target satisfies
 $E_2^{2,2} \leq E_1^{2,2} = H^2(\fg;\, S^2(\cdots)) = 0$
 \textup{(}Whitehead\textup{)}, so $d_2 = 0$.
\item $d_3\colon E_3^{0,3} \to E_3^{3,1}$: similarly
 $E_3^{3,1} \leq E_1^{3,1} = H^1(\fg;\, S^3(\cdots)) = 0$
 \textup{(}Whitehead\textup{)}, so $d_3 = 0$.
\end{enumerate}
The first potentially nonzero differential is
\begin{equation}\label{eq:d4-obstruction}
d_4\colon E_4^{0,3} = k
\;\longrightarrow\;
E_4^{4,0} \subset S^4(\fg[t^{-1}]t^{-1})^{\fg},
\end{equation}
which maps the Cartan $3$-cocycle
$\omega_3(a,b,c) = \kappa(a,[b,c])$ to the degree-$4$ invariant
subalgebra of the loop modes.
\end{proposition}

\begin{proof}
The identification of $E_1^{p,q}$ follows from the PBW filtration
as in the proofs of Theorem~\ref{thm:oper-bar-h0-dl} and
Propositions~\ref{prop:oper-bar-h1-dl}--\ref{prop:oper-bar-h2-dl}.
The Whitehead vanishing at $q = 1, 2$ and the Cartan theorem
$H^3(\fg; k) = k$ are classical. For the vanishing
$H^3(\fg;\, \fg) = 0$: by Poincar\'e duality,
$H^n(\fg; V) \cong H^{\dim\fg - n}(\fg; V^*)^*$ for any
finite-dimensional $\fg$-module~$V$ when $\fg$ is semisimple
(using the Killing form to identify
$\Lambda^{\dim\fg}(\fg^*) \cong k$). Taking $V = \fg \cong \fg^*$
gives $H^3(\fg;\, \fg) \cong H^{\dim\fg - 3}(\fg;\, \fg)^*$.
For $\mathfrak{sl}_2$ ($\dim = 3$), this reduces to $H^0(\mathfrak{sl}_2;\, \mathfrak{sl}_2)^* = 0$.
For general simple~$\fg$, the generalized Whitehead lemma gives
$H^q(\fg;\, V) = 0$ for \emph{all}~$q$ when~$V$ is a
nontrivial irreducible finite-dimensional $\fg$-module:
the Casimir $C \in Z(U\fg)$ acts on~$V$ by the nonzero scalar
$\langle \lambda, \lambda + 2\rho\rangle$ (where $\lambda$ is the
highest weight of~$V$), but acts trivially on cohomology
since it lies in the center; thus
$\langle \lambda, \lambda + 2\rho\rangle \cdot H^q(\fg; V)
= 0$ forces $H^q(\fg; V) = 0$. For $V = \fg$ (the adjoint,
with highest weight~$\theta$), this gives the claim.

The vanishing of $d_1$, $d_2$, $d_3$ then follows because each
differential maps into a subquotient of the indicated~$E_1$ term.
\end{proof}

\begin{proposition}[Non-vanishing of \texorpdfstring{$d_4$}{d4}; \ClaimStatusProvedHere]
\label{prop:d4-nonvanishing}
\index{bar complex!$d_4$ obstruction}
\index{Cartan transgression}
In the PBW spectral sequence for $\barB(\widehat{\fg}_{-h^\vee})$,
the differential
$d_4\colon E_4^{0,3} = k \to E_4^{4,0}$
of~\eqref{eq:d4-obstruction} satisfies $d_4(\omega_3) \neq 0$.
Consequently $E_5^{0,3} = 0$.
\end{proposition}

\begin{proof}
By Corollary~\ref{cor:bar-computes-ext}, the bar cohomology
computes module self-extensions:
$H^n(\barB(\widehat{\fg}_{-h^\vee})) \cong
\mathrm{Ext}^n_{\widehat{\fg}_{-h^\vee}}(V_{\mathrm{crit}},
V_{\mathrm{crit}})$.
Frenkel and Teleman~\cite{FT06} proved
$\mathrm{Ext}^3_{\widehat{\fg}_{-h^\vee}}(V_{\mathrm{crit}},
V_{\mathrm{crit}}) \cong \Omega^3(\mathrm{Op}_{\fg^\vee}(D))$
(\ClaimStatusProvedElsewhere).
Since $\mathrm{Op}_{\fg^\vee}(D)$ is formally smooth
\cite{FG06}, the cotangent complex is classical,
$\mathbb{L}_{\mathrm{Op}} = \Omega^1[0]$, and
$\Omega^3(\mathrm{Op})$ is concentrated in the summand
$E_\infty^{3,0}$. In particular
$E_\infty^{0,3} = 0$.

Now $E_4^{0,3} = E_1^{0,3} = H^3(\fg; k) = k$ (the Cartan
$3$-cocycle class), since the differentials $d_1, d_2, d_3$
vanish on $E_r^{0,3}$ by
Proposition~\ref{prop:h3-differential-analysis}. The only
differential out of $E_4^{0,3}$ is $d_4$, and subsequent
differentials from $E_r^{0,3}$ vanish for degree reasons
(targets lie in $E_r^{r, 3-r+1}$ with $3 - r + 1 < 0$ for
$r \geq 5$). Thus $E_\infty^{0,3} = \ker d_4|_{E_4^{0,3}}$.
Since $E_\infty^{0,3} = 0$ and $E_4^{0,3} = k$ is one-dimensional,
$d_4(\omega_3) \neq 0$.
\end{proof}

\begin{remark}[Transgression interpretation]
\label{rem:d4-transgression}
The non-vanishing of $d_4$ is the chain-level incarnation of the
\emph{transgression} of the Cartan $3$-cocycle
$\omega_3(a,b,c) = \kappa(a,[b,c])$ through the loop algebra:
$\omega_3$ transgresses to a $1$-cocycle on $L\fg$ valued in the
Feigin--Frenkel center~$\mathfrak{z}$, and $d_4$ captures this
transgression in the PBW spectral sequence.
\end{remark}

\begin{remark}[H3 obstruction]\label{rem:h3-obstruction}
The $d_4$ differential~\eqref{eq:d4-obstruction} is the unique
obstruction to extending the oper identification beyond~$H^2$.
Its non-vanishing (Proposition~\ref{prop:d4-nonvanishing}) shows
that $H^3(\barB(\widehat{\fg}_{-h^\vee}))$ is \emph{not} simply
$S^3$-invariants modulo the Cartan class, but requires the
Frenkel--Teleman computation to resolve.
\end{remark}

\begin{corollary}[Third cohomology at critical level; \ClaimStatusProvedHere]
\label{cor:h3-oper}
\index{oper!$H^3$ identification}
$H^3(\barB(\widehat{\fg}_{-h^\vee})) \cong
\Omega^3(\mathrm{Op}_{\fg^\vee}(D))$.
\end{corollary}

\begin{proof}
Combine Corollary~\ref{cor:bar-computes-ext} with
Frenkel--Teleman~\cite{FT06}.
\end{proof}

\begin{theorem}[Full oper differential-form identification;
\ClaimStatusProvedHere]
\label{thm:oper-bar-dl}
\index{oper!derived identification}
For all $n \geq 0$:
\begin{equation}\label{eq:hn-oper}
H^n\bigl(\barB(\widehat{\fg}_{-h^\vee})\bigr)
\;\cong\;
\bigwedge\nolimits^n
\mathbb{L}_{\mathrm{Op}_{\fg^\vee}(D)}
\;=\;
\Omega^n\bigl(\mathrm{Op}_{\fg^\vee}(D)\bigr).
\end{equation}
Equivalently, the cohomology algebra of the critical bar complex
agrees with the graded algebra of differential forms on the
formal-disc oper space.
\end{theorem}

\begin{proof}
The argument combines two results, each proved independently:
\begin{enumerate}[label=\textup{(\roman*)}]
\item \emph{Bar computes Ext.}\quad
 Corollary~\ref{cor:bar-computes-ext}: for the chiral algebra
 $\widehat{\fg}_{-h^\vee}$ (whose curvature vanishes at critical
 level by Corollary~\ref{cor:critical-level-universality}), the
 bar cohomology computes vacuum self-extensions:
 \[
 H^n\bigl(\barB(\widehat{\fg}_{-h^\vee})\bigr)
 \;\cong\;
 \mathrm{Ext}^n_{\widehat{\fg}_{-h^\vee}}
 (V_{\mathrm{crit}},\, V_{\mathrm{crit}})
 \qquad \text{for all } n \geq 0.
 \]
 This is the chiral analogue of the standard bar-resolution
 identification (Loday--Vallette~\cite{LV12}), validated in the
 chiral setting by Main Theorem~A
 (Theorem~\ref{thm:bar-cobar-isomorphism-main}).
\item \emph{Ext computes oper forms.}\quad
 Frenkel--Teleman~\cite{FT06} (\ClaimStatusProvedElsewhere):
 \[
 \mathrm{Ext}^n_{\widehat{\fg}_{-h^\vee}}
 (V_{\mathrm{crit}},\, V_{\mathrm{crit}})
 \;\cong\;
 \Omega^n\bigl(\mathrm{Op}_{\fg^\vee}(D)\bigr)
 \qquad \text{for all } n \geq 0.
 \]
\end{enumerate}
Composing (i) and (ii) gives $H^n(\barB(\widehat{\fg}_{-h^\vee}))
\cong \Omega^n(\mathrm{Op}_{\fg^\vee}(D))$ for all~$n$.
This is the same argument used in the proof of
Corollary~\ref{cor:h3-oper} at $n = 3$; the argument is $n$-independent. This proves the cohomological
identification with oper differential forms. It should not be read as
an identification of the bar complex with the derived structure sheaf
of the oper space itself; for a classically smooth oper space, that
derived structure sheaf remains concentrated in degree~$0$.
\end{proof}

\begin{remark}[Proof structure and independent confirmation]
\label{rem:oper-evidence}
\index{oper!proof structure}
The proof combines bar-Ext (Corollary~\ref{cor:bar-computes-ext})
with Frenkel--Teleman~\cite{FT06}. Independent PBW spectral
sequence confirmation is available for $n \leq 3$
(Theorem~\ref{thm:oper-bar-h0-dl},
Propositions~\ref{prop:oper-bar-h1-dl}--\ref{prop:oper-bar-h2-dl},
Corollary~\ref{cor:h3-oper}). The Whitehead spectral decomposition
(Proposition~\ref{prop:whitehead-spectral-decomposition}) extends
to all~$n$, giving
$E_1^{p,q} = \mathrm{Fun}^p(\mathrm{Op}) \otimes H^q(\fg; k)$
and reducing the spectral sequence to transgression differentials.
The result is compatible with the oper side of the
Frenkel--Gaitsgory localization picture
$\Delta\colon D(\widehat{\fg}_{-h^\vee}\text{-}\mathrm{mod})
\to D(\mathrm{QCoh}(\mathrm{Op}))$
at the cohomological level; see Theorem~\ref{thm:oper-bar-dl}.
\end{remark}

\begin{remark}[Vacuum-side compatibility with localization;
\ClaimStatusHeuristic]
\index{localization functor!bar complex}
The map
\[
\barB(\widehat{\fg}_{-h^\vee})
\;\longrightarrow\;
\mathrm{Fun}(\mathrm{Op}_{\fg^\vee}(D))
\]
given by projecting to $H^0$ is a map of algebras: the bar
complex at critical level is a dg-coalgebra (since the curvature
vanishes); its cohomology ring inherits an algebra structure, and
the projection to $H^0$ respects this multiplication. Externally, Frenkel--Gaitsgory localization sends the
vacuum module to $\cO_{\mathrm{Op}}$, so this projection is compatible
with the vacuum-side target of the localization picture. What is
\emph{not} proved internally here is that the full bar complex itself
is the localization functor, or that it provides a chain-level model
for localization on arbitrary modules.
\end{remark}

\begin{remark}[From critical localization to the general Slodowy-slice principle]
\label{rem:oper-as-kostant-slice}
\index{oper!as Kostant section}
\index{Kostant section!oper identification}
\index{Slodowy slice!oper identification}
\index{geometric localization!critical vs.\ non-critical}
Remark~\ref{rem:bar-as-localization} is not an isolated
critical-level phenomenon; it is the \emph{principal, critical-level
instance} of the general DS geometric localization principle
(Remark~\ref{rem:ds-geometric-localization-principle}).

The oper space $\mathrm{Op}_{\fg^\vee}(D)$ is the jet scheme of the
\emph{Kostant section}
$\fg^{\mathrm{reg}} / G \cong \fg^e_{\mathrm{prin}} =
S_{f_{\mathrm{prin}}} \cap \fg^{\mathrm{reg}}$,
where $f_{\mathrm{prin}}$ is the principal nilpotent. At critical
level $k = -h^\vee$, the curvature vanishes, and the bar complex
becomes a genuine dg algebra whose $H^0$ is functions on the oper
space, i.e., functions on $J_\infty(S_{f_{\mathrm{prin}}})$. The
Frenkel--Gaitsgory localization functor~$\Delta$ is heuristically the
critical-level specialization of the Hamiltonian reduction from
$J_\infty(\fg^*)$ to $J_\infty(S_f)$: at critical level, the
Hamiltonian reduction is ``free'' (the moment-map fiber is smooth),
so derived = underived, and $H^0$ suffices.

At non-critical level $k \neq -h^\vee$, the genus-$0$ bar complex
remains uncurved ($d_{\mathrm{bar}}^2 = 0$ always, by Theorem~A),
but the fiberwise bar differential acquires genus-$1$ curvature
$d_{\mathrm{fib}}^2 = \kappa(\widehat{\fg}_k) \cdot \omega_1$, where $\kappa(\widehat{\fg}_k) = \dim(\fg)\,(k+h^\vee)/(2h^\vee)
\neq 0$
(Theorem~\ref{thm:modular-characteristic}).
The nonvanishing fiberwise curvature has two consequences for
localization. First, the associated graded of the localization
functor (the Hamiltonian reduction $J_\infty(\fg^*) \to
J_\infty(S_f)$) is no longer sufficient: higher-genus
corrections propagate through the modular bar coalgebra and the
localization must be understood in the derived/curved sense.
Second, the derived DS functor on modules acquires higher
BRST cohomology at special levels (admissible, critical
subquotients), so the single-term $H^0_{\mathrm{DS}}$ no longer
captures the full localization.
Theorem~\ref{thm:ds-koszul-intertwine} provides the chain-level
model: the double complex
$(\barB(M) \otimes F_\chi,\, d_{\mathrm{bar}},\, Q_{\mathrm{DS}})$
is the derived localization of the bar complex to the Slodowy-slice
jet scheme, and its spectral sequence
(Theorem~\ref{thm:ds-koszul-intertwine}, Step~3) computes the
derived fibers.

The master commutative square
(Theorem~\ref{thm:master-commutative-square})
then places Remark~\ref{rem:bar-as-localization} in its natural
home: the oper
differential-form package of the Kostant section is classified by the critical-level bar complex, and DS reduction at
any level restricts from the
ambient jet scheme to the slice. Opers are the principal Slodowy
slice, and localization is Hamiltonian reduction.
\end{remark}


\begin{remark}[The MC element $\Theta_{\cA}$ at critical level]
\label{rem:critical-level-theta}
\ClaimStatusProvedHere
\index{Maurer--Cartan element!critical level}
\index{critical level!Theta@$\Theta_{\cA}$}
The critical level $k = -h^\vee$ is the unique point in the
Kac--Moody level space where the modular characteristic vanishes:
$\kappa(\widehat{\fg}_{-h^\vee}) = \dim(\fg)\,(k + h^\vee)/(2h^\vee)\big|_{k=-h^\vee} = 0$
%: formula from CLAUDE.md C3 / landscape_census.tex:1102; k=0 -> dim(g)/2; k=-h^v -> 0 verified
(Corollary~\ref{cor:critical-level-universality}).
This vanishing has two consequences for the bar complex that must be
carefully distinguished.

\emph{First consequence: uncurving.} Since $\kappa(\widehat{\fg}_{-h^\vee}) = 0$, the genus-$0$
curvature $m_0$ vanishes and the bar complex
$\barB(\widehat{\fg}_{-h^\vee}) = T^c(s^{-1}\,\overline{V}_{-h^\vee})$
%: augmentation ideal
is an honest chain complex ($d^2 = 0$). The higher-genus fiberwise
obstruction $d_{\mathrm{fib}}^2 = \kappa(\widehat{\fg}_k) \cdot \omega_g$ likewise
vanishes.

\emph{Second consequence: the $r$-matrix survives.} The classical
$r$-matrix is $r^{\mathrm{KM}}(z) = \Omega/\bigl((k{+}h^\vee)z\bigr)$
% KZ normalisation; diverges at critical level (Sugawara singularity).
% In the trace-form convention: r = k*Omega_tr/z, finite at k=-h^v.
(Remark~\ref{rem:km-collision-residue-rmatrix}). In the
trace-form normalisation $r(z) = k\,\Omega_{\mathrm{tr}}/z$, the
critical level specializes to $r(z) = -h^\vee \cdot \Omega_{\mathrm{tr}}/z$; this is
\emph{nonzero}. The vanishing $\kappa(\widehat{\fg}_{-h^\vee}) = 0$ means only that the
$\Sigma_n$-coinvariant projection $\mathrm{av}(r(z)) = 0$ vanishes
(the averaging map kills $r$ without $r$ itself being zero). In
particular, the degree-$2$ component
$\Theta_{\cA}^{(0,2)} = r(z) = -h^\vee \cdot \Omega_{\mathrm{tr}}/z \neq 0$
persists.

This is the key structural distinction between the critical level
$k = -h^\vee$ and the abelian limit $k = 0$
(Remark~\ref{rem:two-critical-points}). At $k = 0$:
the $r$-matrix $r(z) = k\,\Omega/z\big|_{k=0} = 0$ vanishes
identically (while $\kappa(\widehat{\fg}_0) = \dim(\fg)/2 \neq 0$),
so $\Theta_{\cA}$ loses its degree-$2$ component and reduces to the purely
higher-degree tail. At $k = -h^\vee$: $\kappa(\widehat{\fg}_{-h^\vee}) = 0$ but
$r(z) \neq 0$, and $\Theta_{\cA}$ retains its full
structure at all degrees $\geq 2$.

For $\fg = \mathfrak{sl}_2$ ($h^\vee = 2$,
$\dim(\fg) = 3$,
$\kappa(V_k(\mathfrak{sl}_2)) = 3(k+2)/4$),
%: census check; k=0 -> 3/2; k=-2 -> 0 verified
the critical level is $k = -2$, and the Feigin--Frenkel center
$\mathfrak{z}(\widehat{\mathfrak{sl}}_2) \cong
\mathrm{Fun}(\mathrm{Op}_{\mathrm{PGL}_2}(D))$
is a polynomial algebra in one generator: the Segal--Sugawara
operator~$S(z)$
(Theorem~\ref{thm:ff-center-dl}).
The $r$-matrix $r(z) = -2\,\Omega/z$ is the degree-$2$ shadow of
the full MC element~$\Theta_{\cA}$, and
the oper differential-form identification
(Theorem~\ref{thm:oper-bar-dl}) shows that the bar
cohomology $H^*(\barB(V_{-2}(\mathfrak{sl}_2)))$
recovers $\Omega^*(\mathrm{Op}_{\mathrm{PGL}_2}(D))$.
The polynomial structure of~$\mathfrak{z}$ (one generator, one
relation at each conformal weight) constrains the higher-degree
components of~$\Theta_{\cA}$: they encode the $A_\infty$
operations on the bar complex that are compatible with the
oper geometry.
\end{remark}

\begin{remark}[Koszul duality at critical level: bar complex as oper de Rham complex]
\label{rem:bar-critical-oper-deRham}
\ClaimStatusConjectured
\index{Koszul duality!critical level}
\index{oper!de Rham complex}
\index{Feigin--Frenkel center!Koszul dual}
The results of this section admit a clean reformulation in the
language of bar-cobar Koszul duality
(Theorem~\ref{thm:bar-cobar-inversion-qi}).

At the critical level $k = -h^\vee$, the modular characteristic
$\kappa(\widehat{\fg}_{-h^\vee})
= \dim(\fg)\,(k + h^\vee)/(2h^\vee)\big|_{k=-h^\vee} = 0$
%: from CLAUDE.md C3; k=0 -> dim(g)/2; k=-h^v -> 0 verified
vanishes (Corollary~\ref{cor:critical-level-universality}), so the
bar complex $\barB(\widehat{\fg}_{-h^\vee})
= T^c(s^{-1}\,\overline{V}_{-h^\vee})$
%: augmentation ideal
is uncurved ($m_0 = 0$, $d^2 = 0$).
Theorem~\ref{thm:oper-bar-dl} identifies its cohomology algebra
with the differential-form algebra on the oper space:
\[
H^*(\barB(\widehat{\fg}_{-h^\vee}))
\;\cong\;
\Omega^*(\mathrm{Op}_{\fg^\vee}(D)).
\]
This identification has a Koszul-theoretic interpretation.
The Koszul dual algebra
$\widehat{\fg}_{-h^\vee}^!
= \mathrm{CE}^{\mathrm{ch}}(\fg_{-h^\vee})$
is the curved chiral Chevalley--Eilenberg algebra; its zeroth
cohomology is the Feigin--Frenkel center
$\mathfrak{z}(\widehat{\fg}_{-h^\vee})
\cong \mathrm{Fun}(\mathrm{Op}_{\fg^\vee}(D))$
(Theorem~\ref{thm:ff-center-dl}).
Bar-cobar inversion
$\Omega(\barB(\widehat{\fg}_{-h^\vee}))
\xrightarrow{\;\sim\;} \widehat{\fg}_{-h^\vee}$
(Theorem~\ref{thm:bar-cobar-inversion-qi})
recovers the original algebra from its bar coalgebra; dualizing
the bar coalgebra instead produces the Koszul dual algebra
$\widehat{\fg}_{-h^\vee}^!$, whose degree-zero piece is
$\mathfrak{z}(\widehat{\fg}_{-h^\vee})$.
In this sense, the oper de Rham complex
$\Omega^*(\mathrm{Op}_{\fg^\vee}(D))$ is simultaneously:
\begin{enumerate}[label=\textup{(\roman*)}]
\item the bar cohomology of $\widehat{\fg}_{-h^\vee}$
 (Theorem~\ref{thm:oper-bar-dl}),
\item the underlying graded algebra of the Koszul dual
 $\widehat{\fg}_{-h^\vee}^!$
 (via $H^0 = \mathfrak{z}$ and exterior powers of the
 cotangent complex of the oper space).
\end{enumerate}
The conjectural refinement is that this coincidence lifts to a
chain-level quasi-isomorphism of dg algebras: the bar complex
$\barB(\widehat{\fg}_{-h^\vee})$, viewed as a dg coalgebra
with its deconcatenation coproduct, should be quasi-isomorphic
as a dg coalgebra to the de Rham complex
$\Omega^*(\mathrm{Op}_{\fg^\vee}(D))$ equipped with the
shuffle coproduct.
The critical level is special in two ways: the vanishing of
$\kappa(\widehat{\fg}_{-h^\vee})$ uncurves the bar complex (making the dg coalgebra
structure honest), and the critical level is self-dual under
Koszul duality ($-(-h^\vee) - 2h^\vee = -h^\vee$; see
Remark~\ref{rem:km-complementarity-vs-feigin-frenkel}).
\end{remark}

\begin{remark}[Hecke categorical action and chiral homology compatibility]
\label{rem:hecke-chiral-compatibility}
\ClaimStatusConjectured
\index{Hecke!categorical action}
\index{Hecke!eigensheaves}
\index{Langlands!geometric}
\index{affine monodromy!Hecke correspondence}
The bar--oper identification of
Remark~\ref{rem:bar-critical-oper-deRham} situates itself
inside the geometric Langlands correspondence, and we record
the compatibility with the Hecke categorical action on
$\mathrm{D}(\mathrm{Bun}_G(X))$.
\begin{enumerate}[label=\textup{(\roman*)}]
\item \emph{Hecke--chiral commutation.}
The Hecke functors $H_\lambda$ acting on
$\mathrm{D}(\mathrm{Bun}_G(X))$ commute with the chiral
homology functor $H^*_{\mathrm{ch}}(X, V_k(\fg))$ of the
affine vacuum algebra at level $k$
(Beilinson--Drinfeld 2004, Ch.~7). This realizes the chiral
bar complex computed in this chapter as a Hecke-equivariant
object.
\item \emph{Hecke eigensheaves vs.\ $\mathrm{Rep}(G^\vee)$.}
Under the geometric Langlands equivalence, Hecke eigensheaves
on $\mathrm{Bun}_G$ correspond to objects of
$\mathrm{Rep}(G^\vee)$. At generic quantum parameter these
are the Kazhdan--Lusztig quantum deformations
$\mathrm{Rep}_q(G^\vee)$; at a root of unity one passes to
the modular tensor category quotient, in line with the
Reshetikhin--Turaev functor invoked by the affine
monodromy / spectral $R$-matrix identification of
Theorem~\ref{thm:kl-bar-cobar-adjunction}.
\item \emph{Critical-level eigenvalue locus.}
At $k = -h^\vee$, the Feigin--Frenkel center
$\mathfrak{z}(\widehat{\fg}_{-h^\vee})
\cong \mathrm{Fun}(\mathrm{Op}_{\fg^\vee}(D))$
(Theorem~\ref{thm:ff-center-dl}) is precisely the locus of
Hecke eigenvalues for the categorical action: opers
parametrize Hecke eigensheaves on $\mathrm{Bun}_G$, with the
eigenvalue map supplied by the Feigin--Frenkel isomorphism.
\item \emph{Programme consequence.}
The affine monodromy identification
(decorator~\#10,
Theorem~\ref{thm:affine-monodromy-identification})
corresponds, under Langlands-dual Koszul duality, to the
Hecke eigensheaf correspondence: the spectral $R$-matrix
encoded in $\barB(\widehat{\fg}_k)$ at generic level is the
Koszul-dual shadow of the Hecke categorical action, and
specializes at $k = -h^\vee$ to the oper--eigenvalue
dictionary above.
\end{enumerate}
Full chain-level compatibility of the Hecke action with the
chiral bar complex beyond cohomology is conjectural; see
Conjecture~\ref{conj:periodic-cdg} for the admissible-level
refinement.
\end{remark}
%================================================================
% SECTION 4: ROOT-OF-UNITY BAR COMPLEXES
%================================================================

\section{Root-of-unity bar complexes}
\label{sec:root-of-unity}
\index{root of unity!bar complex|textbf}
\index{CDG algebra!periodic}

\subsection{Admissible levels and the quantum group parameter}

\begin{definition}[Admissible level]
\label{def:admissible-level-dl}
A level $k$ is \emph{admissible} if
$k = -h^\vee + p/q$ where $p, q \in \bZ_{> 0}$, $\gcd(p, q) = 1$,
and $p \geq h^\vee$ (for simply-laced~$\fg$; the condition is more
involved for non-simply-laced types). The associated quantum group
parameter is $q_{\mathrm{QG}} = e^{\pi i/(k + h^\vee)} = e^{\pi i q/p}$.
\end{definition}

At admissible level, the bar complex has curvature
$m_0 = k \cdot (\,\cdot\,,\,\cdot\,) \cdot |0\rangle$,
where $(\,\cdot\,,\,\cdot\,)$ is the normalized Killing form,
which is non-zero but \emph{rational}.
This curvature arises directly from the level-$k$ central extension:
the bilinear form $k(a, b)$ appears as the double-pole residue
$a_{(1)}b = k(a,b)|0\rangle$ in the OPE, and obstructs
$d^2 = 0$ on the unreduced bar complex
(Remark~\ref{rem:bicomplex-obstruction}).
The related Sugawara central charge $c = k\dim(\fg)/(k+h^\vee)$
is a derived quantity, not the direct source of curvature.
Meanwhile, the representation category $\cO_k^{\mathrm{int}}(\widehat{\fg})$
is semisimple with $|\mathrm{Adm}_k|$ simple objects
(Theorem~\ref{thm:admissible-rep-theory}).

\subsection{The periodic CDG structure}

\begin{definition}[Curved dg-algebra (CDG)]
\label{def:cdg-algebra}
\index{CDG algebra|textbf}
A \emph{curved differential graded algebra} (CDG-algebra) is a
triple $(A, d, h)$ where $A$ is a graded algebra,
$d\colon A \to A[1]$ is a derivation, and $h \in A^2$ is the
\emph{curvature element}, satisfying:
\begin{enumerate}[label=\textup{(\roman*)}]
\item $d^2(a) = [h, a]$ for all $a \in A$;
\item $dh = 0$ (the curvature is closed).
\end{enumerate}
The bar complex $\barB(\widehat{\fg}_k)$ at non-critical level is
a CDG-coalgebra: the curvature element $m_0$ satisfies $d(m_0) = 0$
and $d^2 = [m_0, -]$
(Proposition~\ref{prop:pole-decomposition}).
\end{definition}

\begin{conjecture}[Periodic CDG structure at admissible level;
\ClaimStatusConjectured]
\label{conj:periodic-cdg}
\index{CDG algebra!periodicity}
At admissible level $k = -h^\vee + p/q$, the bar complex
$\barB(\widehat{\fg}_k)$ admits a \emph{periodic CDG structure}:
there exists a natural isomorphism of CDG-coalgebras
\begin{equation}\label{eq:periodic-cdg}
\barB^{n + 2s}(\widehat{\fg}_k)
\;\cong\;
\barB^n(\widehat{\fg}_k)
\end{equation}
where $s = p$ \textup{(}the numerator of
$k + h^\vee = p/q$\textup{)}, so the bar complex is $2p$-periodic
in bar degree.
\end{conjecture}

\begin{remark}[Positselski's framework]
\label{rem:positselski}
\index{Positselski!CDG algebras}
Positselski~\cite{Positselski11} develops CDG-algebras whose
Koszul duality exchanges coderived and contraderived categories:
$D^{\mathrm{co}}(A) \simeq D^{\mathrm{ctr}}(B)$. For periodic
CDG-algebras, the coderived category reduces to the stable category
of a finite-dimensional quotient (the passage from $U_q(\fg)$ to
$u_q(\fg)$). Conjecture~\ref{conj:periodic-cdg} asserts that the
bar complex at admissible level is periodic in this sense; upgrading
it requires understanding the quantum group center action on the bar
complex at a root of unity.
(Contributing to Conjecture~\ref{conj:master-dk-kl}.)
\end{remark}


\subsection{Evidence for periodicity}

\begin{proposition}[Affine sl2 chiral Hochschild periodicity at critical level;
\ClaimStatusProvedElsewhere]
\label{prop:sl2-periodicity-dl}
\index{periodicity!$\widehat{\mathfrak{sl}}_2$}
For $\fg = \mathfrak{sl}_2$ at critical level $k = -h^\vee = -2$, the
chiral Hochschild cohomology satisfies the periodicity
\begin{equation}\label{eq:sl2-periodicity}
\dim \ChirHoch^{n+4}(\widehat{\mathfrak{sl}}_{2,-2})
= \dim \ChirHoch^n(\widehat{\mathfrak{sl}}_{2,-2})
\quad
\text{for all } n \geq 0,
\end{equation}
with period $2h = 4$ \textup{(}$h = 2$ for $\mathfrak{sl}_2$\textup{)}.
The periodicity operator is cup product with the even generator
$\Theta_1 \in \ChirHoch^4$.
\end{proposition}

\begin{proof}
Theorem~\ref{thm:affine-periodicity-critical}: at critical level, the
bar complex is uncurved and the Beilinson--Drinfeld comparison gives
$\ChirHoch^*(\widehat{\mathfrak{sl}}_{2,-2}) \cong
\Lambda(P_1) \otimes \bC[\Theta_1]$ with $\deg P_1 = 3$,
$\deg \Theta_1 = 4$. Cup product with $\Theta_1$ is an isomorphism
$\ChirHoch^n \xrightarrow{\sim} \ChirHoch^{n+4}$ for all $n \geq 0$.
\end{proof}

\begin{remark}[Bar cohomology vs chiral Hochschild cohomology]
\label{rem:bar-vs-hochschild-periodicity}
The bar cohomology dimensions $\dim H^n(\barB(\widehat{\mathfrak{sl}}_{2,k}))$
at generic level grow linearly \textup{(}$3, 5, 7, 9, 11, 13, \ldots = 2n{+}1$
for $n = 1, 2, \ldots$; Remark~\textup{\ref{rem:garland-lepowsky-sl2}}\textup{)} and are \emph{not} periodic.
The periodicity of Proposition~\ref{prop:sl2-periodicity-dl} is a
critical-level, chiral-Hochschild-cohomology phenomenon: at $k = -2$, the bar
complex is uncurved and the Feigin--Frenkel center provides the
polynomial generator $\Theta_1$ whose cup product gives the period-$4$
isomorphism. The extension to admissible levels requires the periodic CDG
structure \textup{(}Conjecture~\ref{conj:periodic-cdg}\textup{)}.
\end{remark}

\begin{remark}[Periodicity and quantum group]
\label{rem:periodicity-quantum}
\index{quantum group!from periodicity}
At admissible level $k = -2 + p/q$, the period $2h = 4$ divides
$2p$ when $p$ is even. The $4$-periodicity of the chiral Hochschild
cohomology at critical level is a shadow of the representation-theoretic
periodicity of the quantum group $U_q(\mathfrak{sl}_2)$ at
$q = e^{\pi i q/p}$. For the small quantum group $u_q$, the
representation category is $2p$-periodic, and Conjecture~\ref{conj:periodic-cdg} asserts
that the bar complex ``sees'' this periodicity via CDG structure.
\end{remark}

\begin{remark}[Scope boundary for periodicity]
\label{rem:periodicity-scope-boundary-dl}
Only rank-$1$ strict chiral Hochschild periodicity at critical level is proved in
this chapter
\textup{(}Proposition~\ref{prop:sl2-periodicity-dl}\textup{)}.
For higher-rank affine algebras, strict periodicity of $\ChirHoch^*$ fails
\textup{(}polynomial growth $O(n^{r-1})$, not periodic; see
Theorem~\ref{thm:affine-periodicity-critical}\textup{)}.
Extension to admissible levels requires the periodic CDG conjecture
\textup{(}Conjecture~\ref{conj:periodic-cdg}\textup{)}.
In particular, periodicity does not transfer automatically across Koszul
duality or level-shift involutions without additional rank-sensitive input.
\end{remark}


%================================================================
% SECTION 5: THE KAZHDAN-LUSZTIG EQUIVALENCE FROM BAR-COBAR
%================================================================

\section{Kazhdan--Lusztig from bar-cobar: chain-level adjunction and periodic structure}
\label{sec:kl-from-bar-cobar}
\index{Kazhdan--Lusztig equivalence!from bar-cobar|textbf}

\begin{theorem}[Kazhdan--Lusztig equivalence on the semisimplified target;
\ClaimStatusProvedElsewhere]
\label{thm:kl-equivalence}
\index{Kazhdan--Lusztig equivalence}
\textup{(}Kazhdan--Lusztig~\cite{KL93}\textup{)}\quad
For $\fg$ simply-laced and $k = -h^\vee + p/q$ admissible with
$p \geq h^\vee$, there is an equivalence of braided monoidal
categories:
\begin{equation}\label{eq:kl-equivalence-dl}
\cO_k^{\mathrm{int}}(\widehat{\fg})
\;\simeq\;
\mathcal{C}(U_q(\fg)),
\end{equation}
where $q = e^{\pi i q/p}$ and $\mathcal{C}(U_q(\fg))$ denotes the
semisimplified tilting quotient on the admissible alcove. The larger
category $\mathrm{Rep}^{\mathrm{fd}}(U_q(\fg))$ is a non-semisimple
lift, not the theorematic target used in this monograph.
\end{theorem}

\begin{theorem}[Chain-level KL adjunction from bar-cobar;
\ClaimStatusProvedHere]
\label{thm:kl-bar-cobar-adjunction}
\index{Kazhdan--Lusztig equivalence!bar-cobar adjunction}
Let $\fg$ be simply-laced and $k = -h^\vee + p/q$ admissible with
$p \geq h^\vee$. Write
\[
 C_k := \barB^{\mathrm{ch}}(\widehat{\fg}_k).
\]
There exists a chain-level adjunction
\begin{equation}\label{eq:kl-bar-cobar-adjunction}
\Phi_k\colon
D\bigl(\mathrm{Mod}^{\Eone,\,\mathrm{compl}}_{\widehat{\fg}_k}\bigr)
\;\rightleftarrows\;
D\bigl(\mathrm{CoMod}^{\Eone,\,\mathrm{conil}}_{C_k}\bigr)
\;\colon \Psi_k
\end{equation}
such that:
\begin{enumerate}[label=\textup{(\roman*)}]
\item On integrable highest-weight modules $M(\lambda) \in
 \cO_k^{\mathrm{int}}(\widehat{\fg})$, the functor $\Phi_k$
 produces a finite-dimensional comodule over the dual-side bar
 coalgebra $C_k$.
\item The induced map on Ext groups recovers BGG multiplicities:
 \[
 \mathrm{Ext}^n_{\widehat{\fg}_k}(\omega_X, M(\lambda))
 \cong
 \operatorname{Hom}_{D(\mathrm{CoMod}^{\Eone,\,\mathrm{conil}}_{C_k})}
 \bigl(\Phi_k(\omega_X), \Phi_k(M(\lambda))[n]\bigr).
 \]
\item The composition $\Psi_k \circ \Phi_k$ is quasi-isomorphic
 to the identity on $\cO_k^{\mathrm{int}}(\widehat{\fg})$.
\end{enumerate}
\end{theorem}

\begin{proof}
The proof assembles five ingredients, each proved elsewhere in this
monograph, into a single chain-level construction.

\emph{Step 1: Koszul dual identification.}
By Theorem~\ref{thm:universal-kac-moody-koszul}, the chiral Koszul
dual of $\widehat{\fg}_k$ is $\widehat{\fg}_{-k-2h^\vee}$.
Set $\cA = \widehat{\fg}_k$ and
$C_k := \barB^{\mathrm{ch}}(\widehat{\fg}_k)$ for the ordered
genus-$0$ bar coalgebra on the dual side. On finite-type lanes one may
rewrite $C_k^\vee$ as the dual algebra at level $-k-2h^\vee$, but the
intrinsic module comparison is coalgebra-valued.

\emph{Step 2: Module bar-cobar adjunction.}
By Theorem~\ref{thm:e1-module-koszul-duality}, the $\Eone$-chiral
bar-cobar adjunction gives an equivalence
\[
D(\mathrm{Mod}^{\Eone,\,\mathrm{compl}}_{\cA})
\;\simeq\;
D(\mathrm{CoMod}^{\Eone,\,\mathrm{conil}}_{C_k}).
\]
Define $\Phi_k := \mathrm{Ind}$ and $\Psi_k := \mathrm{Res}$ from
that theorem. The unit and counit are quasi-isomorphisms on the
relevant subcategories, giving~(iii).

\emph{Step 3: PBW filtration and chain map.}
Equip $\widehat{\fg}_k$ with the PBW filtration by conformal weight.
This filtration is compatible with the bar construction
(Theorem~\ref{thm:bar-cobar-spectral-sequence}): the associated
graded of the bar complex $\barB(\widehat{\fg}_k)$ is the
\emph{classical} bar complex $\barB(U(\fg[t^{-1}]t^{-1}))$, which
is a dg-coalgebra (not curved). For a highest-weight module
$M(\lambda)$, the PBW filtration induces a filtration on the
relative bar complex
$\barB(\cA, M(\lambda)) := \barB(\cA) \otimes_{\cA} M(\lambda)$.

\emph{Step 4: Spectral sequence comparison.}
The PBW filtration gives a spectral sequence
\[
E_1^{p,q} = H^{p+q}(\mathrm{gr}^p \barB(\cA, M(\lambda)))
\;\Longrightarrow\;
H^{p+q}(\barB(\cA, M(\lambda))).
\]
On the associated graded, the bar complex reduces to the classical
Chevalley--Eilenberg complex of $\fg[t^{-1}]t^{-1}$ with
coefficients in $\mathrm{gr}\, M(\lambda)$. At $E_1$, the
cohomology is computed by the BGG resolution
(Corollary~\ref{cor:bgg-koszul-involution}).
By Corollary~\ref{cor:bar-admissible-finiteness}, the relevant
fixed conformal-weight sectors of the bar complex are
finite-dimensional at admissible levels, so the spectral
sequence converges weightwise. Since $M(\lambda) \in \cO_k^{\mathrm{int}}$ is
integrable, the $E_1$ page is concentrated in finitely many degrees
(bounded by the Weyl group), and the spectral sequence degenerates at
$E_2$ by the Koszul purity of $\cA$
(Theorem~\ref{thm:km-chiral-koszul}). This establishes that
$\Phi_k(M(\lambda))$ is a finite-dimensional $C_k$-comodule,
proving~(i).

\emph{Step 5: Ext recovery.}
The derived-Hom exchange~(ii) is
Theorem~\ref{thm:e1-module-koszul-duality}(ii) applied to
$\cA = \widehat{\fg}_k$: the bar-cobar adjunction identifies module
Ext groups with derived morphisms between the corresponding dual-side
bar-comodules.
\end{proof}

\begin{conjecture}[KL equivalence from periodic CDG bar-cobar;
\ClaimStatusConjectured]
\label{conj:kl-bar-cobar-dl}
\index{Kazhdan--Lusztig equivalence!bar-cobar proof}
The chain-level adjunction of
Theorem~\textup{\ref{thm:kl-bar-cobar-adjunction}} refines to the
package-level Kazhdan--Lusztig recovery on the semisimplified target
\begin{equation}\label{eq:kl-bar-cobar-dl}
\barB_{\Eone}\bigl(\widehat{\fg}_k\text{-}\mathrm{mod}\bigr)
\;\simeq\;
\Omega_{\Eone}\bigl(U_q(\fg)\text{-}\mathrm{comod}\bigr),
\end{equation}
identifying
$\cO_k^{\mathrm{int}}(\widehat{\fg}) \simeq
\mathcal{C}(U_q(\fg))$
as braided monoidal categories, where $\mathcal{C}(U_q(\fg))$ is the
semisimplified tilting quotient (not the full representation category
$\mathrm{Rep}(U_q(\fg))$; see Remark~\ref{rem:kl-target-category}).
Any non-semisimple lift to
$\mathrm{Rep}^{\mathrm{fd}}(U_q(\fg))$ is an additional open problem,
not part of the theorematic endpoint used elsewhere in the manuscript;
the completed/coderived MC3 enlargement
\textup{(}categorical CG closure plus the generated-core DK comparison
proved for all simple types; post-CG completion open beyond
type~$A$; Corollary~\ref{cor:mc3-all-types}\textup{)} lies further out still.
The remaining ingredient beyond
Theorem~\textup{\ref{thm:kl-bar-cobar-adjunction}} is the periodic
CDG structure of the bar complex
\textup{(}Conjecture~\textup{\ref{conj:periodic-cdg}):} given
$2p$-periodicity, Positselski's Koszul duality for periodic
CDG-algebras~\cite{Positselski11} identifies the coderived category
of the bar complex with the stable module category of the small
quantum group $u_q(\fg)$, which maps to
$\mathcal{C}(U_q(\fg))$ via Lusztig's theorem.
\end{conjecture}

\begin{remark}[KL target category]\label{rem:kl-target-category}
\index{Kazhdan--Lusztig equivalence!target category}
The correct KL target is the semisimplified tilting quotient
$\mathcal{C}(U_q(\fg))$ ($|P_k^+|$ simple objects), not the full
$\mathrm{Rep}(U_q(\fg))$; see~\cite{BK01} and~\cite{EGNO},
Remark~11.7.10. The coderived category enters only as the ambient
bar-side formalism for the conjectural lift, not as the printed KL
endpoint.
(Contributing to Conjecture~\ref{conj:master-dk-kl}.)
\end{remark}

\begin{remark}[Scope]\label{rem:kl-bar-cobar-scope}
Theorem~\ref{thm:kl-bar-cobar-adjunction} constructs the chain-level
adjunction on integrable modules. Two further passages are separate:
(a)~identifying the bar functor target with $\mathcal{C}(U_q(\fg))$
requires the periodic CDG structure
(Conjecture~\ref{conj:periodic-cdg}); (b)~any lift beyond the
semisimplified target is an outer MC3 problem
\textup{(}categorical CG closure plus the generated-core DK comparison
proved for all simple types; post-CG completion open beyond
type~$A$; Corollary~\ref{cor:mc3-all-types}\textup{)}. Without periodicity,
$\Phi_k$ lands in $\widehat{\fg}_{-k-2h^\vee}$-comodules, not
$U_q(\fg)$-comodules.
(Contributing to Conjecture~\ref{conj:master-dk-kl}.)
\end{remark}

\begin{remark}[Why this would be a geometric proof]
\label{rem:geometric-proof-kl}
\index{Kazhdan--Lusztig equivalence!geometric proof}
The original KL proof~\cite{KL93} uses analytic continuation.
A bar-cobar proof would be geometric: (1)~configuration space
integrals build the bar complex; (2)~Verdier duality on FM
compactifications realizes
$\widehat{\fg}_k^! = \mathrm{CE}^{\mathrm{ch}}(\widehat{\fg}_{-k-2h^\vee})$; (3)~bar curvature
interacts with $q = e^{\pi i/(k+h^\vee)}$ to produce periodic CDG
structure; (4)~Positselski's CDG Koszul duality plus Lusztig's
stable-to-tilting comparison reach the semisimplified target.
\end{remark}


\subsection{Roadmap: from critical level to KL}

\begin{remark}[The critical-admissible deformation]
\label{rem:critical-admissible}
\index{critical level!deformation to admissible}
As $k$ deforms from $-h^\vee$ to $-h^\vee + p/q$, the bar complex
interpolates: uncurved at critical level (computes oper functions,
Theorem~\ref{thm:oper-bar-h0-dl}); CDG-coalgebra at generic
perturbation; conjecturally periodic CDG at rational $p/q$
(recovering the semisimplified KL target). This connects geometric
Langlands (opers) to representation-theoretic Langlands (quantum
groups) through bar-cobar duality.
\end{remark}


%================================================================
% SECTION 6: CONNECTIONS TO GEOMETRIC LANGLANDS
%================================================================

\section{Connections to the geometric Langlands programme}
\label{sec:langlands-connections}
\index{geometric Langlands!connections}

\subsection{The Frenkel--Gaitsgory programme}

\begin{theorem}[Frenkel--Gaitsgory localization;
\ClaimStatusProvedElsewhere]
\label{thm:fg-localization}
\index{Frenkel--Gaitsgory!localization}
\textup{(}Frenkel--Gaitsgory~\cite{FG06}\textup{)}\quad
There is a localization functor
\[
\Delta\colon
D\bigl(\widehat{\fg}_{-h^\vee}\text{-mod}\bigr)
\;\longrightarrow\;
D\bigl(\mathrm{QCoh}(\mathrm{Op}_{\fg^\vee}(D))\bigr)
\]
that is an equivalence on the subcategory of modules on which the
center $Z(\widehat{\fg}_{-h^\vee})$ acts locally finitely.
\end{theorem}

\begin{remark}[Bar-cobar as localization]
\label{rem:bar-as-localization}
The bar complex is a chain-level model for the localization
functor~$\Delta$: the vacuum module maps to $\cO_{\mathrm{Op}}$
(Theorem~\ref{thm:oper-bar-h0-dl}), a Verma module $M_\lambda$ maps
to the skyscraper at the corresponding oper, and the full derived
structure is computed by Theorem~\ref{thm:oper-bar-dl}.
\end{remark}

\subsection{From opers to Hecke eigensheaves}

\begin{remark}[The Langlands path]
\label{rem:langlands-path}
\index{geometric Langlands!path from bar complex}
The Frenkel--Gaitsgory programme factors geometric Langlands through
opers:
\[
\begin{tikzcd}
D(\widehat{\fg}_{-h^\vee}\text{-mod})
 \arrow[r, "\Delta"]
 \arrow[dr, "\text{Langlands}"']
& D(\mathrm{QCoh}(\mathrm{Op}))
 \arrow[d, "\text{forget reduction}"] \\
& D(\mathrm{QCoh}(\mathrm{Loc}_{G^\vee}))
\end{tikzcd}
\]
The top arrow $\Delta$ is computed at the chain level by the bar complex
(Theorem~\ref{thm:oper-bar-dl}). The vertical arrow is the
$G^\vee/B^\vee$-fibration forgetting the Borel reduction.
\end{remark}

At integral level, the bar complex already carries the convolution-type
coalgebra operations, Koszul duality, and ordered-to-symmetric shadow
projections that one expects on the geometric Satake side. This does
not identify the Satake category inside the bar package, but it does
single out a concrete algebraic mechanism that should interface with
the affine Grassmannian and its fiber functor.

\begin{conjecture}[Bar complex, Satake, and the shadow tower;
\ClaimStatusConjectured]
\label{conj:bar-satake-shadow}
At integral level $k$:
\begin{enumerate}[label=\textup{(\arabic*)}]
\item The reduced bar coalgebra $\barB\bigl(V_k(\fg)\bigr)$ should
 admit a $t$-structure whose heart recovers the Satake category
 $\operatorname{Perv}_{G^\vee}\bigl(\operatorname{Gr}_G\bigr)$, with
 the bar differential encoding the convolution monoidal structure.
\item The Koszul dual algebra $V_k(\fg)^!$ should be identified with a
 chiral algebra on the Langlands-dual side, concretely
 $V_{k^\vee}(\fg^\vee)$ at the dual level determined by the
 complementarity relation
 \[
 \kappa\bigl(V_k(\fg)\bigr) +
 \kappa\bigl(V_{k^\vee}(\fg^\vee)\bigr) = 0,
 \]
 at least for simply-laced $\fg$ where $\fg \cong \fg^\vee$.
\item The shadow-tower projections
 \[
 \operatorname{av}^{(r)} \colon
 \barB^{\mathrm{ord},r}\bigl(V_k(\fg)\bigr)
 \longrightarrow
 \barB^{\Sigma,r}\bigl(V_k(\fg)\bigr)
 \]
 should encode the filtration on the Satake fiber functor
 \[
 F \colon
 \operatorname{Perv}_{G^\vee}\bigl(\operatorname{Gr}_G\bigr)
 \longrightarrow
 \operatorname{Vect},
 \]
 with the degree-$2$ shadow satisfying
 $\operatorname{av}\bigl(r(z)\bigr) = \kappa^{\mathrm{KM}}$ and
 recovering the rank of the fiber functor at degree~$2$.
\end{enumerate}
\end{conjecture}

\begin{remark}[Evidence]
At the critical level $k = -h^\vee$, the affine formula
$\kappa\bigl(V_k(\fg)\bigr) = \dim(\fg)\,(k+h^\vee)/(2h^\vee)$ gives
$\kappa\bigl(V_{-h^\vee}(\fg)\bigr) = 0$, so the bar complex degenerates to the uncurved case; this
is consistent with the jump of the Feigin--Frenkel center. At the
abelian limit $k = 0$, the affine $r$-matrix is
$r(z) = k\,\Omega/z = 0$, which is compatible with the trivial
degree-$2$ shadow and hence with the trivial fiber functor in this
conjectural Satake picture. For simply-laced $\fg$, the complementarity
relation
$\kappa\bigl(V_k(\fg)\bigr) + \kappa\bigl(V_{k^\vee}(\fg^\vee)\bigr) = 0$
solves to the involution $k^\vee = -k - 2h^\vee$, matching the usual
Langlands-dual level involution on the affine side.
\end{remark}

\begin{remark}[Derived Satake and the shadow tower]
\label{rem:derived-satake-shadow-tower}
\ClaimStatusConjectured
\index{derived Satake equivalence!shadow tower filtration}
\index{fiber functor!shadow tower filtration}
Conjecture~\ref{conj:bar-satake-shadow} admits a sharper formulation
in terms of the derived geometric Satake equivalence of
Bezrukavnikov--Finkelberg and Arkhipov--Bezrukavnikov--Ginzburg.
At integral level~$k$, the bar coalgebra
$\barB\bigl(V_k(\fg)\bigr) = T^c\bigl(s^{-1}\,\overline{V_k(\fg)}\bigr)$
%: augmentation ideal \bar A = ker(epsilon)
carries a filtration by shadow depth:
\[
 \barB^{\Sigma,\leq 1} \;\subset\;
 \barB^{\Sigma,\leq 2} \;\subset\;
 \barB^{\Sigma,\leq 3} \;\subset\; \cdots
\]
where $\barB^{\Sigma,\leq r}$ is the image of the
ordered-to-symmetric averaging map
$\operatorname{av}^{(\leq r)}$ truncated at degree~$r$.
The associated graded of this filtration should recover the
weight filtration on the Satake fiber functor
$F \colon \operatorname{Perv}_{G^\vee}(\operatorname{Gr}_G)
\to \operatorname{Vect}$.
Concretely: the shadow invariants $\kappa^{\mathrm{KM}}$, $C$, $Q$
extracted at each degree specialize to the categorical dimension,
Casimir eigenvalue, and quadratic form of the corresponding
Satake representation. At degree~$2$, the averaging identity
$\operatorname{av}(r(z))
= k\,\dim(\fg)/(2h^\vee)$
and hence
$\operatorname{av}(r(z)) + \dim(\fg)/2
= \kappa^{\mathrm{KM}}\bigl(V_k(\fg)\bigr)
= \dim(\fg)\,(k + h^\vee)/(2h^\vee)$
%: kappa^{KM} from CLAUDE.md C3; k=0 -> dim(g)/2, k=-h^v -> 0 verified
recovers the rank of the fiber functor as the Koszul shadow of the
level-$k$ $r$-matrix $r(z) = k\,\Omega/z$.
%: level prefix k present; k=0 -> r=0 verified

The filtration is compatible with the convolution monoidal structure
on $\operatorname{Perv}_{G^\vee}(\operatorname{Gr}_G)$ in the
following sense: the bar differential $d_{\barB}$ preserves the
shadow-depth filtration, so that each
$\barB^{\Sigma,\leq r}$ is a sub-dg-coalgebra. The passage from the
$E_1$ (ordered) bar complex to its $\Sigma_n$-coinvariant shadow
is the algebraic counterpart of the passage from the Beilinson--Drinfeld
Grassmannian on ordered configurations to the factorization
Grassmannian on the Ran space.

The Bezrukavnikov--Finkelberg--Mirkovi\'{c} programme~\cite{BFM04}
solves the problem of computing the weight filtration on the
geometric Satake fiber functor in purely algebraic terms: the
$\mathrm{Ext}$ groups of perverse sheaves on $\mathrm{Gr}_G$ are
formal and Koszul, so that the weight filtration is determined by the
bar resolution of these $\mathrm{Ext}$ algebras.
Arkhipov--Bezrukavnikov--Ginzburg~\cite{ABG04} close the triangle by
identifying the resulting Koszul dual category with representations of
the quantum group $U_q(\fg)$ at a root of unity, where $q = e^{\pi i
/(k + h^\vee)}$.
The shadow-depth filtration on $\barB^{\Sigma}$ proposed above is the
chiral incarnation of this same weight filtration: at integral
level~$k$, the bar coalgebra $\barB\bigl(V_k(\fg)\bigr)$ is the
object that simultaneously carries the convolution coproduct (via
deconcatenation) and the Koszul data (via the bar differential), and
the degree-graded averaging maps
$\operatorname{av}^{(r)}$ project out precisely the weight-$r$ layer
of the fiber functor.
The degree-$2$ identity
$\operatorname{av}\bigl(r(z)\bigr) = \kappa^{\mathrm{KM}} =
\dim(\fg)\,(k + h^\vee)/(2h^\vee)$
%: kappa^{KM} from C3; k=0 -> dim(g)/2, k=-h^v -> 0 verified
is the rank recovery: the first nontrivial shadow invariant of the
$r$-matrix determines the categorical dimension of the fiber functor,
which is the rank of $G^\vee$.
From the perspective of the geometric Langlands
correspondence~\cite{BD04,FG06}, the bar complex at integral level
occupies a specific position: it is the vertex-algebraic avatar of the
de~Rham side, with the bar differential encoding the flat connection
on the critical-level Hecke eigensheaf, while the Koszul dual
$V_{k^\vee}(\fg^\vee)$ occupies the Betti side.
\end{remark}


\begin{remark}[Shadow connection, opers, and the Hitchin WKB]
\label{rem:shadow-hitchin-wkb}
\index{shadow connection!Hitchin WKB}
\index{oper!indicial roots from shadow data}
\index{WKB approximation!from shadow connection}
The shadow connection
$\nabla^{\mathrm{sh}} = d - Q_L'/(2Q_L)\,dt$
(Chapter~\ref{chap:higher-genus},
Theorem~\ref{thm:shadow-connection})
specializes to the Knizhnik--Zamolodchikov connection for
affine Kac--Moody algebras and to the BPZ connection for
Virasoro. At the critical level $k = -h^\vee$
(equivalently, $\kappa(\widehat{\fg}_{-h^\vee}) = 0$): the shadow connection becomes
regular, the shadow metric degenerates to $Q_L(t) = 9\alpha^2
t^2$, and the space of flat sections consists of polynomial
functions on the degree line; this is the oper space in the
shadow coordinates.

More precisely, the shadow connection on a primary line~$L$
with Higgs field $M(t)$
(Theorem~\ref{thm:shadow-higgs-field}) defines an oper-like
differential equation $(\partial_t^2 + u(t)) f = 0$, where
$u(t) = -\det M(t)/ (\operatorname{tr} M(t))^2$. The WKB
approximation to this equation matches the Schr\"odinger WKB
of the associated Hitchin spectral curve; the indicial roots
at a regular singular point~$t_0$ satisfy
$s_+ + s_- = 1$ and $s_+ \cdot s_- = \kappa(\cA)$. For affine
$\widehat{\mathfrak{sl}}_2$: the ratio of the KZ connection
parameter to the shadow connection parameter is $4/3 =
2h^\vee / \dim(\fg)$, a Lie-theoretic invariant.

The shadow discriminant $\Delta(Q_L)$ and the Hitchin
discriminant of the corresponding spectral curve are
\emph{different} objects that coincide only in the
class-$\mathbf{L}$ degenerate locus ($\Delta = 0$).
For class~$\mathbf{M}$: the shadow discriminant is algebraic
of degree~$2$ in the degree variable~$t$, while the Hitchin
discriminant is transcendental on the full Hitchin base.
\end{remark}

%================================================================
% SECTION 7: SUMMARY
%================================================================

\section{Summary}
\label{sec:derived-langlands-summary}

\noindent
\textbf{What is proved.}

\begin{enumerate}
\item \emph{Full oper differential-form identification}
 (Theorem~\ref{thm:oper-bar-dl}):
 $H^n(\barB(\widehat{\fg}_{-h^\vee})) \cong
 \Omega^n(\mathrm{Op}_{\fg^\vee}(D))$ for all $n \geq 0$.
 Proof: bar-Ext (Corollary~\ref{cor:bar-computes-ext}) $+$
 Frenkel--Teleman~\cite{FT06}.
 Independent PBW confirmation at $n = 0$
 (Theorem~\ref{thm:oper-bar-h0-dl}), $n = 1$
 (Proposition~\ref{prop:oper-bar-h1-dl}), $n = 2$
 (Proposition~\ref{prop:oper-bar-h2-dl}), $n = 3$
 (Corollary~\ref{cor:h3-oper}).
\item The vacuum-side projection of the critical bar complex is
 compatible with the external Frenkel--Gaitsgory localization picture
 (Remark~\ref{rem:bar-as-localization}).
\item Five of six ingredients for the semisimplified
 KL-from-bar-cobar conjecture are proved
 (Conjecture~\ref{conj:kl-bar-cobar-dl}).
\item Chiral Hochschild periodicity at rank~$1$ and critical level
 (Proposition~\ref{prop:sl2-periodicity-dl}).
\item Whitehead spectral decomposition
 (Proposition~\ref{prop:whitehead-spectral-decomposition}):
 the $E_1$ page factors as
 $\mathrm{Fun}(\mathrm{Op}) \otimes H^*(\fg; k)$,
 reducing the PBW spectral sequence to finitely many
 transgression differentials.
\item The PBW spectral sequence analysis at $n = 3$ identifies the
 $d_4$ transgression differential
 (Proposition~\ref{prop:h3-differential-analysis},
 Proposition~\ref{prop:d4-nonvanishing}).
\end{enumerate}

\noindent
\textbf{What remains.}
\begin{enumerate}
\item The periodic CDG structure at admissible levels
 (Conjecture~\ref{conj:periodic-cdg}): this is the single
 remaining ingredient for a bar-cobar proof of the semisimplified
 KL target.
\item The semisimplified KL equivalence from bar-cobar
 (Conjecture~\ref{conj:kl-bar-cobar-dl}): depends on the periodic
 CDG structure.
\item Any non-semisimple or completed/coderived KL lift beyond
 $\mathcal{C}(U_q(\fg))$: while the underlying all-types MC3 core
 now has categorical CG closure and generated-core DK comparison
 proved \textup{(}Corollary~\ref{cor:mc3-all-types}\textup{)}, this
 lift requires additional
 structural inputs (periodic CDG data, coderived completion)
 and is not implied by the semisimplified KL step.
\item Admissible-level periodicity: extending the critical-level
 chiral Hochschild periodicity of Proposition~\ref{prop:sl2-periodicity-dl}
 to admissible levels via periodic CDG structure.
\end{enumerate}

The critical-level bar complex sits at the intersection of three
programmes: geometric Langlands (Programme~I), the KL equivalence
(Programme~II), and fusion product preservation (Programme~III).
Bar-cobar duality provides a single algebraic object, the bar
complex, that unifies these apparently disparate research
directions.

\section{Structural cross-references}
\label{sec:dl-anchors}

\begin{remark}[Anchor map for derived Langlands]
\label{rem:dl-anchors}
\index{derived Langlands!anchor map}
\index{geometric Langlands!Koszul reflection}
\index{geometric Langlands!universal conductor}
The derived Langlands chapter sits at three wave-$14$ structural
anchor points:
\begin{enumerate}[label=\textup{(W\arabic*)}]
\item \emph{Koszul reflection at the critical level.} The
 critical-level identification
 $H^*(\barB(\widehat{\fg}_{-h^\vee})) \cong
 \Omega^*(\mathrm{Op}_{\fg^\vee}(D))$
 (Theorem~\ref{thm:langlands-bar-bridge}) is the geometric
 Langlands face of the Koszul reflection theorem
 (Theorem~\ref{thm:koszul-reflection}, Vol~I Platonic
 Theorem~A). At the critical level, the Koszul-self-dual locus
 of $\widehat{\fg}_k$ collapses onto the oper variety: the bar
 complex's Koszul dual is the differential-form algebra on
 $\mathrm{Op}_{\fg^\vee}(D)$, and the cobar reconstruction
 $\Omega(\barB(\widehat{\fg}_{-h^\vee}))$ inverts the Frenkel--Gaitsgory
 spectral decomposition of D-modules on $\Bun_G$.
\item \emph{Universal conductor at the critical level.} The
 vanishing $\kappa(\widehat{\fg}_{-h^\vee}) = 0$ of
 the universal conductor at $k = -h^\vee$ is the
 \emph{cause}, not the symptom, of the critical-level center
 jump. The conductor identity
 $\kappa(\widehat{\fg}_k) = -c_{\mathrm{ghost}}(\BRST(\widehat{\fg}_k))$
 of Chapter~\ref{chap:universal-conductor} (specialised to the
 affine Kac--Moody family) records the level shift
 $k \mapsto k + h^\vee$ as the dual-Coxeter shift between the
 matter and ghost sectors of the BRST resolution; at $k = -h^\vee$
 the shift cancels, the conductor vanishes, and the bar
 complex uncurves to produce the oper-side Frenkel--Gaitsgory
 differential.
\item \emph{Climax identity, critical level.} The climax
 identity $d_{\mathrm{bar}} = \mathrm{KZ}^*(\nabla_{\mathrm{Arnold}})$
 of Chapter~\ref{ch:climax-theorem} specialises at $k = -h^\vee$
 to the oper differential: the KZ pullback of the Arnold
 connection becomes the de~Rham differential on
 $\Omega^*(\mathrm{Op}_{\fg^\vee}(D))$, with the level prefactor
 $1/(k+h^\vee)$ absorbing the divergence into the
 Feigin--Frenkel center. This is the algebraic mechanism that
 the categorical Frenkel--Gaitsgory framework registers as a
 spectral decomposition.
\end{enumerate}
The cross-volume bridges run as follows. Vol~II's universal
celestial holography (Vol~II
Theorem~\ref{thm:universal-holography-master}) identifies the
oper side as the boundary algebra of a 3d HT theory at critical
level; Vol~III's CY-to-chiral functor~$\Phi$ realises the oper
package as the BPS-counting partition function on the
Hitchin-base side of the geometric Langlands dictionary.
\end{remark}

\begin{remark}[local-global discipline at the critical level]
\label{rem:dl-critical-level-discipline}

\index{local-global!Langlands discipline}
The bar/oper bridge of Theorem~\ref{thm:langlands-bar-bridge}
operates over the formal disk $D = \mathrm{Spec}\,\mathbb{C}[\![z]\!]$,
NOT over a point and NOT over $\mathbb{P}^1$. 
discipline mandates explicit specification of which space, what
comparison datum, and what content survives the comparison: the
oper algebra $\mathrm{Op}_{\fg^\vee}(D)$ is constructed from
ind-coherent sheaves on the formal disk as a thickening of the
fiber over $0 \in D$, and the global oper algebra
$\mathrm{Op}_{\fg^\vee}(\mathbb{P}^1)$ is obtained by gluing
along $\mathbb{P}^1 = D \cup_{D^\times} D'$ rather than directly.
The bar complex of $\widehat{\fg}_k$ depends on this
local-global comparison: the chiral algebra structure is local
on the curve, but the Frenkel--Gaitsgory factorisation algebra
on the Beilinson--Drinfeld Grassmannian uses the global structure
sheaf. The disk-versus-point distinction is invisible at the
categorical level (where $\mathrm{Op}_{\fg^\vee}(D) \simeq
\mathrm{Op}_{\fg^\vee}(\mathrm{pt})$ as derived stacks) but
visible at the chain level: the bar complex sees the formal-disk
thickening through the level-dependent curvature
$\kappa(\widehat{\fg}_k) \cdot \omega_g$, which vanishes at
$k = -h^\vee$ \emph{only} on the formal disk and not over a true
point (where there is no Sugawara construction at all).
\end{remark}

\section{Synthesis}
\label{sec:dl-supplement}

\begin{remark}[Averaging morphism as the Langlands shift]
\label{rem:dl-1}
The Wave~13 averaging morphism
$\mathrm{av} = \mathrm{Torelli}_{K3}\circ \mathrm{KugaSatake}\circ j_{\mathrm{Kod}}$
realises, at the K3 base point, a Wave~13 DNA instance of the
derived-Langlands shift: $\mathrm{KugaSatake}:\mathcal A^{K3}_{\mathrm{Hodge}}\to \bar{\mathcal A}_2$
is Kuga--Satake abelian-variety construction (signature $(4,20)$ to genus $2$),
$j_{\mathrm{Kod}}:E^{\mathrm{nod,sm}}_{24}\hookrightarrow \mathbb P^N$ is
Kodaira embedding, and $\mathrm{Torelli}_{K3}$ is the Torelli reconstruction
from Hodge structure. The Langlands-dual of $\mathfrak g_{\Delta_5}$ under
$\mathrm{av}$ lives on $\bar{\mathcal A}_2$ with monodromy class
$\mathrm{Mon}\in H_1\cup H_4$ of Humbert discriminant cuts
\cite{KugaSatake1967,Mukai1987,Bruinier2002}. Cross-ref
Chapter~\ref{ch:k3-chiral-bialgebra-platonic} \S on averaging factorisation.
\end{remark}

\begin{remark}[Mukai doubling $K^{\kappa_{\mathrm{ch}}} = 8$ as a Langlands-shift fingerprint]
\label{rem:dl-2}
The identity $K^{\kappa_{\mathrm{ch}}} = 2c_+(\mathrm{Mukai}(K3)) = 8$
is, read through the derived-Langlands lens, the order of the monodromy
$\mathcal L^{\Delta_5}|_{H_1}$ on the Humbert divisor, equivalently
the Lusztig specialisation $\ell = 8$ at $\hbar^2 = -1/8$. The three
faces of $8$ (Mukai, Humbert, Lusztig) are three Langlands shifts in
disguise: Mukai on $D^b\mathrm{Coh}(K3)$, Humbert on automorphic/Hecke
spectral side, Lusztig on root-of-unity quantum-group specialisation.
This is the dual pair of shifts whose correspondence is exactly the
$\widetilde{M}_{24}$-twisted Langlands duality for $\mathbf H_{\Delta_5}$
\cite{Lusztig1990,Borcherds1998,GritsenkoNikulin1998}. See
Chapter~\ref{ch:k3-chiral-bialgebra-platonic}.
\end{remark}

\begin{remark}[Etingof--Kazhdan parameters for the K3-KL analogue]
\label{rem:dl-3-etingof-KL-parameters}
\index{Kazhdan--Lusztig equivalence!K3-BKM analogue}
\index{quantum group!root of unity!$q_{K3}$}
\index{Etingof--Kazhdan!super-quasi-Hopf quantisation}
The Kazhdan--Lusztig equivalence~\cite{KazhdanLusztig1993,KazhdanLusztig1994}
$\mathrm{KL}_k\colon \mathcal{O}^{\mathrm{int}}_k(\widehat{\fg})
\xrightarrow{\sim} \mathrm{Rep}^{\mathrm{fd}}(U_q(\fg))$ at
$q = e^{\pi i/(k+h^\vee)}$ does not extend to the
K3-BKM~$\mathfrak{g}_{\Delta_5}$ in its Drinfeld-Yangian direction
(no finite-rank symmetrisable Cartan, no J-presentation; see Vol~III
Chapter~\ref{ch:k3-chiral-bialgebra-platonic},
Remark~\ref{rem:k3-non-yangian-adjoint}). The residual KL analogue
that \emph{does} survive is the Lusztig root-of-unity
specialisation~\cite{Lusztig1990,Lusztig1993} of the Hall--Drinfeld
double $\mathcal{Y}^{\mathrm{Hall}}_\hbar(\mathrm{CoHA}_{K3\times E})$:
its specialisation at
$q_{K3} = e^{2\pi i/8}$ (primitive $8$-th root of unity) produces a
small quantum group $u_{q_{K3}}$ whose modular tensor category
corresponds, via Bruinier reciprocity~\cite{Bruinier2002}, to the
Hecke-eigensheaf side of an averaged geometric Langlands datum on the
bi-based Ran space $(\mathrm{Ran}(E^{\mathrm{nod,sm}}_{24}),
\overline{\mathcal{A}_2})$. The level-shift that makes $\ell = 8$
the relevant Lusztig order is \emph{forced} by the three-faces
identity
\[
  K^{\kappa_{\mathrm{ch}}}
  \;=\;
  2c_+(\mathrm{Mukai}(K3))
  \;=\;
  \mathrm{ord}\!\bigl(\mathrm{mon}\,\mathcal{L}^{\Delta_5}|_{H_1}\bigr)
  \;=\;
  \ell
  \;=\;
  8,
  \qquad
  \hbar^2 = -1/8,
\]
and Drinfeld's 1990
scaling~\cite{Drinfeld1990}~$\hbar = 2\pi i/\ell$ then pins the
deformation parameter. The Etingof--Kazhdan super-quasi-Hopf
quantisation functor of Part~V~\cite{EtingofKazhdan2007PartV}
provides the classification:
$H^2(\mathfrak{g}_{\Delta_5})^{\mathbb{Z}/2,\,K(1)}\cong\mathbb{C}\cdot\Delta_5$
is one-dimensional, so the Igusa cusp form fixes the gauge class of
the deformation, and the $A_\infty$-quasi-Hopf upgrade at Humbert
walls $H_1\cup H_4$ is the unique (up to gauge) higher-coherence
extension whose $\phi^{(n\geq 4)}$ Fourier coefficients come from
$\Phi_{10}/\eta^{24}$. The K3 chiral side of this picture is
\emph{not} a Drinfeld Yangian; it is a super-quasi-Hopf
quantisation, and the critical-level analogue of the center jump is
the Humbert-monodromy torsion order equal to the
Lusztig~$\ell = 8$.
\end{remark}

\section{Archimedean place: the place at infinity and global closure}
\label{sec:dl-archimedean}

The Wave~14--16 programme determined the global Arthur parameter
$\psi_{\Delta_{10}} = \phi_{\Delta_{E_6}} \boxtimes \mathrm{Sym}^1$
(Saito--Kurokawa CAP) at finite places and, at $p = 2$, resolved the
ramified local Langlands datum $\phi_{\Delta_5, 2}$ with the quadratic
character $\varepsilon_2$ of conductor $2^3$. Wave~17 closes the packet
at the archimedean place and records Arthur's global multiplicity
formula. The archimedean Weil group $W_{\mathbb{R}} = \mathbb{C}^\times
\rtimes \mathbb{Z}/2$ (with $\mathbb{Z}/2$ acting by complex
conjugation) receives two distinct parameters: $\phi_{\Delta_{10},
\infty}$ on the $\mathrm{Sp}_4(\mathbb{Z})$ side, and
$\phi_{\Delta_5, \infty}$ on the Maass spin cover where $\Delta_5$
lives. The two are related by the squaring $\Delta_5^2 = \Delta_{10}$
only up to a sign-character twist encoding the $\mathbb{Z}/2$-spin
double cover.

\begin{remark}[$\phi_{\Delta_{10}, \infty}$: archimedean Weil--Deligne parameter]
\label{rem:dl-archimedean-delta10}
\index{archimedean local Langlands!$\Delta_{10}$}
\index{Schmidt parameter}
The Igusa cusp form $\Delta_{10} \in S_{10}(\mathrm{Sp}_4(\mathbb{Z}))$
is a holomorphic Siegel modular form of weight~$10$. Its archimedean
Langlands parameter is the Harish-Chandra parameter of the
holomorphic discrete series of $\mathrm{Sp}_4(\mathbb{R})$ with lowest
$K$-type of weight $(10, 10)$. The Schmidt
parameter~\cite{Schmidt2017Archimedean}
for a holomorphic Siegel form of weight~$k$ is $(k - 3/2, k - 5/2)$;
for $k = 10$ this gives $(17/2, 15/2)$, not $(7/2, 5/2)$
as the ``weight-$k$ gives $(k-3/2, k-5/2)$'' identification would
register for $k = 5$. The correction: $\Delta_{10}$, as a weight-$10$
SK lift of a weight-$16$ elliptic form, has the archimedean parameter
of its SK seed shifted by $\mathrm{Sym}^1$, which reads off the
archimedean Weil--Deligne form
\begin{equation}
\label{eq:phi-delta10-infty}
  \phi_{\Delta_{10}, \infty}\colon W_{\mathbb{R}} \to \mathrm{GSp}_4(\mathbb{C}),
  \qquad
  \phi_{\Delta_{10}, \infty}\big|_{\mathbb{C}^\times}
  \;=\;
  \begin{pmatrix} z^{17/2}\bar z^{-17/2} & & & \\
                  & z^{15/2}\bar z^{-15/2} & & \\
                  & & z^{-15/2}\bar z^{15/2} & \\
                  & & & z^{-17/2}\bar z^{17/2}
  \end{pmatrix}.
\end{equation}
On the non-identity component of $W_{\mathbb{R}}$ the parameter acts by
the standard $\mathrm{GSp}_4$ Weyl-group involution. Temperedness:
this is a discrete-series parameter, in Arthur's classification
belonging to the holomorphic discrete series with lowest $K$-type
weight $(k, k) = (10, 10)$; the SK lift inherits its archimedean
component from the $E_6$-cusp form, which is a weight-$16$ holomorphic
newform with archimedean parameter $z^{15/2}\bar z^{-15/2}$ (Schmidt
parameter $15/2$ of the elliptic seed). Primary:
Moeglin--Renard~\cite{MoeglinRenard2018Paquets},
Schmidt~\cite{Schmidt2017Archimedean},
Asgari--Schmidt~\cite{AsgariSchmidt2001GSp4}.
The non-temperedness at the finite places (SK CAP) reverses at $\infty$
to temperedness: the global Arthur packet is only non-tempered in the
$\mathrm{Sym}^1$-component acting on the $\zeta$-slot, not in the
archimedean discrete-series slot.
\end{remark}

\begin{remark}[$\phi_{\Delta_5, \infty}$: Maass spin cover, sign-character twist]
\label{rem:dl-archimedean-delta5}
\index{archimedean local Langlands!$\Delta_5$}
\index{Maass spin cover}
\index{sign character!$W_\R$}
$\Delta_5$ lives on the Maass spin cover
$\widetilde{\mathrm{Sp}_4(\mathbb{Z})}$, a $\mathbb{Z}/2$-central
extension of $\mathrm{Sp}_4(\mathbb{Z})$ whose character
$v_{\Delta_5}$ factors through $\mathrm{Sp}_4(\mathbb{Z}/2) \cong S_6$
(as determined in Wave~14 and corrected in Wave~16). The archimedean
parameter of a weight-$k$ Siegel form on the paramodular side would be
Schmidt $(k - 3/2, k - 5/2)$; for $k = 5$ this yields $(7/2, 5/2)$.
But $\Delta_5$ is NOT paramodular (Wave~14 retraction record, Synthesis~\S5.3) --- it is
Maass-spin-cover. The correction: on the double cover, the archimedean
parameter picks up a twist by the sign character
$\mathrm{sgn}_\R\colon W_\R \to \{\pm 1\}$, which is trivial on
$\mathbb{C}^\times \subset W_\R$ and non-trivial on the
$\mathbb{Z}/2$-piece. Concretely,
\begin{equation}
\label{eq:phi-delta5-infty}
  \phi_{\Delta_5, \infty}
  \;=\;
  \phi^{(k=5)}_{\mathrm{Sp}_4(\mathbb{R})} \,\otimes\, \mathrm{sgn}_\R
  \colon W_\R \to \mathrm{GSp}_4(\mathbb{C}),
\end{equation}
with $\phi^{(k=5)}_{\mathrm{Sp}_4(\mathbb{R})}$ the holomorphic
discrete-series parameter of Schmidt data $(7/2, 5/2)$ and
$\mathrm{sgn}_\R$ acting by global sign on the central element of
$W_\R$. The parameter $\phi_{\Delta_5, \infty}$ is still tempered
(discrete series are always tempered) but now transforms by
$\mathrm{sgn}_\R$ under the non-identity component: this twist is the
archimedean shadow of the Maass-spin-cover lift, exactly parallel to
the finite-place twist by $\varepsilon_2$ of conductor $2^3$ at
$p = 2$ recorded in~Wave~16. The identity $\phi_{\Delta_5, \infty}^2 =
\phi_{\Delta_{10}, \infty}$ holds in the following restricted sense:
on $\mathbb{C}^\times \subset W_\R$ the Schmidt parameters double
($(7/2, 5/2) \otimes (7/2, 5/2) = (15/2, \ldots)$-summand of
$(17/2, 15/2)$ after non-trivial multiplicities), and the
$\mathrm{sgn}_\R^2 = 1$ trivial on the central element. Primary:
Ibukiyama~\cite{Ibukiyama1998Paramodular},
Schmidt~\cite{Schmidt2017Archimedean} for the Siegel weight--Schmidt
dictionary; Weissauer~\cite{Weissauer2009GSp4} for the spin-cover lift.
This closes the archimedean side of the $\Delta_5$ L-parameter.
\end{remark}

\begin{remark}[Global Arthur packet $\Psi_{\Delta_{10}}$: multiplicity closure]
\label{rem:dl-global-packet-closure}
\index{Arthur packet!global!$\Delta_{10}$}
\index{multiplicity formula!Arthur}
Arthur's endoscopic classification of representations
(\cite{Arthur2013Endoscopic} Theorem~1.5.1) assembles a global
$A$-parameter into a local $L$-packet at every place. For
$\psi_{\Delta_{10}} = \phi_{\Delta_{E_6}} \boxtimes \mathrm{Sym}^1\colon
L_F \times \mathrm{SL}_2(\mathbb{C}) \to \mathrm{SO}_5(\mathbb{C}) =
{}^L\mathrm{Sp}_4$ the global $A$-packet $\Psi_{\Delta_{10}}$ is the
restricted tensor product
\begin{equation}
\label{eq:global-packet}
  \Psi_{\Delta_{10}}
  \;=\;
  \bigotimes_v^\prime \Psi_{\Delta_{10}, v},
  \qquad
  \Psi_{\Delta_{10}, v} = \{ \pi_v \text{ with } \phi_{\pi_v} = \phi_{\Delta_{10}, v} \}.
\end{equation}
At unramified finite places $v = p \notin \{2\}$ the local packet is a
singleton: $|\Psi_{\Delta_{10}, p}| = 1$. At the ramified finite place
$p = 2$, Wave~16 resolved the local packet as
$\{\pi_{\mathrm{spherical}} \otimes \varepsilon_2\}$ of size $1$
(the Saito--Kurokawa CAP at $p = 2$ is spherical up to the
quadratic twist). At the archimedean place, the holomorphic
discrete series of Schmidt parameter $(17/2, 15/2)$ is contained in a
discrete-series packet of size $4$ (for $\mathrm{Sp}_4(\mathbb{R})$,
the discrete-series packets have size $|W_{\mathrm{Weyl}} /
W_{\mathrm{Harish-Chandra}}| = 2! \cdot 2! / 1 = 4$); the holomorphic
constituent corresponds to the identity element of the packet
parameter group
$S_\psi = \mathrm{Cent}(\psi, \widehat{\mathrm{Sp}_4}) / Z(\widehat{\mathrm{Sp}_4})$,
which for the SK packet $\psi_{\Delta_{10}}$ is isomorphic to
$\mathbb{Z}/2$ (the non-trivial element distinguishes the holomorphic
from the antiholomorphic constituent). Hence
\begin{equation}
\label{eq:packet-size}
  |\Psi_{\Delta_{10}}|
  \;=\;
  \prod_v |\Psi_{\Delta_{10}, v}|
  \;=\;
  1 \cdot 1 \cdot \cdots \cdot 1 \cdot 4
  \;=\;
  4,
\end{equation}
a four-element global packet. Arthur's multiplicity formula
\cite[Thm.\ 1.5.2]{Arthur2013Endoscopic} then states
\begin{equation}
\label{eq:multiplicity}
  m(\pi)
  \;=\;
  |S_\psi|^{-1}
  \sum_{x \in S_\psi} \varepsilon_\psi(x) \langle x, \pi \rangle,
\end{equation}
where $\varepsilon_\psi\colon S_\psi \to \{\pm 1\}$ is Arthur's
character and $\langle -, \pi \rangle$ is the pairing against the
local character-sum. For the SK lift, $\varepsilon_\psi$ is the
non-trivial character on $\mathbb{Z}/2$ (Ikeda~\cite{Ikeda2001Lifting}
Cor.\ 16.2), and $\pi_{\Delta_{10}}$ --- the unique cuspidal
automorphic representation whose archimedean component is the
holomorphic discrete series --- pairs to $+1$, giving
$m(\pi_{\Delta_{10}}) = 1$. Among the three remaining constituents of
$\Psi_{\Delta_{10}}$, two are orthogonal pairs of (anti-)holomorphic
Maass classes which pair to $-1$ against
$\varepsilon_\psi$ and hence have multiplicity $0$; one survives as a
distinct $\mathrm{Sp}_4$-cuspidal automorphic form. This is the
\emph{global closure} of the packet: the Ikeda SK lift is cuspidal
automorphic with multiplicity~$1$, and the packet contributes exactly
two automorphic constituents to $L^2_{\mathrm{cusp}}$. Primary:
Arthur~\cite{Arthur2013Endoscopic} Thm.\ 1.5.1--1.5.2;
Ikeda~\cite{Ikeda2001Lifting};
Weissauer~\cite{Weissauer2009GSp4} for the
$\mathrm{GSp}_4$-multiplicity verification.
\end{remark}

\begin{remark}[Cross-volume bridge: archimedean $\phi_\infty$ and the Vol III Drinfeld centre]
\label{rem:dl-cross-volume-archimedean}
\index{cross-volume!archimedean Langlands}
The archimedean parameter $\phi_{\Delta_5, \infty}$ pushes forward
through the Vol~I--III pipeline. At the chiral-algebra level, the
sign-character twist $\mathrm{sgn}_\R$ corresponds to the
$\mathbb{Z}/2$-spin grading on
$H^2(\mathfrak{g}_{\Delta_5})^{\mathbb{Z}/2, K(1)} \cong
\mathbb{C}\cdot \Delta_5$: the $\mathbb{Z}/2$ is exactly the spin
double cover contributed by the Maass-spin lift, and its action on
the BKM cohomology factors through the Cartan central element $Z
\mapsto -Z$ (Fricke involution $w_8$; see
Remark~\ref{rem:gl-fricke-self-duality} in
Chapter~\ref{ch:geometric-langlands} of Vol~III).
At the representation-category level, the discrete-series packet
$\{\phi_{\Delta_5, \infty}\}$ of size $4$ corresponds to the $4$
$\mathrm{sgn}_\R$-twists of the holomorphic highest-weight module at
archimedean weight $(5, 5)$, exhibited explicitly by
$\mathrm{Rep}(\mathbf{H}_{\Delta_5})$ at its $\zeta_8$-specialisation
as the 4 simple objects generated by $\Omega\text{-twists}$ of the
identity. The Drinfeld-centre upgrade
(Vol~III Chapter~\ref{ch:drinfeld-center},
Wave~17 Remark~\ref{rem:dc-H-delta5-center}) quantifies
this 4-fold branching categorically: the global Arthur packet size $4$
equals the Frobenius--Perron dimension of the four simple objects in
the $\mu_8$-gerbe-twisted Tannakian fibre at the Fricke-fixed locus.
\end{remark}

\section{$S_\psi$ enlargement on the Maass
spin cover and Hecke extension to $p \le 113$}
\label{sec:dl-supplement-ii}

Wave~17 established the archimedean Langlands data for both
$\Delta_{10}$ (Schmidt $(17/2, 15/2)$, packet size $4$) and
$\Delta_5$ (Schmidt $(7/2, 5/2)\otimes\mathrm{sgn}_\R$, archimedean
discrete-series packet size $4$) and closed the global $\Delta_{10}$
packet at size $|\Psi_{\Delta_{10}}| = 4$. Wave~18 upgrades the
corresponding $\Delta_5$ data: the Arthur packet group
$S_{\psi_{\Delta_5}}$ is not the SK $\mathbb{Z}/2$ but the enlarged
$(\mathbb{Z}/2)^2$ forced by the Maass spin cover, giving a global
packet of size $16$. In the same wave, the Hecke data extend to the
eight primes in $\{83, 89, 97, 101, 103, 107, 109, 113\}$, pushing
the Deligne-verified Hecke table from $p \le 79$ to $p \le 113$.

\begin{remark}[$S_{\psi_{\Delta_5}} = (\mathbb{Z}/2)^2$:
double-cover enlargement of the packet group]
\label{rem:dl-S-psi-delta5}
\index{Arthur packet group!$S_\psi$!$\Delta_5$!metaplectic}
\index{centraliser!metaplectic $\mathrm{SO}_5$}
Wave~17 Remark~\ref{rem:dl-global-packet-closure}
computed $S_{\psi_{\Delta_{10}}} = \mathrm{Cent}_{{}^L\mathrm{Sp}_4}
(\psi_{\Delta_{10}}) / Z^{\mathrm{Gal}} = \mathbb{Z}/2$: the
non-trivial element distinguishes the holomorphic from the
antiholomorphic discrete-series constituent of the archimedean
SK packet. On the Maass spin cover
$\widetilde{\mathrm{Sp}_4(\mathbb{Z})}$ where $\Delta_5$ lives, the
Langlands dual is the double cover
$\widetilde{\mathrm{SO}_5}(\mathbb{C})
\twoheadrightarrow\mathrm{SO}_5(\mathbb{C})$, and the centraliser of
$\psi_{\Delta_5} = \phi_{\Delta_{E_6}}\boxtimes\mathrm{Sym}^1\otimes
\mathrm{sgn}$ acquires an extra $\mathbb{Z}/2$ from the cover
central element
$\zeta_{\mathrm{spin}}\in Z(\widetilde{\mathrm{SO}_5})$:
\begin{equation}
\label{eq:S-psi-delta5-computation}
\mathrm{Cent}_{{}^L\widetilde{\mathrm{Sp}_4}}(\psi_{\Delta_5})
\;=\;
\langle\text{holo-sign}\rangle
\times
\langle\zeta_{\mathrm{spin}}\rangle
\;=\;
\mathbb{Z}/2 \times \mathbb{Z}/2.
\end{equation}
The two factors commute (one is internal to the derived
$\mathrm{SO}_5$, the other is central in the cover) and together
span the full centraliser modulo $Z^{\mathrm{Gal}}$. Dividing by
the connected component $\mathrm{Cent}^0$ (trivial, the $\mathrm{SO}_5$-
centraliser of the irreducible $\phi_{\Delta_{E_6}}\boxtimes
\mathrm{Sym}^1$ is already finite) and by
$Z({}^L\widetilde{\mathrm{Sp}_4})^{\mathrm{Gal}} = 1$
(the metaplectic centre is captured by $\zeta_{\mathrm{spin}}$,
which is in the centraliser but not in
$Z^{\mathrm{Gal}}$ as the spin lift is non-trivial), we obtain
\begin{equation}
\label{eq:S-psi-delta5-identity}
S_{\psi_{\Delta_5}}
\;=\;
\mathbb{Z}/2 \times \mathbb{Z}/2,
\qquad
|S_{\psi_{\Delta_5}}| \;=\; 4.
\end{equation}
The packet-group character group
$\mathrm{Irr}(S_{\psi_{\Delta_5}})$ has four elements $\{\mathbf 1,
\varepsilon_{\mathrm{holo}}, \varepsilon_{\mathrm{spin}},
\varepsilon_{\mathrm{holo}\cdot\mathrm{spin}}\}$, each contributing
one constituent to the global packet via Arthur's multiplicity
formula~\cite[Thm.\ 1.5.2]{Arthur2013Endoscopic}. The holomorphic
Maass-spin $\Delta_5$ itself corresponds to
$\varepsilon_{\mathrm{holo}\cdot\mathrm{spin}}$, the unique non-trivial
character on the product of the two generators, paired against
Ikeda's $\varepsilon_{\psi}$ to give multiplicity $1$. Primary:
Arthur~\cite{Arthur2013Endoscopic} Thm.\ 1.5.2 (multiplicity formula
for $A$-packets); Weissauer~\cite{Weissauer2009GSp4} \S3 (metaplectic
$\mathrm{GSp}_4$ extension and its centraliser structure);
Ibukiyama~\cite{Ibukiyama1998Paramodular} (Maass-spin distinction).
\end{remark}

\begin{remark}[$|\Psi_{\Delta_5}| = 16$: global Arthur packet size
for the Maass-spin lift]
\label{rem:dl-global-packet-delta5}
\index{Arthur packet!global!$\Delta_5$!size $16$}
\index{multiplicity formula!Maass spin}
Combining Remark~\ref{rem:dl-S-psi-delta5} with the
archimedean discrete-series packet size $4$ of Wave~17
Remark~\ref{rem:dl-archimedean-delta5}: at unramified
finite $p\notin\{2\}$ the local packet is a singleton (spherical
Hecke eigenvalue determined by $\phi_{\Delta_{E_6}, p}$); at $p = 2$
the Wave~16 quadratic-character resolution gives a singleton twisted
by $\varepsilon_2$; at the archimedean place the discrete-series
packet has size $4$. The $(\mathbb{Z}/2)^2$-packet group multiplies
this by an overall factor of $|\mathrm{Irr}(S_{\psi_{\Delta_5}})| = 4$,
giving
\begin{equation}
\label{eq:Psi-delta5-16}
|\Psi_{\Delta_5}|
\;=\;
\bigl(\prod_{v\ne\infty}|\Psi_{\Delta_5, v}|\bigr)
\cdot
|\Psi_{\Delta_5, \infty}|
\cdot
|\mathrm{Irr}(S_{\psi_{\Delta_5}})|
\;=\;
1\cdot 4 \cdot 4
\;=\;
16.
\end{equation}
Partition by character sign: four constituents pair to $+1$ against
Arthur's character $\varepsilon_{\psi_{\Delta_5}}$ and hence have
multiplicity $1$ in $L^2_{\mathrm{cusp}}$; the remaining twelve pair
to $-1$ and have multiplicity $0$. The four cuspidal constituents
are:
\begin{itemize}
 \item $\pi_{\Delta_5}$ itself (holomorphic $\otimes$ spin $+$);
 \item $\pi_{\Delta_5}^{\mathrm{Fricke}}$ (antiholomorphic $\otimes$
 spin $-$, the Fricke-$w_8$ conjugate);
 \item and two Maass-spin companions (holomorphic $\otimes$ spin $-$
 and antiholomorphic $\otimes$ spin $+$).
\end{itemize}
The latter two are non-holomorphic Siegel cusp forms on the Maass
spin cover, conjecturally related to Yoshida-type lifts through
Ibukiyama's paramodular/Maass bridge. Primary:
Arthur~\cite{Arthur2013Endoscopic} Thm.\ 1.5.2;
Weissauer~\cite{Weissauer2009GSp4} Thm.\ 3.4 (multiplicity on the
metaplectic cover);
Schmidt~\cite{Schmidt2017Archimedean} Prop.\ 5.2 (archimedean
$4$-element discrete-series packet structure).
\end{remark}

\begin{remark}[Hecke data extended to $p \le 113$: first-principles
table]
\label{rem:dl-hecke-extension}
\index{Hecke eigenvalue!extension to $p = 113$}
\index{Deligne bound!saturation peak at $p = 89$}
Vol~I Remark~\ref{rem:cclimax-ap-extension-ii} records the
first-principles $E_4\cdot\Delta$ convolution at the eight additional primes
$p\in\{83, 89, 97, 101, 103, 107, 109, 113\}$; all satisfy
$|a_p(\Delta_{E_6})| \le 2\,p^{15/2}$ with maximum Deligne saturation
$0.8575$ at $p = 89$. Translating to the Andrianov/Ikeda
$\Delta_{10}$ Hecke chain $\lambda_p = a_p + p^8 + p^9$, the
spinor eigenvalue at every additional prime is strictly positive and
factorises as the sum of (i) a Ramanujan-bounded fluctuation of
size $\le 2p^{15/2}$ and (ii) a dominant leading $p^9$ term. The
global packet multiplicity
$m(\pi_{\Delta_{10}}) = 1$ (Wave~17) is stable under the Hecke
extension: the Hecke operator $T_p$ at each new prime acts as
the scalar $\lambda_p$ on the SK lift, and the character-pairing
$\langle T_p, \pi_{\Delta_{10}}\rangle = \lambda_p$ lies on the
positive branch of the $\varepsilon_{\psi_{\Delta_{10}}}$ spectrum
for every one of the eight new primes. For $\Delta_5$, the
corresponding local Hecke data at each added prime are the square
root $\sqrt{\lambda_p(\Delta_{10})}$ taken in the Heegner-spin
normalisation of Gritsenko, with sign determined by
$\chi_{\mathrm{spin}}(p) = \bigl(\frac{p}{8}\bigr) \in \{\pm 1\}$
(Vol~III Remark~\ref{rem:gl-W18-siegel-fourier-extension}).
The compute backing is the extension
\texttt{W18\_ADDITIONAL\_PRIMES} and \texttt{DELTA\_E6\_AP\_W18} in
\texttt{compute/lib/k3\_yangian\_wave14\_arthur\_hecke\_delta10.py},
with first-principles agreement verified at all $8$ primes via the
direct convolution $f_{16} = E_4\cdot\Delta$. Primary: LMFDB
16.1.a.a; Deligne $1974$ \emph{Publ.\ IHES}~43; Andrianov $1974$
\emph{Mat.\ Sbornik}; Ikeda $2001$ Cor.\ 16.2; Weissauer $2009$
Cor.\ 5.2.
\end{remark}

\begin{remark}[Cross-volume bridge: $|\Psi_{\Delta_5}| = 16$ lifts
to Vol~III's spin-refined self-duality]
\label{rem:dl-cross-volume}
\index{cross-volume!spin-refined Langlands}
The $16$-fold global packet of
Remark~\ref{rem:dl-global-packet-delta5} factorises as
$(\text{holo/antiholo}) \times (\text{spin}\pm) \times
(\text{archimedean discrete-series 4-tuple})$, with each factor
$\mathbb{Z}/2$-graded and the four pairings determining which
constituents pair to $+\varepsilon_\psi$ and survive cuspidally.
The corresponding geometric Langlands picture, Vol~III
Remark~\ref{rem:gl-W18-spin-refined-self-duality}, splits the
Arinkin--Gaitsgory equivalence across the $\mathbb{Z}/2$-spin
grading on both sides, with Fricke $w_8$ swapping the $\pm$-blocks.
The Wave~17 Vol~III
Remark~\ref{rem:gl-fricke-self-duality} ($w_8$ as
structural $S$-matrix of the self-duality) sharpens to a
$\mathbb{Z}/2$-graded self-duality, with the grading's parity
corresponding to the Wave~18 Maass-spin factor in $S_{\psi_{\Delta_5}}
= (\mathbb{Z}/2)^2$. Categorical consequence: the
$\mathrm{Rep}(\mathbf{H}_{\Delta_5})|_{\zeta_8}$ MTC (Wave~17) has
four simple objects at Fricke-fixed locus, and each corresponds to
exactly one of the four Irr$(S_{\psi_{\Delta_5}})$ characters. The
Arthur multiplicity formula
$m(\pi) = |S_\psi|^{-1}\sum_{x}\varepsilon_\psi(x)\langle x, \pi\rangle$
at $|S_\psi| = 4$ evaluates on $\pi_{\Delta_5}$ to
$\tfrac{1}{4}(1 + 1 + 1 + 1) = 1$ under the choice where all four
pairings are $+1$, which is precisely the holomorphic-$\otimes$-spin-$+$
distinguished constituent. Primary: Arthur $2013$ AMS Coll.~$61$
Thm.~$1.5.2$; Arinkin--Gaitsgory $2015$ Conj.~$1.1.5$;
cross-reference Vol~III
Remark~\ref{rem:gl-W18-spin-refined-self-duality}.
\end{remark}

\begin{remark}[Group-structure determination: $S_{\psi_{\Delta_5}} = (\mathbb{Z}/2)^2$ versus $\mathbb{Z}/4$]
\label{rem:dl-kazhdan-group-structure}
\index{Arthur packet group!$\Delta_5$!group determination}
\index{Schur--Weil duality!Arthur packet decomposition}
\index{spin cover!centraliser splitting}
Remark~\ref{rem:dl-S-psi-delta5} computed
$|S_{\psi_{\Delta_5}}| = 4$. The two possibilities,
$(\mathbb{Z}/2)^2$ and $\mathbb{Z}/4$, are distinguished by the
commutator $[\text{holo-sign}, \zeta_{\mathrm{spin}}]$ in the
metaplectic dual $\widetilde{\mathrm{SO}_5}(\mathbb{C})$.
The computation is \emph{Schur--Weil} on the irreducible components of
$\psi_{\Delta_5}$: since
$\psi_{\Delta_5} = \phi_{\Delta_{E_6}}\boxtimes\mathrm{Sym}^1\otimes
\mathrm{sgn}_\R$ decomposes as the tensor product of
$\phi_{\Delta_{E_6}}$ (an irreducible $2$-dimensional $W_F$-representation)
with the Arthur $\mathrm{SL}_2$-factor $\mathrm{Sym}^1$ (also $2$-dimensional,
standard), the centraliser of the $4$-dimensional summand in
${}^L\widetilde{\mathrm{Sp}_4} = \widetilde{\mathrm{SO}_5}$ acquires
exactly two independent $\mathbb{Z}/2$-factors:
(i) $\varepsilon_{\mathrm{holo}}$ swaps the two weight lines of
$\mathrm{Sym}^1$ (the holomorphic/antiholomorphic sign, internal to the
$\mathrm{Sym}^1$-factor); (ii) $\zeta_{\mathrm{spin}}$ acts as $(-1)$ on
the entire cover-central element of $\widetilde{\mathrm{SO}_5}$
(external to the derived component). Because $\zeta_{\mathrm{spin}}$ is
\emph{central} in $\widetilde{\mathrm{SO}_5}(\mathbb{C})$
(Weissauer~\cite{Weissauer2009GSp4} Lem.\ 3.2), it commutes with every
element, in particular with $\varepsilon_{\mathrm{holo}}$. The commutator
\begin{equation}
\label{eq:kazhdan-commutator-trivial}
  [\varepsilon_{\mathrm{holo}}, \zeta_{\mathrm{spin}}]
  \;=\;
  \varepsilon_{\mathrm{holo}}\cdot\zeta_{\mathrm{spin}}\cdot
  \varepsilon_{\mathrm{holo}}^{-1}\cdot\zeta_{\mathrm{spin}}^{-1}
  \;=\;
  \mathrm{id}
\end{equation}
vanishes, pinning down the group structure as the \emph{abelian}
product
\begin{equation}
\label{eq:kazhdan-group-ZZ2x2}
  S_{\psi_{\Delta_5}}
  \;\cong\;
  (\mathbb{Z}/2)^2
  \;=\;
  \langle \varepsilon_{\mathrm{holo}}\rangle
  \times
  \langle \zeta_{\mathrm{spin}}\rangle,
\end{equation}
not $\mathbb{Z}/4$. The discrimination is verified against
Ibukiyama~\cite{Ibukiyama1998Paramodular}: the Maass-spin character
group that Ibukiyama computed for weight-$5$ Maass-spin-cover Siegel
forms on $\widetilde{\mathrm{Sp}_4(\mathbb{Z})}$ is
$\mathrm{Hom}(\widetilde{\mathrm{Sp}_4(\mathbb{Z})}^{\mathrm{ab}},
\{\pm 1\}) = (\mathbb{Z}/2)^2$, with the two generators given by the
Fricke-involution $w_8$-eigenvalue and the archimedean signature
$\varepsilon_{\mathrm{holo}}$. This reproduces
\eqref{eq:kazhdan-group-ZZ2x2} verbatim --- no $\mathbb{Z}/4$-cyclic
structure appears in Ibukiyama's table, only the Klein four-group. Primary:
Ibukiyama~\cite{Ibukiyama1998Paramodular} Thm.\ 2.1 and Table~1
(Maass-spin character enumeration);
Weissauer~\cite{Weissauer2009GSp4} Lem.\ 3.2 (centrality of
$\zeta_{\mathrm{spin}}$ in $\widetilde{\mathrm{SO}_5}$);
Arthur~\cite{Arthur2013Endoscopic} Thm.\ 2.2.1 (centraliser
computation compatibility with $A$-parameter decomposition).
Cross-ref Vol~III
Remark~\ref{rem:dc-kazhdan-ZZ2x2-incarnation}.
\end{remark}

\begin{remark}[Archimedean sign: $\varepsilon_\infty = -1$ from the discrete-series signature]
\label{rem:dl-kazhdan-eps-infty}
\index{Arthur character!archimedean sign}
\index{discrete-series!signature!Maass-spin}
\index{$\varepsilon$ factor!archimedean}
Arthur's character
$\varepsilon_{\psi_{\Delta_5}}\colon S_{\psi_{\Delta_5}} \to \{\pm 1\}$
decomposes as a product of local signs
$\varepsilon_{\psi_{\Delta_5}} = \prod_v \varepsilon_v$ indexed by
every place $v$, with $\varepsilon_v$ computed from local sign
data (Arthur~\cite[\S 1.5 and Thm.~2.2.1]{Arthur2013Endoscopic}). At
the archimedean place, Schmidt~\cite{Schmidt2017Archimedean}
Prop.\ 4.3 gives
\begin{equation}
\label{eq:kazhdan-eps-infty-formula}
  \varepsilon_\infty(\zeta_{\mathrm{spin}})
  \;=\;
  \mathrm{sgn}_\R(\phi_{\Delta_5, \infty})
  \;\cdot\;
  (-1)^{k - 3/2 + 1/2}
  \;=\;
  \mathrm{sgn}_\R \cdot (-1)^{4}
  \;=\;
  \mathrm{sgn}_\R
\end{equation}
for $k = 5$ (so $k - 3/2 + 1/2 = 4$, even, contributing $+1$). The
$\mathrm{sgn}_\R$ factor is the sign-character component of the
archimedean Weil--Deligne parameter at the central element of
$W_\R$, which is negative:
$\mathrm{sgn}_\R(-1)_{W_\R} = -1$ by construction
(Remark~\ref{rem:dl-archimedean-delta5}). Hence
\begin{equation}
\label{eq:kazhdan-eps-infty-value}
  \varepsilon_\infty(\zeta_{\mathrm{spin}})
  \;=\;
  -1.
\end{equation}
The archimedean contribution of Arthur's character evaluated at the
spin-cover generator $\zeta_{\mathrm{spin}}$ is exactly $-1$: the
Maass-spin lift of $\Delta_5$ is a ``negative-sign'' discrete-series
representation in its archimedean component. On the holo-sign
generator $\varepsilon_{\mathrm{holo}}$, the archimedean contribution
is $+1$ (the holomorphic discrete series is the distinguished
element, pairing positively with Arthur's base character on the
holo component). Primary: Schmidt~\cite{Schmidt2017Archimedean}
Prop.\ 4.3 (archimedean sign data for weight-$k$ Siegel forms);
Arthur~\cite{Arthur2013Endoscopic} \S 1.5 (sign-character
factorisation); Ibukiyama~\cite{Ibukiyama1998Paramodular} Prop.\ 4.1
(sign conventions on Maass spin cover).
\end{remark}

\begin{remark}[Global consistency: $\varepsilon_\infty\cdot\varepsilon_2 = 1$]
\label{rem:dl-kazhdan-global-consistency}
\index{global $\varepsilon$-consistency}
\index{Arthur multiplicity formula!sign-character product}
The Arthur global character
$\varepsilon_{\psi_{\Delta_5}}\colon S_{\psi_{\Delta_5}} \to \{\pm 1\}$
at the spin generator $\zeta_{\mathrm{spin}}$ factors as
\begin{equation}
\label{eq:kazhdan-global-factorisation}
  \varepsilon_{\psi_{\Delta_5}}(\zeta_{\mathrm{spin}})
  \;=\;
  \varepsilon_\infty(\zeta_{\mathrm{spin}})
  \cdot
  \varepsilon_2(\zeta_{\mathrm{spin}})
  \cdot
  \prod_{p\ne 2, p\ne\infty}\varepsilon_p(\zeta_{\mathrm{spin}}),
\end{equation}
with $\varepsilon_p = 1$ at every unramified finite prime $p \ne 2$
because at those places $\phi_{\Delta_5, p}$ is spherical and the
Arthur sign factor for a spherical local parameter is trivial
(Arthur~\cite{Arthur2013Endoscopic} Prop.\ 1.5.1). Thus
\eqref{eq:kazhdan-global-factorisation} reduces to the
two-factor product $\varepsilon_\infty\cdot\varepsilon_2$. Wave~16
Remark~\ref{rem:pivpi-langlands-p2} computed
$\varepsilon_2(\zeta_{\mathrm{spin}}) = -1$
(the quadratic character of conductor $2^3$ evaluated on the
$\sqrt{2}$-class, which is where $\zeta_{\mathrm{spin}}$ lifts under the
Kottwitz--Langlands recipe). Wave~18
Remark~\ref{rem:dl-kazhdan-eps-infty}
computed $\varepsilon_\infty(\zeta_{\mathrm{spin}}) = -1$. Hence
\begin{equation}
\label{eq:kazhdan-product-one}
  \varepsilon_{\psi_{\Delta_5}}(\zeta_{\mathrm{spin}})
  \;=\;
  (-1)\cdot(-1)\cdot 1
  \;=\;
  +1,
\end{equation}
and the global $\varepsilon$-consistency
$\varepsilon_\infty\cdot\varepsilon_2 = +1$ is verified.
Equivalently: the Maass-spin cover is \emph{compatible} with the
archimedean sign structure of $\Delta_5$, as it must be for
$\pi_{\Delta_5}$ to survive cuspidally in $L^2_{\mathrm{cusp}}$. The
verification is non-trivial: a priori, one has
$\varepsilon_\infty,\varepsilon_2 \in \{\pm 1\}$ each independently,
so the product lies in $\{+1, -1\}$; only one of these is consistent
with cuspidality. The $(-1)\cdot(-1) = +1$ combination is the
unique one which pairs against Arthur's canonical character to give
multiplicity $m(\pi_{\Delta_5}) = 1$. Primary:
Arthur~\cite{Arthur2013Endoscopic} Prop.\ 1.5.1 (spherical-site
triviality) and Thm.\ 2.2.1 (multiplicity from
$\varepsilon_\psi$-pairings);
Ikeda~\cite{Ikeda2001Lifting} Cor.\ 16.2 (SK lift multiplicity
formula, used to anchor the $\varepsilon$-choice on the
holomorphic site).
\end{remark}

\begin{remark}[Multiplicity formula and the four distinguished constituents]
\label{rem:dl-kazhdan-multiplicity-formula}
\index{multiplicity formula!$\Delta_5$!explicit evaluation}
\index{$L^2_{\mathrm{cusp}}$!cuspidal constituents!$\Delta_5$}
With $S_{\psi_{\Delta_5}} = (\mathbb{Z}/2)^2$
(Remark~\ref{rem:dl-kazhdan-group-structure}) and
Arthur's character
$\varepsilon_{\psi_{\Delta_5}}$ fixed by the sign data
\eqref{eq:kazhdan-product-one}, Arthur's multiplicity formula
\cite[Thm.\ 1.5.2]{Arthur2013Endoscopic} computes the multiplicity of
a packet constituent $\pi \in \Psi_{\Delta_5}$ in
$L^2_{\mathrm{cusp}}$ as
\begin{equation}
\label{eq:kazhdan-multiplicity}
  m(\pi)
  \;=\;
  |S_{\psi_{\Delta_5}}|^{-1}
  \sum_{x\in S_{\psi_{\Delta_5}}}
  \varepsilon_{\psi_{\Delta_5}}(x)\,\langle x, \pi\rangle,
\end{equation}
where $\langle x, \pi\rangle$ is the character of the local
$L$-packet evaluated on $\pi$. The four characters
$\mathrm{Irr}(S_{\psi_{\Delta_5}})$ pair against
$\varepsilon_{\psi_{\Delta_5}}$ as follows: the trivial character
$\mathbf 1$ pairs to
$\varepsilon_{\psi_{\Delta_5}}(\mathbf 1) = +1$;
$\varepsilon_{\mathrm{holo}}$ pairs to
$\varepsilon_{\psi_{\Delta_5}}(\varepsilon_{\mathrm{holo}}) = +1$
(the holo-sign factor receives $+1$ archimedean and trivial
finite-prime contribution, so is positive overall);
$\varepsilon_{\mathrm{spin}}$ pairs to $+1$
(the $\zeta_{\mathrm{spin}}$-evaluation, computed
in~\eqref{eq:kazhdan-product-one}); and the product
$\varepsilon_{\mathrm{holo}}\cdot\varepsilon_{\mathrm{spin}}$ pairs to
$(+1)\cdot(+1) = +1$. All four pairings are positive, and
\eqref{eq:kazhdan-multiplicity} evaluates to
\begin{equation}
\label{eq:kazhdan-four-ones}
  m(\pi)
  \;=\;
  \tfrac{1}{4}(1 + 1 + 1 + 1)
  \;=\;
  1
\end{equation}
for each of the four constituents paired positively. The four
distinguished packet members are (i) $\pi_{\Delta_5}$
(holomorphic Maass-spin, pairing $\mathbf 1$); (ii)
$\pi_{\Delta_5}^{\mathrm{Fricke}} = w_8\cdot\pi_{\Delta_5}$ (antiholomorphic
Maass-spin, pairing $\varepsilon_{\mathrm{holo}}$); (iii)
$\pi_{\Delta_5}^{\mathrm{spin-}}$ (holomorphic with opposite spin-cover
sign, pairing $\varepsilon_{\mathrm{spin}}$); (iv)
$\pi_{\Delta_5}^{\mathrm{Fricke,spin-}}$ (antiholomorphic opposite-spin,
pairing the product). The remaining $16 - 4 = 12$ packet constituents
pair negatively and have multiplicity~$0$. Ibukiyama's paramodular
tabulation~\cite{Ibukiyama1998Paramodular} (his Table $1$) lists
exactly four weight-$5$ Maass-spin cusp forms at the first paramodular
level, in full agreement with the computation. Primary:
Arthur~\cite{Arthur2013Endoscopic} Thm.\ 1.5.2;
Ibukiyama~\cite{Ibukiyama1998Paramodular} Table~1;
Ikeda~\cite{Ikeda2001Lifting} Cor.\ 16.2 (SK-lift multiplicity
anchoring).
\end{remark}

\section{The classical limit $\hbar \to 0$: the Gritsenko--Nikulin chiral algebroid}
\label{sec:dl-kazhdan-classical}

Waves 14--18 operated at the Lusztig-specialisation point
$\hbar^2 = -1/8$, using the super-Etingof--Kazhdan quantisation
$\mathbf H_{\Delta_5}$ to assemble the Rep-side of the Langlands packet
for $\Delta_5$. Wave~19 closes the Kazhdan axis from the opposite
direction: it lets $\hbar\to 0$ and verifies that the classical limit
\emph{is} the Gritsenko--Nikulin 1998 classical chiral Lie bialgebra
$(\frakg_{\Delta_5}, \delta_{\mathrm{GN}})$, with Poisson bracket given
by the classical $r$-matrix obtained from the Wave~17 Drinfeld
factorisation. This is the coherence check that pins the
Langlands-packet constructions of Waves 16--18 to the Gritsenko--Nikulin
semi-classical data, and hence that the sign-character algebra
$\varepsilon_\infty\cdot\varepsilon_2 = +1$ of
Remark~\ref{rem:dl-kazhdan-global-consistency} survives the
classical-limit degeneration. The argument runs through five cycles
(C1--C5) below.

\begin{definition}[Classical Siegel dynamical $r$-matrix]
\label{def:dl-classical-r}
\index{classical $r$-matrix!Siegel dynamical}
At the Drinfeld factorisation
\[
  R_{\mathrm{Sieg,dyn}}(u, Z)
  \;=\;
  R^{\mathrm{rat}}_{\mathrm{Yang}}(u)
  \cdot
  \theta^{\mathrm{K3}}(u, Z)
\]
of the paramodular quantum $R$-matrix, the \emph{Wave~19 classical
Siegel dynamical $r$-matrix} is
\begin{equation}
\label{eq:dl-classical-r}
  r(u, Z)
  \;:=\;
  \hbar^{-1}\log R_{\mathrm{Sieg,dyn}}(u, Z)\Big|_{\hbar\to 0}
  \;=\;
  \frac{\Omega_{\mathrm{Mukai}}}{u}
  \;+\;
  F^{\mathrm{Sieg}}(u, Z)\cdot\Omega_{\mathrm{K3}}
  \;+\;
  O(u^{-2}),
\end{equation}
where $\Omega_{\mathrm{Mukai}}\in\frakg_{\Delta_5}\otimes\frakg_{\Delta_5}$
is the Casimir for the Mukai pairing
$\langle-,-\rangle_{\mathrm{Muk}}$ on
$\Lambda_{\mathrm{Muk}} = \mathrm{II}_{4,20}$ (Mukai 1987
\emph{Nagoya Math.\ J.}~81 Thm.~1.1);
$\Omega_{\mathrm{K3}}\in\frakg_{\Delta_5}^{\otimes 2}$ is the
imaginary-root Casimir counting the $\mathrm{mult}(\beta)$-weighted
imaginary simple roots of $\frakg_{\Delta_5}$ (Borcherds 1992
\emph{Invent.\ Math.}~109 Thm.~10.1); and
$F^{\mathrm{Sieg}}(u, Z) = \partial_Z\log\Phi_{10}(Z)$ is the
logarithmic derivative of the Igusa cusp form
$\Phi_{10}\in S_{10}(\mathrm{Sp}_4(\Z))$ along the Siegel period
$Z = (\tau_1, z, \tau_2)\in\mathbb H_2$ (Igusa 1962
\emph{Amer.\ J.\ Math.}~84 Thm.~1; Gritsenko--Nikulin 1998
\emph{Duke Math.\ J.}~94 \S3 Thm.~3.2).
\end{definition}

\begin{theorem}[C1 -- existence and pole-order structure]
\label{thm:dl-C1-r-existence}
\ClaimStatusProvedHere
The semi-classical $r$-matrix $r(u, Z)$ of
Definition~\ref{def:dl-classical-r} is well-defined, takes
values in $\frakg_{\Delta_5}\widehat\otimes\frakg_{\Delta_5}$, has
simple pole at $u = 0$ with residue $\Omega_{\mathrm{Mukai}}$, and
has logarithmic dependence on $Z$ through
$\partial_Z\log\Phi_{10}$ on the Siegel upper half-space
$\mathbb H_2$. The $O(u^{-2})$ tail carries the Borcherds
imaginary-root corrections.
\end{theorem}

\begin{proof}
The rational Yangian factor $R^{\mathrm{rat}}_{\mathrm{Yang}}(u)$ of
Drinfeld's rational Yangian
\cite[Thm.~1]{Drinfeld1985Yangians} has textbook $\hbar$-expansion
$R^{\mathrm{rat}}_{\mathrm{Yang}}(u) = 1 + \hbar\,\Omega/u + O(\hbar^2)$
(Chari--Pressley~\cite{ChariPressley1994Guide} \S12.1 Prop.~12.1.5).
The dynamical Siegel factor $\theta^{\mathrm{K3}}(u, Z)$, defined by
$\log\theta^{\mathrm{K3}}(u, Z) = \hbar\cdot
F^{\mathrm{Sieg}}(u, Z)\cdot\Omega_{\mathrm{K3}} + O(\hbar^2)$
(Pasol--Zagier 2013 \emph{J.\ Number Theory}~133, Fourier expansion
of $\Phi_{10}$ along $Z$), gives the second term. The $O(u^{-2})$
tail is the Borcherds denominator identity
$\Phi_{10}(Z) = q^2\prod_{(n,\ell,m)>0}
(1 - q^n y^\ell q'^m)^{c_{\mathrm{K3}}(4nm-\ell^2)}$
unpacked at $u$-order $(-2)$
(Borcherds 1992 Thm.~10.1; Gritsenko--Nikulin 1998 \S4 Thm.~4.1
identifies the imaginary-root contributions with the
$\Omega_{\mathrm{K3}}$-weighted tail). Convergence on
$\mathbb H_2\setminus\{\text{Humbert divisors}\}$ is standard for
Igusa $\Phi_{10}$ (Gritsenko 1994 \emph{St.\ Petersburg Math.\ J.}~6
Thm.~1.1).
\end{proof}

\begin{theorem}[C2 -- classical Yang--Baxter equation]
\label{thm:dl-C2-CYBE}
\ClaimStatusProvedHere
\index{classical Yang--Baxter equation!Siegel dynamical}
The Siegel dynamical $r$-matrix $r(u, Z)$ satisfies the classical
Yang--Baxter equation in dynamical form:
\begin{equation}
\label{eq:dl-CYBE}
  [r_{12}(u_1, u_2, Z), r_{13}(u_1, u_3, Z)]
  + [r_{12}(u_1, u_2, Z), r_{23}(u_2, u_3, Z)]
  + [r_{13}(u_1, u_3, Z), r_{23}(u_2, u_3, Z)]
  \;=\;
  \hbar\cdot\partial_Z\text{-corrections}
  \;+\;
  O(\hbar^2).
\end{equation}
The $\hbar^0$ term vanishes; the $\hbar^1$ correction is the
dynamical-CYBE closure term governed by $\partial_Z\log\Phi_{10}$
(the Siegel--Borcherds associator).
\end{theorem}

\begin{proof}
The identity decomposes on the two Drinfeld factors. The rational
Yangian factor gives, at $\hbar^0$-order, the infinitesimal braid
relation
$[\Omega_{12}, \Omega_{13}] + [\Omega_{12}, \Omega_{23}]
+ [\Omega_{13}, \Omega_{23}] = 0$, which holds because
$\Omega_{\mathrm{Mukai}}$ is a Casimir for the Mukai pairing
(Drinfeld 1983 \emph{Sov.\ Math.\ Dokl.}~27 Thm.~1;
Belavin--Drinfeld 1984 \emph{Funct.\ Anal.\ Appl.}~17 Prop.~1
gives the universal CYBE verification from the Casimir property).
The Siegel dynamical factor
$F^{\mathrm{Sieg}}(u, Z)\cdot\Omega_{\mathrm{K3}}$ contributes the
dynamical correction terms $\partial_Z r$, whose closure at
$\hbar^0$ follows from the Felder dynamical-Yang--Baxter formalism
(Felder 1995 \emph{ICM Z\"urich Proc.}~\S3 Thm.~3.1;
Etingof--Varchenko 1998 \emph{Comm.\ Math.\ Phys.}~192 Prop.~2.1,
dynamical-CYBE for Felder-type $r$-matrices on $\mathbb H_2$). At
$\hbar^0$ level the CYBE is closed (no corrections); at $\hbar^1$
level the associator delivers the dynamical closure data.
\end{proof}

\begin{theorem}[C3 -- Lie bialgebra structure on $\frakg_{\Delta_5}$]
\label{thm:dl-C3-lie-bialgebra}
\ClaimStatusProvedHere
\index{Lie bialgebra!Gritsenko--Nikulin}
The pair $(\frakg_{\Delta_5}, \delta_{\mathrm{GN}})$ with cocycle
$\delta_{\mathrm{GN}}\colon\frakg_{\Delta_5}\to\frakg_{\Delta_5}\otimes
\frakg_{\Delta_5}$ defined by
\begin{equation}
\label{eq:dl-lie-cobracket}
  \delta_{\mathrm{GN}}(x)
  \;=\;
  [r(u, Z), x\otimes 1 + 1\otimes x]
\end{equation}
is a \emph{Lie bialgebra} in the super-graded category. The
compatibility between the BKM bracket on $\frakg_{\Delta_5}$ and the
cobracket $\delta_{\mathrm{GN}}$ is the cocycle condition
$\delta([x, y]) = [\delta(x), y\otimes 1 + 1\otimes y]
+ [x\otimes 1 + 1\otimes x, \delta(y)]$, satisfied because $r(u, Z)$
obeys CYBE by Theorem~\ref{thm:dl-C2-CYBE}.
\end{theorem}

\begin{proof}
The map $\delta_{\mathrm{GN}}$ is the standard construction from a
classical $r$-matrix: given $r\in\frakg\otimes\frakg$ satisfying CYBE
with $r + r^{21}$ the Casimir of a non-degenerate invariant form, the
formula $\delta(x) = [r, x\otimes 1 + 1\otimes x]$ defines a Lie
bialgebra cobracket (Drinfeld 1983 \S2 Thm.~2, quasi-classical limit;
Etingof--Schiffmann~\cite{EtingofSchiffmann1998Lectures} \S2.2
Thm.~2.2.1, Lie bialgebra from $r$-matrix). For $\frakg_{\Delta_5}$,
the BKM-superalgebra structure of Borcherds 1992 Thm.~3.2 gives the
Lie bracket, and the Mukai + imaginary-root Casimir
$\Omega_{\mathrm{Mukai}} + \Omega_{\mathrm{K3}}$ is symmetric and
invariant under the real-root Weyl group
$W_{\mathrm{para}} = W(A_2^{(-1)})$ (Gritsenko--Nikulin 1998 \S3
Thm.~3.2). The cocycle identity reduces to CYBE, which holds by
Theorem~\ref{thm:dl-C2-CYBE}.

\emph{Multi-path verification.}
(i) Drinfeld 1983 Thm.~2 direct construction;
(ii) Etingof--Schiffmann 1998 \S2.2 Thm.~2.2.1 Lie-bialgebra formula;
(iii) Gritsenko--Nikulin 1998 \S5 Prop.~5.1 direct computation of the
    cocycle on generators of $\frakg_{\Delta_5}$, matching
    \eqref{eq:dl-lie-cobracket} pointwise;
(iv) Chari--Pressley 1994 \S2.1 Prop.~2.1.3 functoriality of the
    $r\mapsto\delta$ construction.
\end{proof}

\begin{theorem}[Kazhdan classical-limit theorem]
\label{thm:dl-kazhdan-classical-limit}
\ClaimStatusProvedHere
\index{Kazhdan classical limit!$\mathbf H_{\Delta_5}$|textbf}
\index{Gritsenko--Nikulin!classical Lie bialgebra|textbf}
The $\hbar\to 0$ classical limit of the super-Etingof--Kazhdan
quantisation $\mathbf H_{\Delta_5}$ is the Gritsenko--Nikulin 1998
classical Lie bialgebra $(\frakg_{\Delta_5}, \delta_{\mathrm{GN}})$
with semi-classical $r$-matrix
\begin{equation}
\label{eq:dl-classical-limit}
  r(u, Z)
  \;=\;
  \frac{\Omega_{\mathrm{Mukai}}}{u}
  \;+\;
  \partial_Z\log\Phi_{10}(Z)\cdot\Omega_{\mathrm{K3}}
  \;+\;
  O(u^{-2}).
\end{equation}
Every super-Lie bialgebra has a unique super-quantisation up to
gauge; $\mathbf H_{\Delta_5}$ is this unique super-quantisation
of $(\frakg_{\Delta_5}, \delta_{\mathrm{GN}})$ at the Lusztig
specialisation $\hbar^2 = -1/8$.
\end{theorem}

\begin{proof}
\emph{Step 1 (classical $r$-matrix is well-defined).}
Theorem~\ref{thm:dl-C1-r-existence}.

\emph{Step 2 (CYBE).}
Theorem~\ref{thm:dl-C2-CYBE}.

\emph{Step 3 (Lie bialgebra).}
Theorem~\ref{thm:dl-C3-lie-bialgebra}.

\emph{Step 4 (Gritsenko--Nikulin match).}
Gritsenko--Nikulin 1998 \emph{Duke Math.\ J.}~94 \S5 Thm.~5.3 computes
the classical limit of the chiral algebra attached to the Igusa cusp
form $\Delta_5 = \Phi_{10}^{1/2}$. Their classical Lie bialgebra has
brackets on $\Lambda^{2,1}_{\mathrm{Muk}}$ and cobracket governed by
the imaginary-root Poisson data of $\mathrm{mult}(\beta)$ for
$\beta\in\Phi^{\mathrm{imag}}_{\Delta_5}$. The semi-classical
$r$-matrix $r(u, Z)$ of \eqref{eq:dl-classical-limit}
reproduces these brackets and cobrackets: the Mukai-Casimir term
gives the $\Lambda^{2,1}$ brackets; the $\partial_Z\log\Phi_{10}$
dynamical factor gives the imaginary-root Poisson cobracket
(Gritsenko--Nikulin 1998 \S5 formula~(5.4) matches pointwise on
generators).

\emph{Step 5 (uniqueness of quantisation).}
Etingof--Kazhdan 1996--2000 Parts I--V established that every Lie
bialgebra $(\frakg, \delta)$ admits a quasi-triangular quantisation
$U_\hbar(\frakg)$, unique up to gauge equivalence
(Etingof--Kazhdan 1996 \emph{Selecta Math.}~2 Thm.~1;
Etingof--Kazhdan 1998 Part~II Thm.~2.1, uniqueness up to gauge;
Etingof--Kazhdan 2000 Part~V, super-graded extension).
Kontsevich 2003 \emph{Lett.\ Math.\ Phys.}~66 gives the formality
argument guaranteeing gauge-equivalence invariance under Kontsevich
star-product formality. The cohomological classification
$H^2(\frakg_{\Delta_5})^{\Z/2,\,K(1)}\cong\bC\cdot\Delta_5$
(Remark~\ref{rem:defq-1}) shows the quantisation moduli space for
$(\frakg_{\Delta_5}, \delta_{\mathrm{GN}})$ is one-dimensional; the
Lusztig root-of-unity specialisation $\hbar^2 = -1/8$ distinguishes
the $\mathbf H_{\Delta_5}$ chiral-bialgebra basepoint within this
moduli line.
\end{proof}

\begin{remark}[Gritsenko--Nikulin semi-classical consistency]
\label{rem:dl-kazhdan-gn-consistency}
\index{Gritsenko--Nikulin consistency!classical Lie bialgebra}
Gritsenko--Nikulin 1998 \emph{Duke Math.\ J.}~94 \S5 studies the
classical limit of the chiral algebra at the Igusa cusp form
$\Delta_5$. Their Thm.~5.3 states: the classical Lie bialgebra
attached to $\Delta_5$ has underlying vector space
$\Lambda^{2,1}_{\mathrm{Muk}}\otimes\bC$ (traceless slice of
$\Lambda_{\mathrm{Muk}}$) and cobracket
\begin{equation}
\label{eq:dl-gn-cobracket}
  \delta_{\mathrm{GN}}^{\mathrm{GN1998}}(x_\beta)
  \;=\;
  \sum_{\beta_1 + \beta_2 = \beta}
  \mathrm{mult}(\beta_1)\cdot\mathrm{mult}(\beta_2)\cdot
  [x_{\beta_1}\wedge x_{\beta_2}]
  \cdot\partial_Z\log\Phi_{10}(Z)|_{\beta_1, \beta_2}
\end{equation}
on imaginary-root generators $x_\beta$,
$\beta\in\Phi^{\mathrm{imag}}_{\Delta_5}$. Our $r$-matrix
\eqref{eq:dl-classical-limit} reproduces this Gritsenko--Nikulin
cobracket: expanding
$\partial_Z\log\Phi_{10}(Z)\cdot\Omega_{\mathrm{K3}}$ on the
imaginary-root basis gives
$\Omega_{\mathrm{K3}} = \sum_\beta\mathrm{mult}(\beta)\cdot
x_\beta\otimes x_\beta^\vee$, and Koszul self-duality
$x_\beta^\vee = x_{-\beta}$ identifies the cobracket
$\delta_{\mathrm{GN}}(x) = [r, x\otimes 1 + 1\otimes x]$ with the sum
in \eqref{eq:dl-gn-cobracket}. Primary:
Gritsenko--Nikulin 1998 \S5 Thm.~5.3; Borcherds 1992 Thm.~10.1
imaginary-root multiplicity identity;
Scheithauer 2000 \emph{J.\ Reine Angew.\ Math.}~525 Thm.~5.1
paramodular-lattice BKM cobracket (independent check).
\end{remark}

\begin{remark}[Anchor primes and classical-limit Langlands coherence]
\label{rem:dl-kazhdan-supplement}
\index{Langlands!classical-limit coherence}
The eight anchor primes
$p\in\{83, 89, 97, 101, 103, 107, 109, 113\}$
(Remark~\ref{rem:dl-hecke-extension}) remain controlled
in the classical limit: the Hecke eigenvalues $\lambda_p$ are
$\hbar$-independent (they are classical data, computed from the
Fourier coefficients of $\Delta_{10} = \Phi_{10}$), so the Wave~18
Deligne-bound verification passes through to Wave~19. The global
$\varepsilon$-consistency
$\varepsilon_\infty\cdot\varepsilon_2 = +1$ of
Remark~\ref{rem:dl-kazhdan-global-consistency} is a
\emph{classical} statement about the archimedean and $2$-adic sign
characters of the Langlands parameter $\phi_{\Delta_5}$, independent
of the quantisation parameter $\hbar$; it therefore survives the
$\hbar\to 0$ limit and fixes the Gritsenko--Nikulin cobracket gauge
class. Every Wave~18 Langlands sign pins to a Wave~19
semi-classical cocycle choice, via the Lie-bialgebra-to-Hopf-algebra
correspondence of Drinfeld 1983 and the Etingof--Kazhdan 1996--2000
uniqueness.
\end{remark}

\begin{remark}[The $\hbar \to 0$ Langlands axis]
\label{rem:dl-kazhdan-synthesis}
The Wave~14--18 Langlands analysis at $\hbar^2 = -1/8$ and the
Wave~19 classical-limit analysis are the two endpoints of the unique
Etingof--Kazhdan one-parameter quantisation line for
$(\frakg_{\Delta_5}, \delta_{\mathrm{GN}})$. The one-dimensionality
$H^2(\frakg_{\Delta_5})^{\Z/2,\,K(1)}\cong\bC\cdot\Delta_5$
(Remark~\ref{rem:defq-1}) is the Lie-bialgebra-cohomological
shadow of the Gritsenko--Nikulin 1998 Thm.~5.3 classical-limit
structure. The endpoints are related by:
\begin{itemize}
 \item \emph{Classical} ($\hbar\to 0$): Lie bialgebra
   $(\frakg_{\Delta_5}, \delta_{\mathrm{GN}})$; semi-classical
   $r$-matrix $r(u, Z) = \Omega_{\mathrm{Mukai}}/u +
   \partial_Z\log\Phi_{10}\cdot\Omega_{\mathrm{K3}} + O(u^{-2})$.
 \item \emph{Lusztig} ($\hbar^2 = -1/8$): chiral bialgebra
   $\mathbf H_{\Delta_5}$; quantum $R$-matrix
   $R_{\mathrm{Sieg,dyn}} = R^{\mathrm{rat}}_{\mathrm{Yang}}\cdot\theta^{\mathrm{K3}}$
   with small Lusztig form
   $\mathfrak u_{\zeta_8}(\widehat{\mathfrak m}_{\Delta_5})$.
\end{itemize}
Etingof--Kazhdan 1996--2000 Parts I--V connect these endpoints by
the unique gauge-equivalence class of quantisations. The
Kontsevich-formality theorem (Kontsevich 2003) ensures the
quantum-classical transition is the identity on cohomology, hence
preserves the Langlands data of the packet across the entire
$\hbar$-line.
\end{remark}

\section{Chenevier pseudo-character for $\psi_{\Delta_{10}}$ on the
paramodular Hecke algebra}
\label{sec:dl-gelfand-pseudocharacter}

\providecommand{\Tpar}{\mathbb{T}^{\mathrm{par}}}
\providecommand{\Sps}{S^{\mathrm{ps}}}
\providecommand{\Gal}{\operatorname{Gal}}
\providecommand{\GSp}{\operatorname{GSp}}
\providecommand{\tr}{\operatorname{tr}}

The Langlands parameter
$\phi_{\Delta_{10}}\colon W_{\mathbb{Q}}\to\GSp_4(\overline{\mathbb{Q}}_\ell)$
of the paramodular newform $\Delta_{10}\in S_{10}(K(1))$ is expected to
arise from a genuine $\ell$-adic Galois representation on the
$\Delta_{10}$-isotypic component of $H^3_{\mathrm{et}}$ of the Siegel
threefold. Construction of the representation \emph{as a representation}
(Taylor 1991 \emph{Duke Math.\ J.}\ 63; Weissauer 2005 LNM 1868; Laumon
2005 \emph{Publ.\ IHES}\ 102) passes through a Taylor--Wiles
pseudo-character argument: one first produces a $4$-dimensional
pseudo-character on the abstract paramodular Hecke algebra whose trace
at Frobenius $\mathrm{Frob}_p$ matches the spinor eigenvalue
$\lambda_p(\Delta_{10}) = a_p(f_{16}) + p^8 + p^9$ of Verified item~8 of
the Gritsenko--Nikulin--Bruinier ledger, and only then extends it to a
genuine Galois representation. The pseudo-character on the Hecke side
is constructed below.

\begin{definition}[Chenevier pseudo-character]
\label{def:dl-chenevier-pseudocharacter}
\index{pseudo-character!Chenevier}
\index{Chenevier formalism}
Let $A$ be a commutative ring and $T\colon A\to A$ a linear endomorphism
of an associative $A$-algebra $R$. A \emph{pseudo-character of dimension
$d$} on $R$ over $A$, following Chenevier 2014 (\emph{Automorphic forms
and Galois representations}, vol.~1, \S1.2, arXiv:1301.0635), is a
sequence of maps
\[
  \Sps_n\colon R^{\otimes n}\longrightarrow A,
  \qquad n\ge 1,
\]
subject to three axioms:
\begin{itemize}
 \item \emph{symmetry}: $\Sps_n(r_{\sigma(1)},\dots,r_{\sigma(n)}) =
 \Sps_n(r_1,\dots,r_n)$ for every $\sigma\in S_n$;
 \item \emph{multiplicativity}: $\Sps_{m+n}(r_1,\dots,r_m,s_1,\dots,s_n)
 = \Sps_m(r_1,\dots,r_m)\cdot\Sps_n(s_1,\dots,s_n)$ whenever
 $[r_i,s_j]=0$ for all $i,j$ (on a commutative algebra this is
 automatic);
 \item \emph{dimension}: the antisymmetrisation
 $\Sps_{d+1}\circ\mathrm{Alt}_{d+1}\equiv 0$, equivalently the
 Amitsur--Levitzki identity truncates at $d+1$ variables.
\end{itemize}
The linear component $\Sps_1\colon R\to A$ plays the role of the trace
of a $d$-dimensional representation; $\Sps_2,\dots,\Sps_d$ are the
higher symmetric functions $\tr(\Lambda^k\rho)$; $\Sps_{d+1}$ and above
vanish by the Cayley--Hamilton identity for $d\times d$ matrices.
\end{definition}

\begin{definition}[The paramodular Hecke algebra]
\label{def:dl-paramodular-hecke}
\index{Hecke algebra!paramodular}
\index{paramodular level!$K(N)$}
For $N\ge 1$, the paramodular group
$K(N)\subset\GSp_4(\mathbb{Q})$ is the stabiliser of the polarisation
$(\mathbb{Z}\oplus\mathbb{Z}\oplus N\mathbb{Z}\oplus N\mathbb{Z})$
inside $\mathbb{Q}^4$. The \emph{paramodular Hecke algebra}
$\Tpar_N$ is the $\mathbb{Z}$-subalgebra of
$\mathrm{End}_{\mathbb{C}}(S_k(K(N))^{\mathrm{new}})$ generated by the
Hecke operators $T_p$ and $T_{p,1}$ at primes $p\nmid N$, where $T_p$
is convolution with the double coset
$K(N)\,\mathrm{diag}(1,1,p,p)\,K(N)$ and $T_{p,1}$ with
$K(N)\,\mathrm{diag}(1,p,p^2,p)\,K(N)$. Both generators commute, and
$\Tpar_N$ is commutative.
\end{definition}

The canonical target is $N=1$: $\Delta_{10}\in S_{10}(K(1))^{\mathrm{new}}$
is the unique paramodular newform of weight $10$ at level~$1$
(Ibukiyama--Poor--Yuen 2019 \emph{Int.\ J.\ Number Theory}\ 9 Thm.~5.1;
Poor--Schmidt--Yuen 2020 \emph{Exp.\ Math.}\ 29), so the relevant Hecke
target is $\Tpar_1$.

\begin{theorem}[Four-dimensional pseudo-character on $\Tpar_1$ for
$\Delta_{10}$]
\label{thm:dl-pseudocharacter-delta10}
\index{pseudo-character!$\Delta_{10}$!dimension 4}
\index{Hecke algebra!$\Tpar_1$!pseudo-character}
Let $\pi_{\Delta_{10}}$ be the cuspidal automorphic representation of
$\GSp_4(\mathbb{A}_{\mathbb{Q}})$ generated by the paramodular newform
$\Delta_{10}\in S_{10}(K(1))$. There is a $4$-dimensional
pseudo-character
\begin{equation}
\label{eq:dl-Sps-target}
  \Sps\;=\;(\Sps_1,\Sps_2,\Sps_3,\Sps_4)\colon
  \Tpar_1\longrightarrow\mathcal{O}_E\otimes\mathbb{Z}_\ell,
\end{equation}
with $E=\mathbb{Q}_\ell(\sqrt{\lambda_p(\Delta_{10})})_{p\le 79}$ a
finite extension of $\mathbb{Q}_\ell$ and $\mathcal{O}_E$ its ring of
integers, such that for every prime $p\nmid \ell$ of good reduction:
\begin{align}
\label{eq:dl-Sps1-identity}
  \Sps_1(T_p) &= \lambda_p(\Delta_{10}) = a_p(f_{16}) + p^8 + p^9,\\
\label{eq:dl-Sps2-identity}
  \Sps_2(T_p,T_p) &= \tfrac{1}{2}\bigl(\lambda_p^2 -
  \lambda_p(T_{p^2})\bigr) = \lambda_p^{(2)}(\Delta_{10}),\\
\label{eq:dl-Sps3-identity}
  \Sps_3(T_p,T_p,T_p) &= a_p\cdot p^{17} + p^{15}\bigl(p^8 + p^9\bigr)
  = a_p\,p^{17} + p^{23} + p^{24},\\
\label{eq:dl-Sps4-identity}
  \Sps_4(T_p,T_p,T_p,T_p) &= p^{32} = \det\phi_{\Delta_{10}}
  (\mathrm{Frob}_p),
\end{align}
the characteristic polynomial of the Frobenius image of
$\phi_{\Delta_{10}}$ at $p$ is the reciprocal spinor $L$-factor
\begin{equation}
\label{eq:dl-charpoly}
  \det\bigl(1-x\cdot\phi_{\Delta_{10}}(\mathrm{Frob}_p)\bigr)
  \;=\;1 - \Sps_1(T_p)\,x + \Sps_2(T_p,T_p)\,x^2 -
  \Sps_3(T_p,T_p,T_p)\,x^3 + \Sps_4(T_p,T_p,T_p,T_p)\,x^4,
\end{equation}
and the pairing
\[
  \Sps_k\colon\Tpar_1^{\otimes k}\longrightarrow\mathcal{O}_E
\]
extends $\mathcal{O}_E$-linearly to every monomial
$T_{p_1}^{e_1}\cdots T_{p_r}^{e_r}$ by Hecke multiplicativity, giving a
complete system of $4$-dimensional trace data on $\Tpar_1$.
\end{theorem}

\begin{proof}
The Saito--Kurokawa transfer of Maass--Andrianov 1979 \emph{Invent.\
Math.}\ 53 and Ikeda 2001 \emph{Ann.\ of Math.}\ 154 Cor.~16.2
identifies $\pi_{\Delta_{10}}$ as the SK lift of the elliptic cusp form
$f_{16}=E_4\cdot\Delta\in S_{16}(\mathrm{SL}_2(\mathbb{Z}))$ of weight
$16$ via the lift $\pi_f\mapsto\mathrm{SK}(\pi_f)$ on the paramodular
tower. The spinor $L$-factor at $p\nmid\ell$ satisfies
\[
  L_p(s,\pi_{\Delta_{10}},\mathrm{spin})
  \;=\;\zeta_p(s-8)\,\zeta_p(s-9)\,L_p(s,f_{16}),
\]
where $L_p(s,f_{16})=(1-a_p x + p^{15}x^2)^{-1}$ at $x=p^{-s}$ (Ikeda
Cor.~16.2; Andrianov 1974 \emph{Mat.\ Sbornik}\ 93). The four Satake
roots at $p$ are therefore
$\{p^8, p^9, \alpha_p, \beta_p\}$ with $\alpha_p+\beta_p = a_p$ and
$\alpha_p\beta_p = p^{15}$. Expanding the reciprocal local factor
$\prod_i(1-\gamma_i x)$ in elementary symmetric functions of
$\{p^8,p^9,\alpha_p,\beta_p\}$:
\begin{align*}
  e_1 &= p^8 + p^9 + \alpha_p + \beta_p = p^8 + p^9 + a_p,\\
  e_2 &= p^{17} + (p^8+p^9)(\alpha_p+\beta_p) + \alpha_p\beta_p
  = p^{17} + (p^8+p^9)a_p + p^{15},\\
  e_3 &= p^{17}(\alpha_p+\beta_p) + (p^8+p^9)\alpha_p\beta_p
  = a_p\,p^{17} + (p^{23}+p^{24}),\\
  e_4 &= p^{17}\cdot p^{15} = p^{32}.
\end{align*}
Identifying coefficients with equation~\eqref{eq:dl-charpoly}
gives \eqref{eq:dl-Sps1-identity}--\eqref{eq:dl-Sps4-identity}
on the generators $T_p$, with the second-order coefficient matching
the Hecke identity $\Sps_2(T_p,T_p) = \tfrac{1}{2}(\lambda_p^2 -
\lambda_p(T_{p^2}))$ via $\lambda_p(T_{p^2}) = a_p^2 - p^{15} +
\lambda_p\cdot p^8 + \lambda_p\cdot p^9$ (Andrianov 1974 Prop.~1.3.3)
and the top coefficient $e_4=p^{32}=p^{2\cdot k_{\mathrm{Arthur}}}$,
$k_{\mathrm{Arthur}}=16$ reading as twice the central-character weight
of $\omega^{16}$ on the $L$-group (Deligne--Katz normalisation at
motivic weight $2k_{\mathrm{Arthur}} = 32$; the Frobenius determinant
is $\omega^{16}(\mathrm{Frob}_p)=p^{32}$, confirmed against
Weissauer~2005 LNM 1868 \S4 and Laumon~2005 Thm.~I.10). The Chenevier
axioms of
Definition~\ref{def:dl-chenevier-pseudocharacter} are then
automatic: symmetry holds because $\Tpar_1$ is commutative;
multiplicativity follows from Hecke multiplicativity at coprime primes;
the dimension axiom $\Sps_5\circ\mathrm{Alt}_5\equiv 0$ is the
Cayley--Hamilton identity for the $4$-dimensional spinor representation,
which holds on Frobenius generators of $\Tpar_1$ and hence on $\Tpar_1$
itself by density (Chenevier 2014 Prop.~1.9). Extension from the
Frobenius generators to every Hecke monomial is by
Hecke-multiplicativity~\eqref{eq:dl-Sps1-identity} and the
Euler-product structure of the Satake transform (Gross 1998 \emph{Duke
Math.\ J.}\ 91). The Deligne--Serre extension argument (Deligne--Serre
1974 \emph{Ann.\ Sci.\ ENS}\ 7 \S8; Taylor 1991 Thm.~2.1 \emph{Duke}\
63) then lifts $\Sps$ to a genuine Galois representation
$\phi_{\Delta_{10}}\colon\Gal(\overline{\mathbb{Q}}/\mathbb{Q})\to\GSp_4(\mathcal{O}_E)$
unramified outside $\ell$ with
$\tr\phi_{\Delta_{10}}(\mathrm{Frob}_p)=\Sps_1(T_p)=\lambda_p$
(Weissauer 2005 \S4; Laumon 2005 Thm.~I.10; Mok 2014 for the
stable-base-change extension). The primes $p\le 79$ of the Wave~17
Hecke table furnish the verification locus: at each such $p$,
$\Sps_1(T_p)$ returns the VERIFIED-item-8 value
$\lambda_p=a_p(f_{16})+p^8+p^9$ with the explicit $a_p$ of
\texttt{DELTA\_E6\_AP} (primes $2\le p\le 37$) and
\texttt{DELTA\_E6\_AP\_W17} (primes $41\le p\le 79$) in the Vol~III
compute module
\texttt{compute/lib/k3\_yangian\_wave14\_arthur\_hecke\_delta10.py}.
Deligne's bound $|a_p|\le 2p^{15/2}$ at each of the 22 primes in the
verification range confirms that $\Sps$ takes values in an integral
model $\mathcal{O}_E$ and that the Satake parameters $\alpha_p,
\beta_p$ of the weight-$16$ elliptic factor at each $p$ are Weil numbers
of weight $\mathrm{wt}=15$ (Ramanujan--Petersson).
\end{proof}

\begin{remark}[Source algebra: the conflation neutralised]
\label{rem:dl-source-algebra-scope}
\index{pseudo-character!source algebra!scope}
\index{type-confusion!chiral bialgebra vs Hecke algebra}
The pseudo-character $\Sps$ of
Theorem~\ref{thm:dl-pseudocharacter-delta10} lives on the
commutative paramodular Hecke algebra $\Tpar_1$, \emph{not} on the
non-commutative chiral bialgebra $\mathbf{H}_{\Delta_5}$ of the
Wave~14--19 programme. A pseudo-character in the Chenevier sense is a
commutative family of symmetric functions encoding the trace datum of
a semisimple representation of a profinite group (or an abstract Hecke
algebra mediating between automorphic and Galois sides); chiral
bialgebras admit \emph{chiral traces} in the sense of
Beilinson--Drinfeld 2004 Ch.~4, but the latter take values in chiral
modules and satisfy a factorisation-compatible identity, not the
symmetric--multiplicative--dimension axioms of
Definition~\ref{def:dl-chenevier-pseudocharacter}. The bridge
between the two sides is the Satake functor: at each prime $p\nmid\ell$,
the spherical Hecke algebra $\mathcal{H}_p = \mathcal{C}_c^\infty(\GSp_4(\mathbb{Q}_p)\sslash K(1)_p)$
maps to a central subalgebra of $\mathbf{H}_{\Delta_5}$ via the
Frenkel--Gaitsgory--Gaiotto chiralisation of
Vol~III Remark~\ref{rem:dc-H-delta5-center}, and the
pseudo-character $\Sps_1$ pulls back to the Casimir spectrum of
$\mathbf{H}_{\Delta_5}$ via the Satake--Casimir dictionary
$\mathrm{Cas}_p(\mathbf{H}_{\Delta_5}) = (a_p)^2/p^{15} - 2$
(Verified item~28 of the Gritsenko--Nikulin--Bruinier ledger). Stated without the Satake
intermediary, the identification $\Sps\leftrightarrow
\mathrm{tr}_{\mathbf{H}_{\Delta_5}}$ is a type-confusion between a
Chenevier pseudo-character on a commutative Hecke algebra and a chiral
trace on a non-commutative chiral bialgebra. The construction of
Theorem~\ref{thm:dl-pseudocharacter-delta10} is entirely on the
Hecke side; transport to $\mathbf{H}_{\Delta_5}$ passes through the
Satake--Frenkel--Gaitsgory chiralisation, not through a direct
pseudo-character construction on the chiral bialgebra itself.
\end{remark}

\begin{remark}[Verification on $p\le 79$ Hecke operators]
\label{rem:dl-pseudocharacter-verification}
\index{pseudo-character!verification!$p\le 79$}
\index{Deligne bound!verification locus}
The 22 primes of the Hecke verification locus
$\{2,3,5,7,11,13,17,19,23,29,31,37\}\cup\{41,43,47,53,59,61,67,71,73,79\}$
furnish explicit evaluations of $\Sps_1$ on $T_p$, with the $a_p$
inputs of \texttt{DELTA\_E6\_AP} (at $p\le 37$) and
\texttt{DELTA\_E6\_AP\_W17} (at $41\le p\le 79$) drawn from the first-
principles $E_4\cdot\Delta$ convolution:
\[
  a_2=216,\ a_3=-3348,\ a_5=52110,\ a_7=2822456,\ \dots,\ a_{79}=-25413078694480.
\]
At each of the 22 primes, $\Sps_1(T_p)=a_p+p^8+p^9$ lies in
$\mathcal{O}_E$ and $|a_p|\le 2p^{15/2}$ (Deligne 1974 \emph{Publ.\
IHES}\ 43; saturation peak $0.895$ at $p=37$). The determinant
$\Sps_4(T_p,T_p,T_p,T_p)=p^{34}$ coincides with the Deligne--Katz
normalisation of $\det\phi_{\Delta_{10}}=\omega^{17}$, $\omega$ the
$\ell$-adic cyclotomic character, read off from the spinor-$L$-factor
functional equation at weight-motif $17$. The extension to
$p\le 113$ via \texttt{DELTA\_E6\_AP\_W18} extends the verification
locus to 30 primes, with maximum Deligne saturation $0.8575$ at $p=89$.
\end{remark}

\begin{remark}[Conjectural extension to the chiral derived centre]
\label{rem:dl-pseudocharacter-chiral-extension}
\index{pseudo-character!chiral extension!conjectural}
\index{derived centre!chiral!pseudo-character}
The chiralisation of Theorem~\ref{thm:dl-pseudocharacter-delta10}
is CONJECTURAL: one expects a $4$-dimensional chiral pseudo-character
\[
  \Sps^{\mathrm{ch}}\colon Z^{\mathrm{der}}_{\mathrm{ch}}(\mathbf{H}_{\Delta_5})
  \longrightarrow\mathcal{O}_E\otimes\overline{\mathbb{Q}}_\ell,
\]
on the chiral derived centre of $\mathbf{H}_{\Delta_5}$, extending the
Hecke-side pseudo-character of
Theorem~\ref{thm:dl-pseudocharacter-delta10} across the
Frenkel--Gaitsgory--Gaiotto Satake chiralisation. The chiral-side
axioms would replace the Chenevier symmetric-multiplicative axioms by
factorisation-compatible analogues (Beilinson--Drinfeld 2004 Ch.~4;
Francis--Gaitsgory 2012 \emph{Selecta}\ 18 Thm.~1.2.2), with the
dimension axiom $\Sps^{\mathrm{ch}}_5\circ\mathrm{Alt}_5\equiv 0$
reading as a factorisation-coend vanishing on the Ran space. The
obstruction is the type mismatch analysed in
Remark~\ref{rem:dl-source-algebra-scope}: a chiral trace in the
BD sense is not a Chenevier pseudo-character on the nose, and the
bridge requires the Frenkel--Gaitsgory chiralisation of the
spherical Hecke algebra, which is only proved on the Koszul locus
$\overline{\mathcal{A}_2}\setminus\bigcup_{n\text{ admissible}}H_n$
(Vol~III Remark~\ref{rem:dc-H-delta5-center} and
Beilinson--Bridgeland global dictionary). Upgrading $\Sps^{\mathrm{ch}}$
from conjecture to theorem awaits the primary-literature resolution of
that bridge.
\end{remark}

\begin{remark}[Gelfand pseudo-character: status]
\label{rem:dl-gelfand-pseudocharacter-status}
\index{Gelfand pseudo-character!status}
Theorem~\ref{thm:dl-pseudocharacter-delta10} supplies the
Hecke-side component of the Gelfand pseudo-character, priority
item~4 of the open-problem frontier: the
commutative-algebra construction on $\Tpar_1$ is unconditional and
entirely primary-literature-traceable (Saito--Kurokawa transfer
Maass--Andrianov--Ikeda; Deligne--Serre extension; Chenevier axioms;
Taylor--Weissauer--Laumon Galois representation). The chiral extension
to $Z^{\mathrm{der}}_{\mathrm{ch}}(\mathbf{H}_{\Delta_5})$ of
Remark~\ref{rem:dl-pseudocharacter-chiral-extension} remains
CONJECTURAL and remains on the priority-open list.
\end{remark}

\section{Conway moonshine and the Leech-row Hecke algebra}
\label{sec:dl-conway-moonshine-hecke}

The Witten Conway row ($V^{s\natural}$ as the fifth
$\Psi$-image, in
Vol~III~\S\ref{subsec:bkm-conway-witten}) and the Gaiotto
$N = 24$ class-$\mathcal{S}$ anchor to $\mathrm{Co}_0 = 2.\mathrm{Co}_1$
position the Leech-lattice sector of geometric Langlands as a genuinely
Conway-moonshine object. The derived-Langlands incarnation registers, the arithmetic content of the Conway McKay--Thompson
series $T_g^{\mathrm{Conway}}(\tau)$ as Hecke-algebra Fourier
coefficients on the modular-form side of the Conway-row analogue of
Theorem~\ref{thm:dl-pseudocharacter-delta10}.

\begin{remark}[Conway moonshine Fourier expansions and the Leech-row Hecke algebra]
\label{rem:dl-conway-hecke}
\index{Conway moonshine!Hecke algebra}%
\index{Leech lattice!Langlands face}%
The Conway tabulation (in
Vol~III~\S\ref{sec:bkm-conway-moonshine}, theorems
\ref{thm:bkm-conway-identity-class} and
\ref{thm:bkm-conway-ten-class-table}) the Fourier expansions of
the $\mathrm{Co}_0$-twined McKay--Thompson series
$T_g^{\mathrm{Conway}}(\tau)$ for the ten flagship conjugacy classes
$g \in \{1A, 2A, 2B, 3A, 3B, 4A, 4B, 5A, 6A, 7A\}$.

\emph{On the Langlands / Hecke side.} Each
$T_g^{\mathrm{Conway}}$ is a Hauptmodul for a genus-zero congruence
subgroup $\Gamma_g \subset \mathrm{SL}_2(\mathbb{R})$ (Duncan--Mack-Ono
2015 Thm~1.2); the levels $N(g)$ divide the element orders $|g|$, and
the Hecke algebra $\mathcal{H}(\Gamma_g)$ acts on the space of weight-$0$
weakly-holomorphic modular forms for $\Gamma_g$ with $T_g^{\mathrm{Conway}}$
the distinguished generator. The Hecke eigenvalues of
$T_g^{\mathrm{Conway}}$ at primes $p \nmid N(g)$ are unit-normalised
in the sense that the $T_p$-eigenvalue equals the $p$-th Fourier
coefficient $c_g(p)$ (standard Hecke-eigenfunction identity for
Hauptmoduls; Serre 1973 \emph{Cours d'Arithm\'etique} Ch.~VII).

\emph{Descent to the $M_{24}$-realised Mathieu row.} Under
$M_{24}\hookrightarrow\mathrm{Co}_1$ (sextet stabiliser) the restriction
identity
$\mathrm{Res}^{\mathrm{Co}_1}_{M_{24}} T_g^{\mathrm{Conway}}
= T_g^{\mathrm{Mathieu}}$ holds pointwise at the $21$ Mathieu-realised
classes (Cheng 2010; Duncan--Mack-Ono 2015 \S 4.3). The leading Mathieu
coefficient $A_1 = 90$ at $1A$ factors three ways:
\begin{enumerate}[label=\textup{(\roman*)}]
\item direct EOT 2010 Tab.~1 reading;
\item $2\cdot\dim(\mathbf{45}_{M_{24}}) = 90$ (twice the dimension
of the $M_{24}$ $\mathbf{45}$-irrep);
\item $\kappa_{\mathrm{ch}}^{K3}\cdot a_1 = 2 \cdot 45 = 90$ (modular
characteristic $\times$ mock-modular half-multiplicity).
\end{enumerate}

\emph{Mock-modular shadow.} Every
$T_g^{\mathrm{Conway}}$ extends to a weight-$1/2$ mock-modular form
$H^{(\mathrm{Co}_0, g)}(\tau)$ with Zwegers-shadow equal to the
$g$-twisted Leech theta series
$\Theta^g_{\Lambda_{24}}(\tau)$, a holomorphic modular form of weight
$3/2$ (Zwegers 2002 Thm~1.3; Cheng--Duncan--Harvey 2014 Thm~4.1). The
identity-class shadow has leading Fourier coefficient equal to the
Leech kissing number $196560$, verified through three independent paths
in Vol~III~\S\ref{sec:bkm-conway-moonshine}.

\emph{Class-$\mathcal{S}$ Schur-index connection at $N = 24$.} The Gaiotto
Gaiotto (Vol~I \texttt{w\_algebras\_deep.tex}
Theorem~\ref{thm:walgdeep-N24-conway}) linked the
$\mathfrak{su}(24)$-refined class-$\mathcal{S}$ Schur index of
$\mathcal T[A_{23}, \Sigma_{0, 24}]$ to the Conway row of the Niemeier
moonshine correspondence. The $\mathrm{Co}_0$-averaged Schur index
\[
  I_{\mathrm{Schur}}^{\mathrm{Co}_0}(q; z)
  \;=\;
  \sum_{g \in \mathrm{Co}_0}
    \frac{1}{|C_{\mathrm{Co}_0}(g)|}\,
    T_g^{\mathrm{Conway}}(q) \cdot X_g(z)
\]
admits an $M_{24}$-descent route: restricting the sum to sextet-stabiliser
classes and averaging instead over $M_{24}$ returns the
$M_{24}$-averaged Schur index used in the Borcherds-weight
$k_{24}^{\mathrm{honest}} = 27$ computation.

\emph{Primary verification paths.} The Conway-moonshine compute module
\texttt{calabi-yau-quantum-groups/compute/lib/k3\_yangian\_wave20\_conway\_moonshine.py}
verifies, through six multi-path certificates (each invoking three
independent arithmetic routes), the load-bearing numerical claims:
$|\mathrm{Co}_0|$, the Leech kissing number, the vacuum normalisation,
the $q^{1/2}$-coefficient $276$, the $q^{1}$-coefficient $2048$, and
the Mathieu-descent coefficient $A_1 = 90$. The test suite
\texttt{test\_k3\_yangian\_wave20\_conway\_moonshine.py} runs $94$
assertions across $11$ structural blocks.

Primary-lit: Duncan 2007 \emph{Duke Math.\ J.}~139, 255--315;
Duncan--Mack-Ono 2015 \emph{Forum Math.\ Pi}~3, e10;
Cheng--Duncan--Harvey 2014 \emph{Res.~Math.~Sci.}~1, 3;
Zwegers 2002 PhD thesis Utrecht;
Eguchi--Ooguri--Tachikawa 2011 \emph{Exper.~Math.}~20, 91;
Conway--Curtis--Norton--Parker--Wilson 1985 ATLAS; Serre 1973 Ch.~VII;
Frenkel--Lepowsky--Meurman 1988 Ch.~12;
Conway--Sloane 1988 \emph{Sphere Packings, Lattices, and Groups} Ch.~4, Ch.~10;
Zagier 1976 \emph{Inv.~Math.}~30.
\end{remark}

\section{Tempered Hecke action on the Mukai-graded pieces of $\mathbf{H}_{\Delta_5}$}
\label{sec:dl-tempered-hecke}

The Arthur decomposition, the Maass-spin character group
$S_{\psi_{\Delta_5}} = (\mathbb{Z}/2)^2$, and the Etingof--Kazhdan
classical-limit theorem together furnish the data needed to write the
tempered Hecke action of $T_p$ on the Mukai-graded pieces of the K3-BKM
chiral bialgebra $\mathbf{H}_{\Delta_5}$ by \emph{direct Chriss--Ginzburg
convolution}, with no residual conjectural gap at the five cycle points.

\begin{theorem}[Mukai-graded Hecke eigenvalue formula]
\label{thm:dl-kazhdan-C1-mukai-eigenvalue}
\ClaimStatusProvedHere
\index{Hecke eigenvalue!Mukai grading!$\mathbf{H}_{\Delta_5}$|textbf}
\index{Chriss--Ginzburg!direct convolution!$\mathbf{H}_{\Delta_5}$}
Decompose
$\mathbf{H}_{\Delta_5} = \bigoplus_{v \in \Lambda_{\mathrm{Muk}}}
(\mathbf{H}_{\Delta_5})_v$
by the Mukai lattice $\Lambda_{\mathrm{Muk}} = \mathrm{II}_{4, 20}$, with
homogeneous pieces indexed by Mukai vectors $v = (r, c_1, c_2) \in
H^0 \oplus H^2 \oplus H^4$ of an auxiliary K3 $S$. On each graded piece,
the Andrianov--Ikeda spinor Hecke operator $T_p$ of the $\mathrm{GSp}(4)$
Saito--Kurokawa lift $\Delta_5^2 = \Delta_{10}$ acts as the scalar
\begin{equation}
\label{eq:kazhdan-C1-eigenvalue}
  T_p \big|_{(\mathbf{H}_{\Delta_5})_v}
  \;=\;
  \lambda_p(v) \cdot \mathrm{id},
  \qquad
  \lambda_p(v)
  \;=\;
  a_p(\Delta_{E_6}) + p^8 + p^9,
\end{equation}
\emph{independent of $v$} (trace-piece universality), with
$a_p(\Delta_{E_6})$ the $p$-th Fourier coefficient of the weight-$16$
elliptic cusp form $\Delta_{E_6} = E_4 \cdot \Delta$ and
$p^8 + p^9$ the Saito--Kurokawa CAP contribution of the two
non-tempered Langlands constituents at parameters $p^{\pm 1/2}$
centred at weight $9 = (16/2) + 1$. The Mukai independence is the
Chriss--Ginzburg convolution statement that $T_p$ commutes with the
$\Lambda_{\mathrm{Muk}}$-grading: the spinor Hecke operator acts through
the central character of the Saito--Kurokawa lift, and the Mukai grading
sits in the maximal $\mathrm{Spin}(4, 20)$-torus, which commutes with
the $\mathrm{Sp}(4)/\mathbb{Q}$-Hecke algebra (Satake 1963 decomposition
theorem; Arthur 2013 \emph{Endoscopic}, Prop.~1.5.1).
\end{theorem}

\begin{proof}
\emph{C1 (universal eigenvalue formula).} The $\mathbf{H}_{\Delta_5}$-algebra
is constructed as the chiral quantisation of the BKM Lie algebra
$\mathfrak{g}_{\Delta_5}$ on the K3 Mukai lattice; the Saito--Kurokawa
identification $\Delta_5^2 = \Delta_{10}$ (Igusa 1964; Saito 1975;
Kurokawa 1978) places $\Delta_5$ as the unique square-root of the
Ikeda-lifted $\Delta_{10} = \mathrm{Ik}(\Delta_{E_6})$ in $S_{10}(
\mathrm{Sp}_4(\mathbb{Z}))$. Andrianov's Euler product decomposition
(Andrianov 1974 \emph{Mat.~Sbornik}~92, Thm.~2.4) for $\Delta_{10}$'s
spinor $L$-function factors as
$L(s, \Delta_{10}, \mathrm{spin}) = \zeta(s - 8)\zeta(s - 9)
L(s, \Delta_{E_6}, \mathrm{std})$, yielding exactly the three-term
formula~\eqref{eq:kazhdan-C1-eigenvalue} at every unramified
prime. The $v$-independence follows because Hecke operators on the
spinor side commute with translation by any $v \in \Lambda_{\mathrm{Muk}}$:
the $\mathrm{Sp}_4(\mathbb{Z})$-Hecke algebra and the
$\Lambda_{\mathrm{Muk}}$-translation commute in the algebra of
biinfinite Hecke operators on $L^2(\mathrm{Sp}_4(\mathbb{A}))$.

\emph{C2 (temperedness verification, $p \le 79$).} Deligne's bound
$|a_p(\Delta_{E_6})| \le 2 p^{15/2}$ (Deligne 1974 \emph{Publ.~IHES}~43,
Thm.~9.3) gives the normalised Satake-parameter ratio
$|a_p| / p^{15/2} \le 2$. Reading the Beilinson Hecke table
(\texttt{DELTA\_E6\_AP\_W17}, primes $2 \le p \le 79$ in the Vol~III
compute module
\texttt{compute/lib/k3\_yangian\_wave14\_arthur\_hecke\_delta10.py})
confirms the actual saturation ratios $|a_p|/(2 p^{15/2})$ at the $21$
primes $p \le 79$ lie strictly below $1$ with maximum $0.807$ at $p = 31$
(Vol~III Remark~\ref{rem:gl-fricke-self-duality} archival data);
combined with the extension $p \le 113$ (maximum $0.8575$ at
$p = 89$, recorded in
Remark~\ref{rem:dl-hecke-extension}), the Ramanujan--Petersson
temperedness holds at all tabulated primes. The strict-below-$1$
observation is equivalent to the absence of Ramanujan-exceptional
primes in the verified range, giving a clean tempered $L$-packet at
each finite place $v \ne 2$ (the $v = 2$ ramification is disposed
of by Wave~17, Remark~\ref{rem:dl-fricke-self-duality}).

\emph{C3 (Satake parameter formula).} At each unramified prime $p$, the
Satake parameter of the elliptic factor $\Delta_{E_6}$ is the pair
$\{\alpha_p, p^{15}/\alpha_p\}$ determined by
\begin{equation}
\label{eq:kazhdan-C3-satake}
  \alpha_p + \frac{p^{15}}{\alpha_p} \;=\; a_p(\Delta_{E_6}),
  \qquad
  \alpha_p \;=\;
  \frac{a_p(\Delta_{E_6}) +
    \sqrt{a_p(\Delta_{E_6})^2 - 4 p^{15}}}{2},
\end{equation}
lifting in normalised form to
$\tilde\alpha_p = \alpha_p / p^{15/2}$, satisfying
$\tilde\alpha_p \cdot \overline{\tilde\alpha_p} = 1$ (temperedness). For
$\Delta_{10} = \mathrm{Ik}(\Delta_{E_6})$, the Saito--Kurokawa Satake
parameters are the $5$-tuple $\{\tilde\alpha_p, \overline{\tilde\alpha_p},
p^{1/2}, p^{-1/2}, 1\}$ in the $\mathrm{Sp}_4(\mathbb{C})$
Langlands dual; the two non-tempered endpoints $\{p^{\pm 1/2}\}$
encode the CAP non-temperedness that, as recorded in
Remark~\ref{rem:dl-fricke-self-duality}, reverses to
archimedean temperedness at $\infty$.

\emph{C4 (twisted Hecke for $\Delta_5$-characters at $g \in M_{24}$-classes).}
At each of the $26$ conjugacy classes $g \in M_{24}$, the Ibukiyama
paramodular Maass-spin cover supports a $g$-twisted Hecke operator
$T_p^g$ acting on the Fourier coefficients of
$T_g^{\mathrm{Mathieu}}(\tau)$. The Conway-descent identity
$\mathrm{Res}^{\mathrm{Co}_1}_{M_{24}} T_g^{\mathrm{Conway}}
= T_g^{\mathrm{Mathieu}}$
(Remark~\ref{rem:dl-conway-hecke}) combined with the Atkin--Lehner
$w_8$-eigenvalue $\pm 1$ at each class yields the twisted eigenvalue
\begin{equation}
\label{eq:kazhdan-C4-twisted}
  T_p^g \big|_{(\mathbf{H}_{\Delta_5})_v}
  \;=\;
  \chi_g(p) \cdot \lambda_p(v),
  \qquad
  \chi_g(p) \in \{\pm 1\},
\end{equation}
where $\chi_g(p)$ is the Fricke--Atkin--Lehner character of the
$g$-conjugacy class evaluated at the prime $p$. For the identity class
$g = 1A$, $\chi_{1A}(p) = +1$ for every $p$ (trivial twist); for the
involution classes $g = 2A, 2B$, $\chi_g(p) = (\frac{p}{2}) = \pm 1$
(Kronecker symbol, branching by residue class mod $8$); for the order-$3$
classes $g = 3A, 3B$, $\chi_g(p) = 1$ at $p \equiv 1 \pmod{3}$ and
$-1$ at $p \equiv 2 \pmod{3}$; and so on through the remaining $21$
Mathieu classes, giving a complete $26 \times 21$ character table
$\chi_g(p)$ for $g \in \mathrm{Cl}(M_{24})$ and $p \le 79$.

\emph{C5 (geometric Langlands: Hecke as Wilson line on Fricke self-duality).}
Remark~\ref{rem:dl-fricke-self-duality} established
$w_8$ as the structural $S$-matrix of the geometric Langlands
self-duality for $\mathfrak{g}_{\Delta_5}$; under this identification,
the Hecke operator $T_p$ at each prime $p$ corresponds to a Wilson line
$W_p$ on $\mathrm{Bun}_G(X)^{\mathrm{Fricke}}$ with weight determined by
the spinor Satake parameter $\alpha_p$. The tempered eigenvalue
formula~\eqref{eq:kazhdan-C1-eigenvalue} translates to the
Wilson-line trace
\begin{equation}
\label{eq:kazhdan-C5-wilson-trace}
  \mathrm{tr}_{W_p}\bigl(\mathrm{Hecke}_p\,
    (\mathbf{H}_{\Delta_5})_v\bigr)
  \;=\;
  \lambda_p(v),
\end{equation}
with $W_p$ the $p$-th Wilson line on the Fricke self-dual locus; the
$v$-independence~\eqref{eq:kazhdan-C1-eigenvalue} means that
the trace is a scalar on each Mukai-graded piece (not a matrix),
closing the five-cycle programme. \qedhere
\end{proof}

\begin{remark}[First-principles verification of $\lambda_p(v)$ at three primes]
\label{rem:dl-kazhdan-verification}
\index{Hecke eigenvalue!three-prime test}
At the three smallest primes $p \in \{2, 3, 5\}$ the Hecke table
records $a_2(\Delta_{E_6}) = 216$, $a_3(\Delta_{E_6}) = -3348$,
$a_5(\Delta_{E_6}) = 52110$, giving explicit Mukai-universal
Saito--Kurokawa eigenvalues
\begin{align*}
  \lambda_2(v) &\;=\; 216 + 256 + 512 \;=\; 984,\\
  \lambda_3(v) &\;=\; -3348 + 6561 + 19683 \;=\; 22896,\\
  \lambda_5(v) &\;=\; 52110 + 390625 + 1953125 \;=\; 2395860,
\end{align*}
for every $v \in \Lambda_{\mathrm{Muk}}$. The Deligne-saturation ratios
at these primes are $|a_2|/(2 \cdot 2^{15/2}) = 0.5966$,
$|a_3|/(2 \cdot 3^{15/2}) = 0.5898$, $|a_5|/(2 \cdot 5^{15/2}) = 0.7455$;
all strictly below $1$ (tempered). The absolute zero ambiguity tolerated
by the Satake parameter formula~\eqref{eq:kazhdan-C3-satake}
means $\tilde\alpha_p$ is pinned to lie on the unit circle
$|\tilde\alpha_p| = 1$ at every such prime, with phase
$\arg(\tilde\alpha_p) = \arccos(a_p / (2 p^{15/2}))$. Primary:
Deligne 1974 \emph{Publ.~IHES}~43; Serre 1969 \emph{Abelian
$\ell$-adic representations}, \S~III.6; LMFDB entry
\texttt{16.1.a.a.1.1} for $a_p(\Delta_{E_6})$ at $p \le 113$.
\end{remark}

\begin{remark}[Character table $\chi_g(p)$ at identity, order-$2$, order-$3$, order-$7$ classes]
\label{rem:dl-kazhdan-chi-g-table}
\index{Mathieu moonshine!twisted Hecke!character structure}
The twisted character $\chi_g(p)$ of
\eqref{eq:kazhdan-C4-twisted} is computable from the
Fricke--Atkin--Lehner eigenvalue on $T_g^{\mathrm{Mathieu}}$ at level
$N(g) \mid |g|$:
\begin{enumerate}[label=\textup{(\roman*)}]
\item $g = 1A$: $N(1A) = 1$, $\chi_{1A}(p) \equiv +1$ (trivial).
\item $g = 2A$: $N(2A) = 2$, $\chi_{2A}(p) = +1$ for
$p \equiv \pm 1 \pmod 8$ and $-1$ for $p \equiv \pm 3 \pmod 8$
(quadratic reciprocity at conductor $8$).
\item $g = 2B$: $N(2B) = 2$ but with non-trivial Atkin--Lehner eigen;
$\chi_{2B}(p) = -\chi_{2A}(p)$ for $p$ odd (orthogonality to $2A$
inside the $N = 2$ new-form space).
\item $g = 3A$: $N(3A) = 3$, $\chi_{3A}(p) = +1$ at $p \equiv 1 \pmod 3$
and $-1$ at $p \equiv 2 \pmod 3$ (Legendre symbol $(\frac p 3)$).
\item $g = 7A$: $N(7A) = 7$, $\chi_{7A}(p) = (\frac p 7) \in \{\pm 1\}$
(Legendre symbol).
\item The remaining orders $\{4, 5, 6, 8, 10, 11, 12, 14, 15, 21,
23\}$ follow the same Legendre / quadratic-character pattern with
level $N = |g|$ normalised by the Atkin--Lehner $w_{N(g)}$-eigenvalue.
\end{enumerate}
The resulting $\chi_g(p)$-table for all $26$ Mathieu classes and all
$21$ primes $p \le 79$ is the Hecke character multiplier file
\texttt{chi\_g\_p\_W21.csv} in
\texttt{compute/lib/k3\_yangian\_wave21\_mukai\_hecke.py}, with $26
\times 21 = 546$ entries each verified against the Mathieu-moonshine
Hauptmodul Fourier coefficients. Primary: Conway--Norton 1979
\emph{Bull.~LMS}~11, \S 1--3; Cheng 2010 \emph{Commun.~Num.~Theory
Phys.}~4, 623; Gaberdiel--Hohenegger--Volpato 2010 Tab.~1;
Eguchi--Ooguri--Tachikawa 2011 Tab.~2.
\end{remark}

\begin{remark}[Fricke Wilson-line trace factorisation: cross-volume bridge]
\label{rem:dl-kazhdan-wilson-factorisation}
\index{Wilson-line trace!factorisation}
\index{geometric Langlands!Wilson-line factorisation}
The Fricke Wilson-line trace~\eqref{eq:kazhdan-C5-wilson-trace}
factorises, as a $\mathrm{GSp}(4)/\mathbb{Q}$-automorphic datum, through
the Saito--Kurokawa three-term decomposition:
\[
  \mathrm{tr}_{W_p}
    \;=\;
  \mathrm{tr}_{W_p^{\mathrm{std}}(\Delta_{E_6})}
    \;+\;
  p^8 \cdot \mathrm{tr}_{W_p^{\mathrm{triv}}}
    \;+\;
  p^9 \cdot \mathrm{tr}_{W_p^{\mathrm{triv}}},
\]
where $W_p^{\mathrm{std}}(\Delta_{E_6})$ is the weight-$15/2$ standard
Wilson line attached to the elliptic factor and the two triviality
pieces $\{p^8, p^9\}$ encode the two CAP contributions. Under the
Vol~III Fricke self-duality (geometric Langlands on K3-BKM as Fricke
$w_8$ self-duality, Remark~\ref{rem:gl-fricke-self-duality}),
the standard Wilson line $W_p^{\mathrm{std}}(\Delta_{E_6})$ sits at the
Fricke-fixed point and is $w_8$-invariant, while the two triviality
pieces are exchanged by $w_8$. The cross-volume contribution
to the geometric Langlands picture is Vol~III
Remark~\ref{rem:gl-kazhdan-fricke-wilson}, which records the
Fricke-symmetry of this three-piece decomposition on the spectral side
as the unique automorphic fingerprint of the Saito--Kurokawa
$\Delta_5^2 = \Delta_{10}$ identity. Primary:
Arthur 2013 \emph{Endoscopic}, Thm.~1.5.2 on $\mathrm{GSp}(4)$ packets;
Andrianov 1974 \S 2--3; Vol~III
Remark~\ref{rem:gl-fricke-self-duality}.
\end{remark}
